\documentclass[11pt]{article}
\usepackage[T1]{fontenc}
\usepackage{lmodern}
\usepackage[letterpaper,margin=1in]{geometry}
\usepackage{amsmath,amssymb,amsthm,mathtools,esint}
\usepackage{microtype}
\usepackage{enumitem,longtable,array,xurl}
\usepackage[hidelinks]{hyperref}
\hypersetup{pdftitle={A forward formalization blueprint for the trilinear Hilbert transform estimate},pdfauthor={Mathematical blueprint based on the supplied manuscript},pdfsubject={Self-contained proofs and forward declaration dependencies}}
\setlength{\emergencystretch}{2em}
\setlength{\parskip}{2pt}
\setcounter{tocdepth}{1}
\setlist[enumerate]{itemsep=2pt,topsep=4pt,leftmargin=*}
\newtheorem{theorem}{Theorem}[section]
\newtheorem{lemma}[theorem]{Lemma}
\newtheorem{proposition}[theorem]{Proposition}
\newtheorem{corollary}[theorem]{Corollary}
\theoremstyle{definition}
\newtheorem{definition}[theorem]{Definition}
\newcommand{\R}{\mathbb R}
\newcommand{\Z}{\mathbb Z}
\newcommand{\Q}{\mathbb Q}
\newcommand{\C}{\mathbb C}
\newcommand{\T}{\mathbb T}
\newcommand{\E}{\mathbb E}
\newcommand{\one}{\mathbf 1}
\newcommand{\avint}{\fint}
\newcommand{\phys}{\mathrm{phys}}
\newcommand{\supp}{\operatorname{supp}}
\newcommand{\osc}{\operatorname{osc}}
\newcommand{\rank}{\operatorname{rank}}
\newcommand{\spanq}{\operatorname{span}_{\Q}}
\newcommand{\sym}{\operatorname{sym}}
\newcommand{\poly}{\operatorname{poly}}
\newcommand{\dist}{\operatorname{dist}}
\newcommand{\eqtag}[1]{\refstepcounter{equation}\tag{\theequation}\label{#1}}
\numberwithin{equation}{section}

\newcommand{\bp}[3]{\par\smallskip\begingroup\footnotesize
\noindent\hypertarget{bp:#1}{\textsf{Declaration #1.}}\quad\nolinkurl{#2}\par
\noindent\textsf{Dependencies:} #3.\par\endgroup\smallskip}
\title{A forward formalization blueprint\\for the $L^3$ trilinear Hilbert transform estimate\\[5pt]
\large Complete mathematical proofs and declaration dependencies}
\author{Based on the supplied manuscript and its revised argument}
\date{8 October 2026}
\begin{document}
\maketitle
\begin{abstract}
This document reorganizes the full supplied argument into forward dependency order and supplies the additional structural proofs needed for a formalization using mathlib-level foundations and standard Hardy--Littlewood maximal theory. A finite cyclic $U^3$ inverse theorem with quasipolynomial bounds is proved directly by sampling, additive energy, local symmetry and explicit lattice coordinates; no nilsequence inverse theorem is assumed. An elementary nonsharp transference estimate replaces Minkowski's second theorem. The short-block comparison uses one vector-valued maximal norm and does not import extrapolation. The singular-integral construction, its parameter choices, the normalized absorption, and the passage to the principal-value operator are included. Declaration identifiers, proposed Lean names, and dependency records accompany the proofs. This is a mathematical blueprint, not a Lean formalization or a claim of kernel verification.
\end{abstract}
\tableofcontents
\clearpage
\section{Foundations, conventions, and proof interfaces}\label{sec:foundations}
\paragraph{Scope.}
This is a mathematical blueprint for formalization, not a Lean source file or a claim of kernel verification. Each named mathematical dependency is either proved before its use or belongs to the foundation list below. Proposed declaration names in the declaration notes are new names for the eventual project, not names asserted to exist in mathlib. The supplied \nolinkurl{shorter_argument.tex} \cite{R}, derived from \cite{M}, is the source for the singular-integral construction. The additional finite inverse argument and the elementary lattice estimates below replace its external degree-two nilsequence inverse theorem and its appeal to Minkowski's second theorem. No result about nilmanifolds is imported.

\paragraph{Library reference.}
The reference mathlib revision is
\begin{center}\small\texttt{6be3d352cec3ceabd28bf85266d28646e02f1975}.\end{center}
The documentation's source link for \nolinkurl{Mathlib.MeasureTheory.Group.GeometryOfNumbers} identifies this revision \cite{mathlibGeometry}. The inventory below records mathematical interfaces, rather than pretending that a single existing declaration has exactly the notation of this manuscript. In particular that geometry-of-numbers file supplies the first Minkowski theorem, not the successive-minima transference estimate used here. We prove the latter. The final formalization must pin its Lean toolchain and this mathlib revision together; compiling this \LaTeX{} document performs no Lean check.

\paragraph{Permitted mathematical foundations.}
We use the following mathlib-level material: arithmetic of natural numbers, integers, rationals, reals and complexes; finite sets, sums and products; induction and finite well-founded recursion; finite-dimensional linear algebra, determinants, Gram--Schmidt, and Euclidean compactness; Lebesgue and Haar measure, Bochner integration, Tonelli, dominated and monotone convergence, H\"older, Minkowski and Jensen; complete $L^p$ spaces, density of bounded simple functions and of compactly supported continuous functions for $1\le p<\infty$, and density of Schwartz functions in $L^3(\R)$; Hilbert-space adjoints, orthogonal projections, the Riesz representation theorem, and real Hahn--Banach separation of a point from a closed convex set. These are supplied by the measure, integration, $L^p$, projection, duality and separation modules of mathlib. Measurability and almost-everywhere equivalence are always included in an $L^p$ assertion.

For Fourier theory the imported statements are precisely these. With $e(t)=\exp(2\pi it)$, the Fourier transform $\widehat f(\xi)=\int f(x)e(-x\cdot\xi)\,dx$ extends to a unitary map on $L^2(\R^a)$, agrees with the integral transform on integrable $L^2$ functions, and has the usual inverse. On $(\R/\Z)^a$, the integer characters are a complete orthonormal basis of $L^2$; an absolutely summable Fourier series of a continuous function converges uniformly to that function. The corresponding Hilbert-valued statements follow by orthonormal expansion and Tonelli. Verified library interfaces include \nolinkurl{MeasureTheory.Lp.norm_fourier_eq}, \nolinkurl{MeasureTheory.Lp.inner_fourier_eq}, \nolinkurl{SchwartzMap.toLp_fourier_eq} in \nolinkurl{Mathlib.Analysis.Fourier.LpSpace}, and \nolinkurl{UnitAddTorus.mFourierBasis}, \nolinkurl{UnitAddTorus.hasSum_sq_mFourierCoeff}, and \nolinkurl{UnitAddTorus.hasSum_mFourier_series_apply_of_summable} in \nolinkurl{Mathlib.Analysis.Fourier.AddCircleMulti}; see \cite{mathlibLp,mathlibTorus}. Finite cyclic Fourier identities will be proved explicitly.

The only additional analytic import is the following permitted part of Hardy--Littlewood maximal theory. For the uncentered maximal operator $M$ on $\R$, the scalar weak $(1,1)$ inequality and strong $L^r$ inequality hold for $1<r<\infty$. The vector-valued form, for $1<r,s<\infty$ and any finite measure space $(\Omega,\nu)$, is
\[
 \left\|\left(\int_\Omega |Mh_\omega|^s\,d\nu(\omega)\right)^{1/s}\right\|_{L^r(\R)}
 \le C_{r,s}\left\|\left(\int_\Omega |h_\omega|^s\,d\nu(\omega)\right)^{1/s}\right\|_{L^r(\R)}.
\]
The constant is independent of the measure space and of the number of coordinates. This is the Fefferman--Stein maximal inequality in its standard integral form. The dyadic versions follow by domination by $M$ (or by the elementary dyadic stopping proof). No weighted extrapolation theorem is assumed. The permission to use Calder\'on--Zygmund and Littlewood--Paley theory is not otherwise needed as a separate theorem import.

\paragraph{Representatives and norms.}
All scalar and vector functions are strongly measurable. A local normalized norm means
\[
 \|f\|_{p,I}=\left(|I|^{-1}\int_I\|f(x)\|^p\,dx\right)^{1/p}
 \quad(1\le p<\infty),
\]
with the essential supremum at $p=\infty$. An unqualified norm uses ordinary Lebesgue measure. Inner products in the written argument are linear in their first variable; translate to mathlib's convention by conjugating the displayed pairings, not by changing a projection equation. Finite sums of exceptional null sets can always be removed at once. In a fixed finite tree every list and every family of operator rows is finite, although its cardinality need not be uniformly bounded. Classical choice is permitted; effectiveness of selection is not asserted.

\begin{lemma}[Quantitative closure rules]\label{lem:budget-rules}
For $t\ge2$, call a positive quantity quasipolynomially bounded in $t$ when it is at most $\exp(C(1+\log t)^d)$ for some fixed $C,d\ge1$. A fixed number of products, sums, positive fixed powers, and compositions preserves this property. Polynomial bounds in the logarithms are preserved by a fixed number of sums, products and compositions. These assertions remain true with a coefficient $C_q$ depending on $q$, while every exponent is fixed independently of $q$.

If $D\le\exp(C_q k^d)$, $m\le C\log(2k)$, and
\[
 F\le\bigl(C(m+1)\log(2D)\bigr)^{C(m+1)},
\]
then $F\le\exp(C_q(1+\log(2k))^{d'})$ for an absolute $d'$. If $P=\lceil\exp(\alpha k^\theta)\rceil$ with fixed $\alpha,\theta>0$, every such $F$ and every fixed power of $k$ is at most $P^\epsilon$ for every fixed $\epsilon>0$ once $k$ is sufficiently large.
\end{lemma}
\begin{proof}
Write $u=1+\log t$. For a sum of $n$ quantities bounded by $e^{C u^d}$, the logarithm is at most $\log n+Cu^d$. Products add these logarithms; fixed powers multiply them. For a composition, $1+\log(e^{Cu^d})\le(1+C)u^d$, so the next bound has logarithm at most $C'(1+C)^{d'}u^{dd'}$. Iterate this calculation only the stated fixed number of times. The polynomial assertions are proved by the same degree arithmetic. When coefficients depend on $q$, multiplying and exponentiating them changes only that coefficient.

For the second assertion take logarithms of the bound for $F$. Its logarithm is at most
\[
 C(m+1)\bigl(\log C+\log(m+1)+\log\log(2D)\bigr)
 \le C_q(1+\log(2k))^{d'}.
\]
Finally $(\log k)^A/k^\theta\to0$ for fixed $A,\theta>0$, as follows on putting $u=\log k$ and comparing $u^A$ with the exponential $e^{\theta u}$. Comparing logarithms proves the claimed $P^\epsilon$ estimates. All threshold choices are made after fixing their coefficients.
\end{proof}
\bp{budget-rules}{THT.Quantitative.closure}{library arithmetic, real exponential limits}

\paragraph{How existential constants are formalized.}
An assertion $D\le\exp(k^{O(1)})$ in a construction means that the construction exports an integer degree $d$, independent of $q$, and a coefficient and lower threshold depending on $q$, such that $\log D\le C_q k^d$ above the threshold. The degree is returned before the later choice of $q$. The proof of each construction supplies it by the preceding closure rules and the displayed finite recurrences. A statement with this meaning is not the weaker assertion ``for each $q$ some degree exists.'' The latter would not suffice for the final absorption. Bounds written $P^C$ similarly export an absolute degree. This convention eliminates implicit changes in the order of quantifiers.

\begin{lemma}[A fixed Gevrey cutoff and partition]\label{lem:gevrey}
There are nonnegative $C^\infty$ compactly supported cutoffs equal to one on any prescribed smaller interval inside a larger interval, and a nonnegative $C^\infty$ function $\rho$ supported in $(-1,1)$ with $\sum_{n\in\Z}\rho(t-n)=1$. They can be chosen of one fixed Gevrey order: on a fixed pair of intervals their $m$th derivatives are at most $(C(m+1))^{C(m+1)}$. Rescaling a transition width by $a$ multiplies the derivative bound by at most $a^{-m}$, with a fixed additional normalization factor if an integral is normalized. Product cutoffs in dimension $d$ have the corresponding product derivative bounds.
\end{lemma}
\begin{proof}
Set $g(t)=e^{-1/t}$ for $t>0$ and $g(t)=0$ otherwise. On $t>0$, its $m$th derivative is $e^{-1/t}P_m(1/t)$, where $P_0=1$ and
\[
 P_{m+1}(z)=z^2(P_m(z)-P_m'(z)).
\]
The degree is at most $2m$ and the sum of the absolute coefficients is at most $\prod_{j=0}^{m-1}(1+2j)\le(2m)^m$ for $m\ge1$. For $0\le z\le1$ bound the polynomial by its coefficient sum; for $z\ge1$ bound it by that sum times $z^{2m}$. Since $\sup_{z\ge0}z^{2m}e^{-z}=(2m/e)^{2m}$, the derivative is at most $(C(m+1))^{3(m+1)}$. Each derivative tends to zero as $t\downarrow0$; induction using the definition of the derivative proves that the zero extension is $C^\infty$.

The function $g(t)g(1-t)$ is a positive bump on $(0,1)$. Leibniz's formula gives the same type of derivative bound. Normalize its positive integral, rescale, and integrate once to obtain a smooth increasing step function with fixed transition interval. Products of two translated steps give the stated cutoffs. For the partition, take a nonnegative probability bump $h$ supported in $(-1/4,1/4)$ and set
\[
 \rho(t)=\int_{-1/2}^{1/2}h(t-u)\,du.
\]
Its support lies in $[-3/4,3/4]$, its derivatives have the required bounds, and Tonelli gives
\[
 \sum_{n\in\Z}\rho(t-n)
 =\int_{\R}h(t-u)\sum_{n\in\Z}\one_{[n-1/2,n+1/2]}(u)\,du=1.
\]
The possible double counting at integer half-points is on a null set. All sums are locally finite. The scaling and product assertions follow from the chain and product rules.
\end{proof}
\bp{gevrey}{THT.Cutoff.existsGevreyPartition}{library calculus}

\begin{lemma}[Fourier tails and elementary oscillatory rows]\label{lem:fourier-foundation}
A smooth function on a $d$-torus whose derivatives of order $m$ are bounded by $A_m$ has Fourier coefficients satisfying
\[
 |\widehat F(n)|\le A_m\left(\frac{\sqrt d}{2\pi\|n\|}\right)^m\quad(n\ne0).
\]
Its dyadic Fourier shells therefore have absolute mass at most $C_d A_m T^{d-m}d^{m/2}$ at $T\le1+\|n\|<2T$. In particular all the uniform Fourier series and truncations used below follow from an explicitly chosen derivative order exceeding the dimension and the requested tail power.

If real frequencies $\xi_\nu$ are $\delta$-separated, then on an interval $J$ of length $r$,
\[
 \int_J\left|\sum_\nu c_\nu e(\xi_\nu y)\right|^2dy
 \le C(r+\delta^{-1})\sum_\nu|c_\nu|^2.
\]
If a real quadratic polynomial $p$ satisfies $|p'|\ge A>0$ on $J$, then $|\int_J e(p)|\le C/A$.
\end{lemma}
\begin{proof}
Select a coordinate with $|n_j|\ge\|n\|/\sqrt d$ and integrate by parts $m$ times; periodic boundary terms cancel. There are at most $(CT)^d$ integer points in the indicated shell. Summing the resulting geometric series proves the first assertion, including uniform convergence by the imported torus Fourier theorem.

For the second, majorize $\one_J$ by a nonnegative fixed smooth bump rescaled to length $r$. Two integrations by parts bound its Gram entry by $Cr(1+r|\xi_\nu-\xi_\mu|)^{-2}$. In the ordered frequency list, separation makes every absolute row and column sum at most
\[
 Cr\left(1+2\sum_{j\ge1}(1+r\delta j)^{-2}\right)\le C(r+\delta^{-1}).
\]
For a matrix $(a_{ij})$ with both absolute row and column sums at most $B$, Cauchy--Schwarz gives
\[
 \sum_i\left|\sum_j a_{ij}z_j\right|^2
 \le\sum_i\left(\sum_j|a_{ij}|\right)\sum_j|a_{ij}|\,|z_j|^2
 \le B^2\sum_j|z_j|^2.
\]
This proves the required Gram bound. The same computation proves the operator-valued use with scalar entry bounds. A common bounded multiplier costs its supremum norm, and duality gives the analysis-row form.

Finally integrate by parts using $e(p)/(2\pi i p')$. Boundary terms cost $C/A$ and the integral term costs $C\int_J |p''|/|p'|^2\le C/A$, because the affine derivative has constant sign and is monotone.
\end{proof}
\bp{fourier-foundation}{THT.Fourier.uniformTailsAndRows}{gevrey, library Fourier and integration by parts}

\paragraph{Cube definitions.}
On a compact abelian group with Haar probability, let $\mathcal C$ denote complex conjugation and, for bounded functions indexed by $\omega\in\{0,1\}^j$, put
\[
 \Lambda_j(f_\omega)=\E_{x,h_1,\ldots,h_j}
 \prod_{\omega\in\{0,1\}^j}\mathcal C^{|\omega|}f_\omega(x+\omega\cdot h).
\]
For a single repeated function the value is nonnegative: disintegrating in $h_j$, the variables $x$ and $x+h_j$ are independent, so its diagonal equals
\[
 \E_{h_1,\ldots,h_{j-1}}
 \left|\E_x\prod_{\omega\in\{0,1\}^{j-1}}
 \mathcal C^{|\omega|}f(x+\omega\cdot h)\right|^2.
\]
Define $\|f\|_{U^j}$ to be the nonnegative $2^j$th root of this diagonal. Write $\Lambda_\square=\Lambda_3$. Only $j=2,3$ will be needed. Finite group expectations use probability counting measure; convolution is $(f*g)(x)=\E_y f(y)g(x-y)$. These definitions do not yet assume a triangle inequality; the cube lemma below proves the one used in the argument.
\bp{cube-definition}{THT.Gowers.cubeAverage}{library Haar, finite products, nonnegative roots}

\begin{lemma}[Signs, symmetrization, and finite Fourier identities]\label{lem:finite-fourier}
For independent signs $\epsilon_i=\pm1$ of equal probabilities and complex numbers $a_i$, one has, for every $p\ge2$,
\[
 \bigl(\E|\sum_i\epsilon_i a_i|^p\bigr)^{1/p}
 \le C\sqrt p\left(\sum_i|a_i|^2\right)^{1/2}.
\]
Two independent sign randomizations of a double-indexed array satisfy every fixed moment estimate with a constant depending only on that moment. If $a_1,\ldots,a_n$ are independent uniform samples of a nonempty set $A$ in a finite abelian group and $p\ge2$, then
\[
 \E\left\|n^{-1}\sum_{i=1}^n f(\cdot-a_i)-f*\mu_A\right\|_p^p
 \le\left(C\sqrt{p/n}\,\|f\|_p\right)^p,
\]
where convolution and $L^p$ norms use probability counting measure and $\mu_A=\one_A/(|A|/|G|)$.

For $G=\Z/q\Z$, define $\widehat f(\xi)=\E_x f(x)e(-\xi x/q)$. Then Fourier inversion, Parseval, and $\widehat{f*g}=\widehat f\widehat g$ hold, with counting measure on the dual. In particular
\[
 \|f\|_{U^2(G)}^4=\sum_\xi|\widehat f(\xi)|^4,
 \qquad
 \|f\|_{U^3(G)}^8=\E_h\|f(\cdot+h)\overline f\|_{U^2(G)}^4.
\]
\end{lemma}
\begin{proof}
For real coefficients, expand the $2m$th moment. Only index lists in which every index occurs an even number of times survive. Pairing the $2m$ positions overcounts these lists and bounds their sum by
\[
 \frac{(2m)!}{2^m m!}\left(\sum_i a_i^2\right)^m
 \le (2m)^m\left(\sum_i a_i^2\right)^m.
\]
Use the real and imaginary parts and the triangle inequality for complex coefficients. For general $p$, take $m=\lceil p/2\rceil$ and use probability-space norm monotonicity. Applying this result conditionally to each sign family and using Minkowski for the intervening square norm proves the double-index assertion.

For sampling, introduce independent copies $a_i'$ and apply Jensen to replace the mean by their average. The differences $f(\cdot-a_i)-f(\cdot-a_i')$ have a joint distribution invariant under independent sign changes. Apply the sign estimate pointwise, integrate, and use Minkowski in $L^{p/2}$ to obtain
\[
 \left\|\left(\sum_i|f(\cdot-a_i)-f(\cdot-a_i')|^2\right)^{1/2}\right\|_p
 \le 2\sqrt n\,\|f\|_p.
\]
Divide by $n$ and absorb the fixed factor into $C$.

For the finite Fourier assertions, the geometric sum gives $\E_x e(\xi x/q)=1$ when $\xi=0$ in $G$, and zero otherwise. Substituting this identity proves inversion and Parseval by finite expansion. Changing variables proves the convolution identity. Expanding the four vertices of a square gives the $U^2$ identity, while separating the last cube increment gives the displayed $U^3$ recursion. These computations apply to every positive integer $q$.
\end{proof}
\bp{finite-fourier}{THT.FiniteFourier.samplingAndCubeIdentities}{cube-definition, library finite sums, Holder, Minkowski}

\begin{lemma}[Cube Cauchy--Schwarz and $L^2$ continuity]\label{lem:cube-foundation}
On a compact abelian probability group, the order-three cube expression is nonnegative on the diagonal and its eighth root is a seminorm. Its mixed form satisfies
\[
 |\Lambda_\square(g_\omega)|\le\prod_{\omega\in\{0,1\}^3}\|g_\omega\|_{U^3}.
\]
On a real torus, including $\T_{16}$, it also satisfies
\[
 |\Lambda_\square(g_\omega)|\le\prod_\omega\|g_\omega\|_2.
 \eqtag{eq:cube-L2}
\]
Consequently it extends continuously to eight $L^2$ inputs.
\end{lemma}
\begin{proof}
Write the cube vertices as $a_1^{\omega_1}+a_2^{\omega_2}+a_3^{\omega_3}$, with six independent Haar variables. This has the same distribution as $(x+\omega\cdot h)_\omega$: use $x=a_1^0+a_2^0+a_3^0$ and $h_i=a_i^1-a_i^0$. For arbitrary functions of the three chosen variables, integrate first in $a_i^0,a_i^1$. Cauchy--Schwarz in the other four variables bounds the square of the mixed integral by the product of the two integrals obtained by repeating, on both sides of coordinate $i$, respectively its $0$-face and its $1$-face. The two resulting expressions are integrals of absolute squares. Apply this inequality to coordinate $1$, then to coordinate $2$ in both factors, and finally to coordinate $3$ in all four factors. The outcome is
\[
 |\Lambda_\square(g_\omega)|^8
 \le\prod_\omega\Lambda_\square(g_\omega,g_\omega,\ldots,g_\omega).
\]
Each diagonal on the right is nonnegative: it equals the average over two increments of the squared absolute value of the average of the corresponding second multiplicative difference. Taking eighth roots proves the mixed inequality. Expansion into the $2^8$ choices in the cube of $f+g$ bounds it by $(\|f\|_{U^3}+\|g\|_{U^3})^8$, proving the triangle inequality; homogeneity is immediate.

For the $L^2$ bound split the vertices into $\{000,100,010,001\}$ and $\{111,011,101,110\}$. In each set the four linear forms in $(x,h_1,h_2,h_3)$ have integer determinant $1$ or $-1$. Their Haar distributions are therefore independent. Cauchy--Schwarz between the two four-factor groups factors their squared integrals into the eight squared $L^2$ norms. Telescoping a difference in the eight inputs proves continuity for bounded functions, and truncation proves it for all $L^2$ inputs.
\end{proof}
\bp{cube-foundation}{THT.Gowers.cubeCSAndL2}{cube-definition, library Haar, Holder}

\begin{lemma}[The maximal norm used later]\label{lem:vector-maximal}
For $2<q<3$ and $p>q$, on any finite measure space $\Omega$,
\[
 \|M_q f\|_{L_y^pL_\omega^4}\le C_{p,q}\|f\|_{L_y^pL_\omega^4},
 \qquad M_qf=(M|f|^q)^{1/q}.
\]
The constants are independent of the number of random signs, lattice indices, and torus variables included in $\Omega$.
\end{lemma}
\begin{proof}
Apply the permitted vector maximal theorem with $r=p/q>1$, $s=4/q>1$, and $h=|f|^q$. The exact power identities for the two iterated norms, followed by a $q$th root, give the assertion. Product spaces of signs and torus variables are a single measure space $\Omega$; no interchange of unequal norms is used.
\end{proof}
\bp{vector-maximal}{THT.Maximal.oneMixedNorm}{permitted vector Hardy--Littlewood theorem}


\section{Elementary lattice and torus facts}\label{sec:lattices}
A full lattice in a $d$-dimensional Euclidean space means the integer span of a real basis. Its dual is the set of vectors having integer inner product with every lattice vector. A subgroup $L\subset\Z^d$ is saturated when $nv\in L$, $n\in\Z\setminus\{0\}$, implies $v\in L$.

\begin{lemma}[Integer bases and saturated subgroups]\label{lem:integer-basis}
A primitive integer vector extends to an integer basis of $\Z^d$. A saturated subgroup of $\Z^d$ is spanned by an initial part of some integer basis of $\Z^d$. Consequently its quotient is free, and a rational subspace intersected with $\Z^d$ is saturated.
\end{lemma}
\begin{proof}
For integers $a,b$, not both zero, let $g=\gcd(a,b)>0$ and choose integers $u,v$ with $ua+vb=g$. The matrix with rows $(u,v)$ and $(-b/g,a/g)$ has determinant one and sends $(a,b)$ to $(g,0)$. Successively apply this operation to the nonzero coordinates of a vector, changing a sign if necessary. A primitive vector is thereby sent to the first coordinate vector by an integer matrix with determinant $\pm1$. The inverse matrix is integral by the cofactor formula and supplies the asserted basis.

For the subgroup assertion induct on $d$. If the subgroup is zero there is nothing to do. Otherwise choose a nonzero vector in it and divide by the greatest common divisor of its coordinates; saturation keeps the resulting primitive vector in the subgroup. Change integer coordinates to make this vector the first coordinate vector. Subtract the first coordinate from every other subgroup element. The remaining subgroup in $\{0\}\times\Z^{d-1}$ is saturated: a multiple belonging to it first implies membership in the original subgroup, and the zero first coordinate is already given. Apply induction to this subgroup. This proves the basis assertion and the description of the quotient. If $V$ is a rational vector subspace and $nv\in V\cap\Z^d$, then $v=n^{-1}(nv)\in V\cap\Z^d$, proving the last assertion.
\end{proof}
\bp{integer-basis}{THT.Lattice.saturatedAdaptedBasis}{integer Bezout identity, determinants, induction}

\begin{lemma}[A size-reduced shortest-projection basis]\label{lem:reduced-basis}
A full Euclidean lattice $\Lambda$ has a basis $b_1,\ldots,b_d$ whose Gram--Schmidt vectors $b_i^*$, of lengths $a_i>0$, satisfy
\[
 b_j=b_j^*+\sum_{i<j}\mu_{ij}b_i^*,\qquad |\mu_{ij}|\le\tfrac12,
 \qquad a_{i+1}\ge\tfrac{\sqrt3}{2}a_i.
\]
The vector $b_i^*$ is shortest nonzero in the projection of $\Lambda$ orthogonal to $\operatorname{span}(b_1,\ldots,b_{i-1})$. With $c=2/\sqrt3$, the basis and its dual basis $q_i$ obey
\[
 \|b_i\|\le c^{i-1}a_i,\qquad \|q_i\|\le 4^d/a_i.
\]
\end{lemma}
\begin{proof}
The intersection of a lattice with a bounded set is finite, by applying the inverse of a matrix of generators and using finiteness of bounded sets of integer vectors. Thus a shortest nonzero lattice vector exists. It is primitive in lattice coordinates, since division by a nontrivial common divisor would shorten it. Lemma \ref{lem:integer-basis} extends it to a lattice basis. Project the remaining basis vectors orthogonally to its perpendicular space. They are real-linearly independent there, so their integer span is again a full lattice. Induct on dimension in that projected lattice and lift the resulting basis. For each successive lift subtract integer multiples of the already chosen vectors, in decreasing index order, to place its Gram--Schmidt coefficients in $[-1/2,1/2]$. This does not change any later orthogonal projection. It gives the required basis and shortest-projection property.

In the lattice obtained by projecting off the first $i-1$ vectors, the projected vector $b_{i+1}$ has squared length $a_{i+1}^2+\mu_{i,i+1}^2a_i^2$. Shortness of $b_i^*$ gives $a_i^2\le a_{i+1}^2+a_i^2/4$. This proves the adjacent-length inequality. In particular $a_i\le c^{j-i}a_j$ for $i<j$, and
\[
 \|b_j\|^2\le a_j^2+\tfrac14\sum_{i<j}a_i^2
 \le a_j^2\left(1+\tfrac14\sum_{r=1}^{j-1}c^{2r}\right)
 =c^{2(j-1)}a_j^2.
\]
For the dual bound write the basis matrix in orthonormal coordinates as $DU$, where $D$ is diagonal with entries $a_i$ and $U$ is upper triangular with diagonal one and other entries $\mu_{ij}$. The finite geometric expansion of $(1+(U-1))^{-1}$ shows, for $i<j$, that
\[
 |(U^{-1})_{ij}|\le\sum_{r=1}^{j-i}\binom{j-i-1}{r-1}2^{-r}
 =\tfrac12(\tfrac32)^{j-i-1}.
\]
The coordinates of $q_i$ are $(U^{-1})_{ij}/a_j$, $j\ge i$, and zero otherwise. Use $a_j\ge c^{-(j-i)}a_i$ and sum the squares of the resulting geometric majorant. The sum is at most $16^d/a_i^2$, including $j=i$. Taking square roots proves the displayed bound.
\end{proof}
\bp{reduced-basis}{THT.Lattice.shortestProjectionBasis}{integer-basis, Gram--Schmidt, finite lattice intersections}

\begin{lemma}[Nonsharp successive-minima transference]\label{lem:transference}
Let $K$ be a compact symmetric convex body with nonempty interior in $\R^d$, $d\ge1$, and let $K^\circ=\{y:|\langle x,y\rangle|\le1\ (x\in K)\}$. Write $\lambda_j(K,\Lambda)$ for the least $t>0$ such that $tK$ contains $j$ real-linearly independent vectors of a full lattice $\Lambda$. Then
\[
 1\le \lambda_j(K,\Lambda)\lambda_{d+1-j}(K^\circ,\Lambda^*)
 \le d16^d\le C_d,\qquad C_d=(100d)^{10d}.
\]
In particular no form of Minkowski's second theorem is needed below.
\end{lemma}
\begin{proof}
The minima are finite and attained: a sufficiently large dilate contains a lattice basis, and in a fixed larger dilate there are only finitely many lattice vectors. For the lower bound choose the two independent lists realizing the indicated minima. Their spans have dimensions adding to $d+1$. They cannot be mutually orthogonal, so one cross inner product is a nonzero integer. Its absolute value is at least one and, by polarity, at most the product of the minima.

First let $K$ be the Euclidean unit ball. The first $j$ vectors of Lemma \ref{lem:reduced-basis} have norms at most $c^{j-1}a_j$. The independent dual vectors $q_j,\ldots,q_d$ have norms at most $4^dc^d/a_j$, since $a_i\ge c^{-(i-j)}a_j$ for $i\ge j$. Therefore the product of the two minima is at most $4^dc^{2d}\le16^d$.

For general $K$, maximize the absolute determinant of $d$ columns $v_i\in K$. Compactness gives a maximum, and the interior hypothesis makes it positive. Let $A$ be the matrix of maximizing columns and let $E=A(B_2)$. Convexity and symmetry give $A(B_1)\subset K$. Replacing one column by $x\in K$ and using maximality and Cramer's rule gives $\|A^{-1}x\|_\infty\le1$, hence $K\subset A(B_\infty)$. The elementary inclusions between these three norm balls imply
\[
 d^{-1/2}E\subset K\subset d^{1/2}E,
 \qquad d^{-1/2}E^\circ\subset K^\circ\subset d^{1/2}E^\circ.
\]
Apply the Euclidean assertion to $A^{-1}\Lambda$, whose dual is $A^T\Lambda^*$. Each of the two minima for $K,K^\circ$ is at most $\sqrt d$ times its corresponding ellipsoid minimum. This proves the upper bound. The final numerical inequality follows directly from $d\ge1$.
\end{proof}
\bp{transference}{THT.Lattice.successiveMinimaProduct}{reduced-basis, determinants, compact convex sets}

\begin{lemma}[A small lattice ball in a coefficient box]\label{lem:lattice-box}
For every full lattice $\Lambda\subset\R^d$, $d\ge1$, there are part of a lattice basis $b_1,\ldots,b_m$ and integers $L_i\ge1$ such that
\[
 \Lambda\cap B_2(0,\epsilon_d)\subset
 \left\{\sum_{i=1}^m n_i b_i:n_i\in\Z,\ |n_i|\le L_i\right\}
 \subset B_2(0,1/4),\qquad \epsilon_d=(100d)^{-10d}.
\]
Every partial coordinate sum in this box also belongs to $B_2(0,1/4)$. The empty box is interpreted as $\{0\}$.
\end{lemma}
\begin{proof}
Use Lemma \ref{lem:reduced-basis}, put $c=2/\sqrt3$ and $t_0=(100dc^{2d})^{-1}$, and take $m$ to be the last index with $a_m\le t_0$, or zero if none exists. For $i\le m$,
\[
 \|b_i\|\le c^{i-1}a_i\le c^{m-1}a_m\le c^dt_0\le(100d)^{-1}.
\]
Set $L_i=\lfloor(4d\|b_i\|)^{-1}\rfloor$. These integers are positive and satisfy
\[
 L_i\ge(8d\|b_i\|)^{-1}\ge(8dc^da_i)^{-1},\qquad
 \sum_{i=1}^m L_i\|b_i\|\le1/4.
\]
This proves the second inclusion and the partial-sum assertion. If $v\in\Lambda$ has norm at most $\epsilon_d<t_0$, its projection beyond the first $m$ vectors must vanish. Otherwise it would have norm at least the shortest nonzero projected length $a_{m+1}>t_0$; when $m=d$ this step is unnecessary. Thus $v=\sum_{i\le m}n_ib_i$ with integer coefficients. By the dual bound,
\[
 |n_i|=|\langle v,q_i\rangle|\le4^d\epsilon_d/a_i
 \le(8dc^da_i)^{-1}\le L_i.
\]
The middle inequality and $\epsilon_d<t_0$ follow from the displayed explicit choices for every $d\ge1$. This proves the first inclusion, including $m=0$.
\end{proof}
\bp{lattice-box}{THT.Lattice.smallBallCoefficientBox}{reduced-basis}

\begin{lemma}[Rational annihilators and progression Haar measure]\label{lem:torus-annihilator}
Let $V\subset\Q^d$ be a rational subspace, put $L=V\cap\Z^d$, and set
\[
 H=\{x\in(\R/\Z)^d:n\cdot x=0\pmod1\text{ for every }n\in L\}.
\]
Then $H$ is a connected torus. Its Haar integral of the character $e(n\cdot x)$ is one for $n\in L$ and zero otherwise. For distinct integers $i,j$ and any integer $l$, the map
\[
 H^2\longrightarrow H^2,\qquad (u,v)\longmapsto(u+(i-l)v,u+(j-l)v)
\]
is surjective and preserves Haar probability. These statements include the zero-dimensional torus.
\end{lemma}
\begin{proof}
The subgroup $L$ is saturated. Apply Lemma \ref{lem:integer-basis}; the induced integer coordinate automorphism of the torus sends $H$ to the coordinate torus with its first $\dim V$ coordinates zero. This proves connectedness and the character integral by one-dimensional Haar integration in each remaining coordinate. Multiplication by a nonzero integer on a circle is surjective: divide a real representative by that integer. It is therefore surjective on a coordinate torus and on $H$. Given outputs $a,b$, solve $(j-i)v=b-a$ in $H$ and set $u=a-(i-l)v$. This proves surjectivity. The pushforward of Haar probability under a surjective continuous homomorphism is translation invariant and has mass one, so uniqueness of Haar probability proves the measure assertion. In the coordinate-torus model this can equivalently be checked first on characters and then by the complete Fourier basis.
\end{proof}
\bp{torus-annihilator}{THT.Torus.rationalAnnihilator}{integer-basis, Haar uniqueness, torus Fourier basis}


\section{Local forms, finite constants, and stopping}\label{sec:reduction}
A dyadic grid is a nested sequence of partitions of $\R$ into half-open intervals, with lengths $r_s=r_02^{-s}$, $s\in\Z$, in which each interval is partitioned by its two children. Statements about endpoints are always understood up to their null set. The ordinary grid is realized by $r_0=1$ and intervals $[2^{-s}m,2^{-s}(m+1))$; translations and dilations give further grids. A finite rooted tree is a finite ancestor-closed set of descendants of one root, completed by terminal cells to a partition whenever a stopping partition is used. A collection of summation cells need not be ancestor-closed.

For $y_j=x+jt$, $0\le j\le3$, and a cell $I$ of length $r$, define
\[
 H_I(z_0,z_1,z_2,z_3)=r^{-2}\int_{\R^2}w_I(x,t)\prod_{j=0}^3(z_j)_{\phys}(x+jt)\,dx\,dt.
 \eqtag{eq:local}
\]
Fix $c_0>0$, $C_0\ge1$. An admissible kernel is smooth, supported where all four $x+jt$ lie in $I$ at distance at least $c_0r$ from its endpoints, and satisfies
\[
 \int_\R w_I(u-jt,t)\,dt=0\quad(u\in\R,\ 0\le j\le3).
 \eqtag{eq:marginals}
\]
Each derivative of total order $n$ is bounded by $r^{-n}(C_0(n+2))^{C_0(n+2)}$. Kernels may be chosen separately at different cells and multiplied by scalars of modulus at most one. Every later constant may depend on this fixed class. The class will be fixed large enough for the explicit compact kernel constructed at the end.

\begin{lemma}[The positive local form estimate]\label{lem:absolute}
If $p_j\in[1,\infty]$ and $\sum_jp_j^{-1}\le2$, then
\[
 r^{-2}\int|w_I(x,t)|\prod_j\|z_j(x+jt)\|\,dx\,dt
 \le C\prod_j\|z_j\|_{p_j,I}.
 \eqtag{eq:absolute}
\]
The same statement, with fixed constants, holds for the corresponding progression integral on a connected torus and for a local integer-lattice progression sum normalized at its scale, provided the kernel is bounded and supported in a fixed scale-sized box.
\end{lemma}
\begin{proof}
For distinct $i,j$, the map $(x,t)\mapsto(x+it,x+jt)$ has determinant $j-i\ne0$. The support condition and this change of variables give the estimate when slots $i,j$ have exponent one and the others exponent infinity. The all-infinity estimate follows from the bounded area of the support and the kernel bound. A single exponent one is treated by integrating in one coordinate and using the scale-sized support in the other. These bounds have a common fixed constant.

Here is the positive interpolation step explicitly. The polytope $0\le a_j\le1$, $\sum_ja_j\le2$ is the convex hull of the vectors with entries zero or one and at most two ones. To verify this without an interpolation theorem, if at least two coordinates are strictly between zero and one, perturb one upward and the other downward until one reaches an endpoint in each of the two directions. The original vector is a convex combination of the two resulting vectors, each with fewer fractional coordinates, and the total sum is unchanged. Repeat. With one fractional coordinate left, interpolate between its two endpoints; both endpoints still have sum at most two, because all other coordinates are integers. This is a finite induction.

Write $(p_j^{-1})_j=\sum_v\theta_v v$, with $\theta_v\ge0$ and $\sum_v\theta_v=1$, over these endpoint vectors. For finite $p_j$, apply the endpoint bound in configuration $v$ to the $j$th function $\|z_j\|^{p_jv_j}$; a zero exponent means the constant one. For $p_j=\infty$ the corresponding factor is bounded first and removed. The pointwise product of these endpoint integrands, to powers $\theta_v$, is the original product, because $p_j\sum_v\theta_vv_j=1$. H\"older with exponents $1/\theta_v$ and the positive measure $r^{-2}|w_I|\,dx\,dt$ proves the desired bound. Zero or infinite norms are treated by truncation and monotone convergence.

On a connected torus the two-coordinate map preserves Haar probability by Lemma \ref{lem:torus-annihilator}, or by the same coordinate-torus argument for the full torus. Thus the endpoint bounds hold with constant one for a bounded multiplier of supremum one. On an integer lattice the pair map is injective onto its image, so a sum over its image is bounded by the full product sum. The scale normalization and bounded support give the other endpoints. The identical positive H\"older argument completes both variants.
\end{proof}
\bp{absolute}{THT.LocalForm.absoluteBound}{torus-annihilator, change of variables, Holder}

Let $M_qf=(M|f|^q)^{1/q}$. For $2<q<3$ and a positive integer $N$, define $C_N(q)$ to be the least real constant at least one such that
\[
 \sum_{I\in\mathcal I}|I|\,|H_I(f)|\le C_N(q)\int_\R\prod_{j=0}^3M_qf_j
 \eqtag{eq:CN}
\]
for every grid, admissible kernel family, compactly supported $L^q$ inputs, and finite collection in at most $N$ consecutive depths. To define it with real rather than extended-real quotients, take the supremum of the ratios with positive finite denominator and then its maximum with one. A zero denominator forces every relevant local norm product to vanish, hence every summand to vanish. Infinite denominators give a trivial extended-real inequality because $C_N\ge1$.

\begin{lemma}[Finiteness before absorption]\label{lem:finite-CN}
The preceding constant is well-defined and $1\le C_N(q)\le1+CN$. Its defining inequality holds at the least constant itself.
\end{lemma}
\begin{proof}
The all-$L^2$ case of Lemma \ref{lem:absolute}, followed by normalized norm monotonicity, bounds $|H_I(f)|$ by $C\prod_j\|f_j\|_{q,I}$. For every $x\in I$, each of these four local norms is at most $M_qf_j(x)$. At a fixed depth the intervals are disjoint; integrate on them, and then sum the at most $N$ depths. Every ratio in the definition is at most $CN$, so its supremum is finite. A supremum majorizes every member of its defining set. The zero and infinite cases were dealt with before taking ratios, proving the last assertion.
\end{proof}
\bp{finite-CN}{THT.LocalForm.finiteDepthConstant}{absolute, scalar maximal definition}

\paragraph{Finite data and operator families.}
The normalization used from this point until the outer reduction is $\|f_j\|_{q,I}\le1$ on every cell of the chosen finite subtree. It is a hypothesis on the original inputs, not on arbitrary restrictions of a later side. All selection procedures are completed before an operator estimate is applied. An estimate is uniform over each resulting fixed family; no linearity of the map selecting that family from an input is asserted. Fresh private labels can be taken from a fixed countably infinite set, with disjoint injections for the separate catalogs. Every finite construction uses only finitely many labels, so its proofs reduce to the corresponding finite coordinate space. Norms of the resulting operators do not depend on that cardinality.

\begin{lemma}[Rational heights and finite linear systems]\label{lem:rational-linear}
For a rational number its height is the maximum of the absolute numerator and positive denominator in reduced form. Vector and matrix heights are entrywise. A matrix of size at most $d$ and height at most $B\ge2$ has at most $(3B^2)^{d^2}$ possible values. Nonsingular inverses, rational orthogonal projections determined by such equations, and right inverses on their ranges have heights and operator norms at most $\exp(C(d+1)^4\log(2B))$. A basis chosen from a bounded-height spanning list obeys the same bound without a factor for the length of that list.
\end{lemma}
\begin{proof}
There are at most $3B^2$ fractions of height at most $B$. Clearing all at most $d^2$ denominators multiplies entries by an integer at most $B^{d^2}$ and gives an integer matrix whose entries are at most $B^{d^2+1}$. The determinant expansion bounds every minor by $d!B^{d(d^2+1)}$. A nonzero integer determinant has absolute value at least one. Cramer's formula therefore bounds both numerators and denominators of an inverse, and its entries' absolute values, by exponentials of the stated form. For a system of rank $r$, choose $r$ independent rows and $r$ pivot columns and solve on that nonsingular minor, setting the other variables to zero. This gives a right inverse on the range. If $S$ has independent rows, the orthogonal projection onto its row space is $S^T(SS^T)^{-1}S$; the same bounds and a fixed number of products give the assertion. The orthogonal complement is obtained by subtracting this projection from the identity. Selecting independent rows or a basis never increases the entrywise bound and uses at most $d$ members of a list.
\end{proof}
\bp{rational-linear}{THT.Rational.heightLinearAlgebra}{determinants, finite-dimensional projections}

All uses of quasipolynomial bounds have the quantifier convention of Lemma \ref{lem:budget-rules}. In particular the degree is independent of $q$, physical frequencies, lengths, and depth. For nonnegative cell data the exact packing identity is
\[
 \sum_{I\subseteq R}|I|a_I=\int_R\sum_{I\ni y}a_I\,dy.
 \eqtag{eq:path}
\]
It follows by writing $|I|=\int_R\one_I$ and applying finite Tonelli. An \emph{attempt} is a construction from a normalized root. If the total length of its first exceptional subroots is at most $\rho|R|$, $0\le\rho<1$, the total length of all roots under restarting is at most $|R|/(1-\rho)$: generation $n$ has total length at most $\rho^n|R|$ by induction. For several rules take maximal cells in their union and add their length bounds. Stops remove original summands at and below the first crossing, not the fixed auxiliary functions used above that crossing. Continue an old inherited column family below an unused cell by retaining its columns and refitting coefficients. Each fresh attempt proves its own root-normalized energy bounds; restricting a previously obtained global energy bound is not a substitute.
\bp{path-packing}{THT.Tree.pathIntegralAndRestarts}{finite Tonelli, geometric series}

\begin{lemma}[First-crossing cuts]\label{lem:cuts}
Let $0<\rho<1$ and let all data be fixed on a finite tree rooted at $R$.
If $a_I\ge0$ and $\sum_I|I|a_I\le A|R|$, removing the first cells where the ancestor sum exceeds $A/\rho$ costs at most $\rho|R|$. The retained ancestor sums are at most $A/\rho$ everywhere, including below the removed cells.

There is also the following version for fixed functions $f_\nu$ and specified hulls of cells. Suppose the maximal cells of each hull have union $\Omega_\nu$,
\[
 \sum_\nu\|f_\nu\|_2^2\le E|R|,\qquad
 \sum_\nu|\Omega_\nu|\le M|R|.
\]
First-crossing cuts of total length at most $\rho|R|$ can make the squared maxima over the retained hulls satisfy
\[
 \sum_\nu\sup_{I\ni x\text{ retained in hull }\nu}
       \avint_I\|f_\nu\|^2
 \le C_\rho E\bigl(1+\log(2+M)\bigr)
 \quad(x\in R).
 \eqtag{eq:hull-cap}
\]
An empty supremum is zero. Hull intervals and maxima are kept on their entire cells; only summands, not the functions $f_\nu$, are removed.
\end{lemma}
\begin{proof}
For the first assertion the disjoint first-crossing cells lie in the set where the full path sum exceeds $A/\rho$. Its measure is at most $\rho|R|$ by \eqref{eq:path} and Markov. A point below a removed cell retains only the ancestors strictly before its first crossing, proving the everywhere assertion.

For the second let $m_\nu$ be the full squared hull maximum. The dyadic weak $(1,1)$ bound and the support bound give
\[
 \sum_\nu|\{m_\nu>t\}|\le\min(M,CE/t)|R|.
\]
Take $T=C_\rho E$ large enough that individual crossings of $T$ cost at most $\rho|R|/2$. Integration yields
\[
 \sum_\nu\int_R\min(m_\nu,T)
 \le CE\bigl(1+\log(2+MT/E)\bigr)|R|.
\]
This formula also holds if $T<CE/M$, by the bound $MT\le CE$; zero energy is trivial. Markov bounds crossings of $C_\rho E(1+\log(2+M))$ by the remaining $\rho|R|/2$, after enlarging $C_\rho$. Implement both rules on partial hull maxima using only ancestors through the current depth. Before an individual crossing the full partial sum equals its truncated version. The full-tree distribution bounds dominate these partial quantities. Exclude a violating cell before including its value. As in the first assertion, below a cut only earlier ancestors remain, so \eqref{eq:hull-cap} holds on whole retained cells. Taking maximal cells in the union combines the rules.
\end{proof}

\bp{cuts}{THT.Tree.firstCrossingAndHullCuts}{path-packing, permitted weak maximal inequality}


\section{Column projections and bounded-height frequencies}\label{sec:columns}
On a tree rooted at $J$, attach to each label $v$ a birth cell $J_v$, a scalar function $\phi_v$, and a private number $\epsilon_v>0$. Work in $\mathcal H=\C\oplus\ell^2(V)$ with columns
\[
 \Phi_v(y)=(\phi_v(y),\sqrt{\epsilon_v}\,e_v),\qquad y\in J_v.
\]
Different labels always have different private coordinates. The common budget $D\ge2$ means
\[
 \sqrt{\epsilon_v}\ge D^{-1},\qquad \|\Phi_v\|_\infty\le D,
 \qquad \sum_v\one_{J_v}(y)\le D.
 \eqtag{eq:columnbudget}
\]
The last condition counts all available labels on a path; there is no bound on their total number on disjoint branches. A mask $V^i(I)\subseteq\{v:I\subseteq J_v\}$ is inherited if it increases on descendants. Write $P_I^i$ for projection onto constant-coefficient combinations of its restricted columns, extended by zero off $I$. Direct sums $P_\pi^i$ on refining stopping partitions are nested orthogonal projections, because an old combination remains an allowed combination on each smaller cell. Adaptive masks are held fixed in every operator estimate.

\paragraph{Column data interface.}
The data consist of a finite label type, birth cells, measurable scalar columns and positive private weights, together with the three inequalities in \eqref{eq:columnbudget}. Each mask is a map from cells to available finite label sets with the stated inheritance relation. A projection is defined from its finite-dimensional range, not from an algorithm for choosing a basis. Direct-sum nesting follows by restricting a constant-coefficient combination to each refined cell, so it is a proved consequence of the mask data. A row-family interface consists of bounded linear maps, their square-sum inequality, and equality with right multiplication by the indicator of the decision cell. Later compositions must retain that exact cell-locality equation.
\bp{column-data}{THT.Columns.InheritedSystem}{finite dyadic data, library Hilbert projections}

\begin{lemma}\label{lem:coeff}
For available coefficients,
\[
 D^{-1}\|c\|_{\ell^2}\le\left\|\sum_vc_v\Phi_v(y)\right\|_{\mathcal H}
 \le D^{3/2}\|c\|_{\ell^2}.
\]
Their projection $P_{I,A}$ satisfies
\[
 \|P_{I,A}z\|_{\infty,I}\le D^5\avint_I\|z\|.
\]
For $F=(f,0)$ and $P_{I,A}F=\sum_vc_v\Phi_v$,
\[
 \avint_I\bigl(f-(P_{I,A}F)_{\phys}\bigr)\overline{\phi_v}=\epsilon_v c_v.
 \eqtag{eq:predictor-moment}
\]
\end{lemma}
\begin{proof}
Private coordinates give the lower bound; Cauchy--Schwarz and at most $D$ columns give the upper bound. If $E_yc=\sum c_v\Phi_v(y)$ and $G=\avint_I E_y^*E_y\,dy$, then $G\ge D^{-2}\mathrm{Id}$ and $\|E_y\|\le D^{3/2}$. The projection kernel $E_yG^{-1}E_t^*$ is bounded by $D^5$. Orthogonality to $\Phi_v$ gives the last identity, and the lower bound also proves uniqueness of coefficients.
\end{proof}
\bp{coeff}{THT.Columns.coefficientsAndKernel}{column-data, library Hilbert projections}


\begin{lemma}[Nested products]\label{lem:prefix}
Let $P_n^i$, $0\le n\le T$, $1\le i\le m$, be nested orthogonal projections on a Hilbert space, and $Z_n=P_n^m\cdots P_n^1$. If $\sum_\nu\|B_\nu h\|^2\le\|h\|^2$, then, for arbitrary indices $n(\nu)$,
\[
 \left(\sum_\nu\|B_\nu Z_{n(\nu)}g\|^2\right)^{1/2}
 \le L_T^m\|g\|,\qquad L_T=1+\lceil\log_2(T+1)\rceil,
 \eqtag{eq:prefix}
\]
and
\[
 \sum_{n=1}^T\|(Z_n-Z_{n-1})g\|^2\le m^2L_T^{2(m-1)}\|g\|^2.
 \eqtag{eq:product-jump}
\]
The empty product has row bound one. For a fixed family of disjoint index intervals $[r,t]$, the same jump estimate holds with $Z_t-Z_r$ and $L_T$ on the right.
\end{lemma}
\begin{proof}
Set $\Delta_0^i=P_0^i$ and $\Delta_n^i=P_n^i-P_{n-1}^i$. Decompose a prefix of $\{0,\ldots,T\}$ into dyadic blocks, at most one of each length. At a fixed length the sums $\Delta_A^i=\sum_{n\in A}\Delta_n^i$ are orthogonal. Expand each factor of $Z_{n(\nu)}$. For a fixed tuple of lengths, group rows with the same actual blocks and apply the row hypothesis to their common vector. The sum over all block tuples is at most
\[
 \sum_{A_1,\ldots,A_m}\|\Delta_{A_m}^m\cdots\Delta_{A_1}^1g\|^2\le\|g\|^2
\]
by successive Bessel inequalities, with no commutation between layers. The triangle inequality over at most $L_T^m$ length tuples proves \eqref{eq:prefix}.

Telescope
\[
 Z_n-Z_{n-1}=\sum_{i=1}^mP_n^m\cdots P_n^{i+1}\Delta_n^iP_{n-1}^{i-1}\cdots P_{n-1}^1.
\]
Discard the left contractions and use Cauchy--Schwarz in $i$. For fixed $i$, the middle increments are Bessel rows; \eqref{eq:prefix} applied to the $i-1$ right factors gives \eqref{eq:product-jump}. For disjoint intervals use $P_t^i-P_r^i$ as the middle rows and right prefix index $r$. These increments are again orthogonal.
\end{proof}
\bp{prefix}{THT.Projections.nestedPrefixRows}{library orthogonal projections, dyadic index decomposition}


\begin{lemma}[Spatial compression]\label{lem:compression}
For inherited column projections satisfying \eqref{eq:columnbudget}, suppose local rows satisfy
\[
 \sum_\nu\|B_\nu z\|^2\le\|z\|_{L^2(J;\mathcal H)}^2,
 \qquad B_\nu=B_\nu\one_{u(\nu)},
\]
where $u(\nu)$ are arbitrary tree cells. Then
\[
 \left(\sum_\nu\|B_\nu P_{u(\nu)}^m\cdots P_{u(\nu)}^1z\|^2\right)^{1/2}
 \le\bigl(C(m+1)^C\log(2D)\bigr)^{C(m+1)}\|z\|_2.
 \eqtag{eq:compression}
\]
Complements may replace factors. Fixed averages of such compositions cost only the total variation of the averaging measure. The estimate is independent of spatial depth and total label count.
\end{lemma}
\begin{proof}
The empty product is the hypothesis. Add one generation below the last row cells. Stop the averages $a_K=\avint_K\|z\|$ whenever they first exceed twice the average at the preceding stopping root; write $\mathcal S$ for these roots. Assign $u$ to its deepest such root $K$. Its size segment consists of the intervening cells. Cut it at the next roots and at final leaves, obtaining $\pi_K$. Every cut cell has average at most $4a_K$, by its parent's stopping bound.

On each cut cell project $z$ onto all first-layer columns used by the segment's ancestors of that cell, and call the resulting function $g_K$, zero off $K$. Lemma \ref{lem:coeff} gives, with $G=4D^5$,
\[
 \|g_K\|_\infty\le Ga_K,\qquad \|g_K\|_2^2=\langle z,g_K\rangle,
 \qquad P_u^1g_K=P_u^1z.
 \eqtag{eq:moment-representative}
\]
The last identity follows by adding the preserved moments over the cut partition of $u$. Since stopping averages double along a path,
\[
 \sum_Ka_K^2|K|\le\frac43\|M_d\|z\|\|_2^2\lesssim\|z\|_2^2.
\]
We also need $S=\sum_K\|g_K\|_2^2\lesssim\log(2D)\|z\|_2^2$, not a power of $D$. At a point split the roots into $a_K<M_d\|z\|/(GD^2)$ and its complement. In the first class the sum of $\|g_K(y)\|$ is at most $2D^{-2}M_d\|z\|(y)$. The second class has at most $L_D=2+\lceil\log_2(GD^2)\rceil$ members. By \eqref{eq:moment-representative}, pointwise Cauchy--Schwarz and then integration, $S\le4D^{-2}\|z\|_2^2+\sqrt{L_D}\,\|z\|_2\sqrt S$. Solving this quadratic inequality proves the required bound.

Fix a segment and abbreviate $g=g_K$, $a=a_K$. Put $\tau=D^{-10}$, $b=2D(m+1)$, and $T=\lceil b^{128(m+1)^2}\rceil$. At each cell let $c(u)$ be the last-layer coefficient vector of $P_u^m\cdots P_u^1g$, extended by zero to all labels. From a root $R$ declare an anchor whenever $\|c(u)-c(A)\|_2>\tau a$ for the preceding anchor $A$. After $T$ transitions send that cell to a new attempt with the same $g,a$.

Complete each anchor generation to a stopping partition of $R$ by retaining terminal cuts; on those cuts extend the original masks by refitting and then keep them fixed. The resulting projections $P_n^i$ are nested. On a restart root $R'$ every successive anchor difference has pointwise norm greater than $D^{-1}\tau a$, by private-coordinate recovery. Consequently \eqref{eq:product-jump} gives
\[
 TD^{-2}\tau^2a^2\sum_{R'}|R'|\le m^2L_T^{2m}\|\one_Rg\|_2^2.
\]
Using $\|g\|_\infty\le4D^5a$ once for lengths and once for the energy on $R'$ yields both
\[
 \sum_{R'}|R'|\le E|R|,\qquad
 \sum_{R'}\|\one_{R'}g\|_2^2\le E\|\one_Rg\|_2^2,
 \qquad E=\frac{16D^{32}m^2L_T^{2m}}{T}<\frac12.
\]
For the last inequality, $b\ge8$, $L_T\le129(m+1)^2\log_2b\le b^4$, so $E\le16b^{34+8m-128(m+1)^2}<1/2$. Thus all attempt roots in the segment satisfy $\sum_R|R|\le2|K|$ and $\sum_R\|\one_Rg\|_2^2\le2\|g\|_2^2$.

A row assigned to $R$ has a last anchor $A$ of generation $j<T$, and $u(\nu)$ lies wholly in its partition atom. Its input product is the restriction of $Z_j(\one_Rg)$ plus a coefficient error $e_{\nu,v}$ of $\ell^2$ norm at most $\tau a$. The anchor rows have squared sum at most $L_T^{2m}\|\one_Rg\|_2^2$ by \eqref{eq:prefix}, regardless of the physical depths of the rows. Write $\widetilde\Phi_v=\one_{J_v}\Phi_v$. Locality gives the error as
\[
 \sum_ve_{\nu,v}B_\nu(\one_R\widetilde\Phi_v).
\]
Cauchy--Schwarz in $v$ and the original row inequality bound its squared row sum by
\[
 \tau^2a^2\sum_v\|\one_R\widetilde\Phi_v\|_2^2
 \le\tau^2a^2D^3|R|.
\]
Here the final inequality is exactly \eqref{eq:columnbudget} integrated over $R$, so repeated row times and labels on disjoint branches are not counted.

Sum segments and attempts in the row Hilbert space. The preceding bounds give
\[
 \left(\sum_\nu\|B_\nu P_{u(\nu)}^m\cdots P_{u(\nu)}^1z\|^2\right)^{1/2}
 \lesssim\bigl(L_T^m\sqrt{\log(2D)}+D^{-17/2}\bigr)\|z\|_2.
\]
Since $L_T\lesssim(m+1)^3\log(2D)$, this proves \eqref{eq:compression}. Expand each complement into identity minus projection, giving at most $2^m$ terms, and use the integral triangle inequality for averages.
\end{proof}
\bp{compression}{THT.Projections.spatialCompression}{coeff, prefix, path-packing, permitted scalar maximal theory}


\begin{lemma}[Height bins]\label{lem:height}
For $\xi\in\R^d$ and $B\ge2$, the nonzero numbers $|a\cdot\xi|$ with $a\in\Q^d$ of height at most $B$ have logarithms in at most $d$ intervals of length $C(d+1)^3\log(2B)$. They meet at most $C(d+1)^4\log(2B)$ dyadic magnitude bins. If $r_n=r_02^{-n}$, $u\in\{1,2\}$ and $0<\alpha\le\beta$, then
\[
 \#\{n:\alpha\le r_n^u|a\cdot\xi|\le\beta\text{ for some }a\}
 \le C(d+1)^4\bigl(\log(2B)+\log(2\beta/\alpha)\bigr),
\]
and, with empty maxima zero,
\[
 \sum_n\max\bigl(\{r_n^u|a\cdot\xi|:0<r_n^u|a\cdot\xi|\le\epsilon\}\cup\{0\}\bigr)
 \le C(d+1)^4\log(2B)\epsilon.
 \eqtag{eq:tiny-sum}
\]
The same conclusions hold for suprema of such coordinates. For Euclidean norms of at most $h$ coordinates add $Cd\log(2h)$ to the bin counts and tail coefficient.
\end{lemma}
\begin{proof}
A minimally dependent set $a_1,\ldots,a_k$ has $k\le d+1$. Clear each vector's denominators by $q_i\le B^d$. Its integer entries are at most $B^d$. Signed maximal minors on $k-1$ independent rows give a dependence $\sum c_ia_i=0$ with all $c_i$ nonzero integers and $|c_i|\le H=d!B^{d^2+d}$. A zero minor would contradict minimality. Set $R=4dH$. Vectors whose nonzero evaluation magnitudes successively increase by more than $R$ are independent: in a minimal dependence the largest term has coefficient of absolute value at least one, while all smaller terms sum to at most $dH/R$ times that magnitude. Greedily cover the evaluation set by intervals $[s,Rs]$ starting at its smallest uncovered element. There are at most $d$ such intervals, and $\log_2R\le C(d+1)^3\log(2B)$. This proves the first assertion and, by translating and dilating logarithms, the scale count. In one occupied bin $[2^j,2^{j+1})$, sum the geometric tail starting at $r_n^u2^j\le\epsilon$; it is at most $4\epsilon$. Summing bins proves \eqref{eq:tiny-sum}. A coordinate maximum is one of its coordinates, and a Euclidean norm is between that maximum and $\sqrt h$ times it, giving the stated enlargement. The zero-dimensional case is empty and is interpreted accordingly.
\end{proof}
\bp{height}{THT.Frequencies.heightBins}{rational-linear, finite independence, geometric series}




\section{Quadratic chart definitions and closure}\label{sec:chart-definitions}
\begin{definition}[Quadratic chart]\label{def:chart}
A quadratic chart atom is
\[
 \phi(y)=e(\lambda y^2+\beta y)\sum_{n\in\Z^d}
 \chi(\theta y-n)e\bigl((\theta y)^TC(\theta y-n)\bigr).
 \eqtag{eq:atom}
\]
Here $\theta\in\R^d$, $\lambda,\beta\in\R$ are unrestricted, $C$ is rational, and $\chi$ is smooth and compactly supported. Complexity at most $H\ge e^2$ means $d\le\log H$, height$(C)\le H$, support in $[-H,H]^d$, and
\[
 \|\partial^\alpha\chi\|_\infty
 \le H^{c_0(|\alpha|+1)}(|\alpha|+1)^{c_0(|\alpha|+1)}
\]
for a fixed Gevrey order $c_0$. A bounded atom has supremum norm at most one. The case $d=0$ includes ordinary quadratic phases. There are at most $(2H+3)^d$ summands at a point, so normalization costs a quasipolynomial. Scalar sizes not included in these complexity bounds are always external coefficients.

\end{definition}
\bp{chart-definition}{THT.Chart.QuadraticAtom}{gevrey, finite rational matrices}

\begin{lemma}[Elementary chart-operations]\label{lem:chart-operations}
Complex conjugation, a fixed finite product, and affine changes of the real argument preserve quadratic charts, with quasipolynomial complexity in the original complexities. The bound is independent of the real translation and dilation parameters. Boundedness is preserved by conjugation and affine substitution, and by products of bounded atoms.
\end{lemma}
\begin{proof}
Conjugation and bounded products preserve this class: negate phases, or concatenate horizontal coordinates and use block diagonal matrices. Affine changes of argument also preserve it, with quasipolynomial complexity independent of their real coefficients. To verify translation, clear the denominators of $C$ by $P\le H^{d^2}$ and write a horizontal offset as $m_0+a_0$, with $m_0\in P\Z^d$ and $\|a_0\|_\infty\le P$. Reindexing, with $w=\theta y-n$, gives modulo integers
\[
 (\theta y+m_0+a_0)^TC(w+a_0)
 =(\theta y)^TCw+(\theta y)^TCa_0+m_0^TC\theta y+a_0^TC(w+a_0)+m_0^TCa_0.
\]
The middle linear terms and final constant go into unrestricted phases; the remaining amplitude $\chi(w+a_0)e(a_0^TC(w+a_0))$ has controlled derivatives. Scaling only changes $\theta,\lambda,\beta$. The same calculation works with any real linear map $V:\R^k\to\R^d$ in place of $y\mapsto\theta y$.

\end{proof}
\bp{chart-operations}{THT.Chart.affineAndProductClosure}{chart-definition, rational-linear, gevrey}

For $z\in\R^k$, $1\le k\le3$, write
\[
 Q(z)=e(z^T\Lambda z+b^Tz)\sum_{n\in\Z^d}\chi(Vz-n)e\bigl((Vz)^TC(Vz-n)\bigr),
\]
where $\Lambda$ is symmetric. Put $C_{\rm sy}=(C+C^T)/2$, $C_{\rm sk}=(C-C^T)/2$ and $A=C-C^T$.



\section{Finite additive arguments for the inverse theorem}\label{sec:additive}
In this section a finite group average is normalized by the group cardinality. Sumsets and their cardinalities use ordinary counting. All groups in Fourier arguments are cyclic; the sumset lemmas hold in every abelian group. For a finite set $A$, $mA$ is its $m$-fold sumset, with $0A=\{0\}$. The sampling method follows the Croot--Sisask mechanism \cite{CS}; the derivative-frequency and local-symmetry organization follows the degree-two inverse method \cite{GT}. Every quantitative statement used here is proved below rather than imported from those papers.

\begin{lemma}[Four-step paths in a dense bipartite graph]\label{lem:four-paths}
Let a bipartite graph have parts of size at most $n$, with at least $\eta n^2$ edges, $0<\eta\le1$. There is a set $A'$ in the first part with $|A'|\ge c\eta n$ such that any ordered pair of its vertices has at least $c\eta^5n^3$ four-step paths between them. Repeated vertices are allowed. If edges are properly colored, then the start vertex and the ordered four colors determine the whole path.
\end{lemma}
\begin{proof}
Pad the second part by isolated vertices to size $n$ and choose a uniform vertex there. Let $Y$ be its neighborhood. Then $\E|Y|\ge\eta n$ and $\E|Y|^2\ge\eta^2n^2$. Call an ordered pair in the first part bad if its codegree is less than $c_0\eta^2n$, where $c_0=10^{-5}$ is sufficient. Its probability of belonging to $Y^2$ is its codegree divided by $n$, so the expected number of bad ordered pairs in $Y^2$ is at most $c_0\eta^2n^2$. Some choice of $Y$ therefore satisfies
\[
 |Y|^2-100\#\{\text{bad ordered pairs in }Y^2\}\ge\tfrac12\eta^2n^2.
\]
In particular $|Y|\ge\eta n/\sqrt2$ and fewer than $|Y|^2/100$ ordered pairs are bad. Delete each vertex having more than $|Y|/4$ bad partners in $Y$. Fewer than $|Y|/25$ vertices are deleted. Given two retained vertices, at least $|Y|/2$ vertices of $Y$ are good partners of both. For every such middle vertex, independently choose a common neighbor with the first endpoint and a common neighbor with the second. This gives at least
\[
 (|Y|/2)(c_0\eta^2n)^2\ge c\eta^5n^3
\]
four-step paths, with distinct choices giving distinct paths. The retained set has the claimed size. In a proper coloring, at a specified vertex there is at most one incident edge of a specified color. Apply this observation successively along four edges to obtain the last assertion. The argument also covers nonintegral lower bounds: the displayed inequalities are inequalities for integer counts.
\end{proof}
\bp{four-paths}{THT.Additive.denseGraphFourPaths}{finite probability, Cauchy--Schwarz, deletion count}

\begin{lemma}[An energy-to-small-difference-set lemma]\label{lem:BSG}
Let $A$ be a nonempty finite subset of an abelian group, $|A|=n$, and suppose
\[
 \#\{(a,b,c,d)\in A^4:a+b=c+d\}\ge n^3/K,\qquad K\ge1.
\]
There is $A'\subset A$ with $|A'|\ge cK^{-1}n$ and $|A'-A'|\le CK^{10}|A'|$.
\end{lemma}
\begin{proof}
Let $r(s)$ count pairs $(a,b)\in A^2$ with sum $s$, and set $D=\{s:r(s)\ge n/(2K)\}$. Since $\sum_s r(s)=n^2$, one has $|D|\le2Kn$. The contribution of $s\notin D$ to $\sum_s r(s)^2$ is at most $n^3/(2K)$. Consequently the graph on two copies of $A$, with edge $(a,b)$ when $a+b\in D$, has at least $n^2/(2K)$ edges: each popular contribution obeys $r(s)^2\le nr(s)$. Color this graph by $a+b$. The coloring is proper by cancellation in the group.

Apply Lemma \ref{lem:four-paths} with $\eta=1/(2K)$. For each element of $A'-A'$ choose one ordered endpoint pair realizing it. Every four-step path between that pair has an ordered color tuple whose alternating sum is that difference. The paths for one pair have distinct color tuples, and tuples for different differences cannot coincide. Thus
\[
 |A'-A'|\,c\eta^5n^3\le |D|^4\le(2Kn)^4.
\]
It follows that $|A'-A'|\le CK^9n\le C'K^{10}|A'|$, as claimed.
\end{proof}
\bp{BSG}{THT.Additive.energySmallDifference}{four-paths}

\begin{lemma}[Difference-set growth from minimal growth]\label{lem:sumset-growth}
If $A$ is finite nonempty and $|A-A|\le K|A|$, $K\ge1$, then
\[
 |mA-nA|\le K^{m+n}|A|\qquad(m,n\ge0).
\]
For finite nonempty sets $X,B,C$, one also has the triangle inequality
\[
 |B-C|\,|X|\le |B-X|\,|X-C|.
\]
\end{lemma}
\begin{proof}
For the triangle inequality choose, for each $d\in B-C$, one pair $(b_d,c_d)$ with $d=b_d-c_d$. The map
\[
 (d,x)\longmapsto(b_d-x,x-c_d)
\]
is injective: the sum recovers $d$, and the chosen representative then recovers $x$. Count its domain and codomain.

Choose nonempty $X\subset A$ minimizing $r=|X-A|/|X|$. Then $r\le K$. For a finite list of translates $C=\{c_1,\ldots,c_t\}$, let $X_j$ consist of the $x\in X$ for which $x+c_j$ is not in an earlier translate $X+c_i$. The number of new points in $(X-A)+c_j$ is at most
\[
 |X-A|-|(X\setminus X_j)-A|\le r|X|-r|X\setminus X_j|=r|X_j|.
\]
The inequality follows from minimality when the complement is nonempty and is immediate when it is empty. To justify subtracting its entire sumset, if $x+c_j=x'+c_i$ for $i<j$, then $x-a+c_j=x'-a+c_i$ for every $a\in A$. Summing the new-point bounds proves
\[
 |X-A+C|\le r|X+C|.
\]
Induction with $C=-(m-1)A$ gives $|X-mA|\le K^m|X|$ for all $m\ge0$. Apply the triangle inequality with $B=mA$ and $C=nA$. Its right side is at most $K^{m+n}|X|^2$, and $|X|\le|A|$, proving the assertion.
\end{proof}
\bp{sumset-growth}{THT.Additive.minimalGrowthIteration}{finite minimal choice, injection counting}

For $G=\Z/q\Z$, a set $S\subset G$ of frequencies and $0<\rho<1/2$, write
\[
 B(S,\rho)=\{x\in G:\|sx/q\|_{\R/\Z}<\rho\text{ for every }s\in S\}.
\]
A zero frequency can be removed. A repeated frequency has no effect.

\begin{lemma}[Bohr cardinality and a spectrum generating set]\label{lem:Chang}
For $|S|=d$ and $0<\rho<1/2$,
\[
 |B(S,\rho)|\ge(\rho/3)^d|G|.
\]
If $X\subset G$ has density $\beta>0$ and $0<\eta\le1$, its spectrum
\[
 \{\xi:|\widehat{\mu_X}(\xi)|\ge\eta\},\qquad \mu_X=\one_X/\beta,
\]
is contained in the set of sums $\sum_{\lambda\in\Lambda}\epsilon_\lambda\lambda$, $\epsilon_\lambda\in\{-1,0,1\}$, for a set $\Lambda$ with
\[
 |\Lambda|\le 4\eta^{-2}\log(1/\beta).
\]
When $\beta=1$, the spectrum consists of zero and $\Lambda$ is empty.
\end{lemma}
\begin{proof}
Partition each circle into $\lceil2/\rho\rceil$ half-open intervals of length at most $\rho/2$. One fiber of the map $x\mapsto(sx/q)_{s\in S}$ has at least $|G|/\lceil2/\rho\rceil^d\ge(\rho/3)^d|G|$ points. Subtract one of these points. All resulting points belong to $B(S,\rho)$, proving the first bound.

Choose a maximal dissociated subset $\Lambda$ of the nonzero spectrum: a sum with coefficients in $\{-1,0,1\}$ vanishes only when all its coefficients are zero. Maximality implies the asserted containment, for otherwise adjoining another spectral character could not create such a relation. Choose unit complex numbers $\omega_\lambda$ so that $\E_X\operatorname{Re}(\omega_\lambda e(\lambda x/q))\ge\eta$. Put $t=\eta/2$ and
\[
 R(x)=\prod_{\lambda\in\Lambda}\bigl(1+t\operatorname{Re}(\omega_\lambda e(\lambda x/q))\bigr).
\]
Every factor is positive. Dissociation and finite Fourier orthogonality imply $\E_G R=1$, because only the constant selection in the product expansion has frequency zero. For $|u|\le1$, $0\le t\le1/2$, one has $\log(1+tu)\ge tu-t^2$. This follows, for example, by differentiating in $u$ and checking the endpoints of the possible minimum. Hence
\[
 \E_X\log R\ge |\Lambda|(t\eta-t^2)=|\Lambda|\eta^2/4.
\]
Jensen and positivity give $\beta\exp(\E_X\log R)\le\beta\E_XR\le1$. Taking logarithms proves the size bound.
\end{proof}
\bp{Chang}{THT.Additive.spectrumDissociatedSpan}{finite-fourier, Jensen, finite maximal subsets}

\begin{lemma}[Many almost periods from sampling]\label{lem:almostperiods}
Let $A\subset G$ have density $\alpha>0$, let $p\ge2$ and $0<\epsilon<1$, and put $h=f*\mu_A$. There is a set $X\subset G$ of density
\[
 \beta\ge\alpha^n/2,\qquad n=\lceil C p\epsilon^{-2}\rceil,
\]
such that
\[
 \|h(\cdot+t-t')-h\|_p\le2\epsilon\|f\|_p\qquad(t,t'\in X).
\]
For every $u\in kX-kX$ the corresponding bound is $2k\epsilon\|f\|_p$.
\end{lemma}
\begin{proof}
Choose $C$ so that Lemma \ref{lem:finite-fourier} and Markov's inequality make at least half the tuples $(a_1,\ldots,a_n)\in A^n$ satisfy
\[
 \left\|n^{-1}\sum_i f(\cdot-a_i)-h\right\|_p\le\epsilon\|f\|_p.
\]
There are $|G|^{n-1}$ cosets of the diagonal subgroup in $G^n$. One such coset, written $(z_1,\ldots,z_n)+(t,\ldots,t)$, contains at least $\alpha^n|G|/2$ good tuples. Let $X$ consist of their parameters $t$. The empirical functions for $t,t'$ are translates of one another. Translate the first approximation, subtract the second, and use invariance of the $L^p$ norm under translation to obtain the asserted bound. For a sum of $k$ differences, telescope the successive translations and use the triangle inequality.
\end{proof}
\bp{almostperiods}{THT.Additive.samplingAlmostPeriods}{finite-fourier, finite cosets, Markov}

\begin{lemma}[A quasipolynomial dense Bogolyubov lemma]\label{lem:dense-Bohr}
There is an absolute constant $C$ with the following property. If $A\subset\Z/q\Z$ has density $\alpha>0$, then $2A-2A$ contains $B(S,\rho)$, where
\[
 |S|\le C(1+\log(1/\alpha))^4,
 \qquad \rho^{-1}\le C\alpha^{-1}(1+\log(1/\alpha))^4.
\]
In particular its Bohr set has density at least $\exp(-C(1+\log(1/\alpha))^5)$.
\end{lemma}
\begin{proof}
Put $B=A+A$, $h=\mu_A*\one_{-B}$ and $F=\mu_A*h$. Then $F\ge0$, $F(0)=1$, and $\supp F\subset2A-2A$. Choose
\[
 p=\max(2,\log(1/\alpha)),\quad
 k=\lceil 10+\log_2(1/\alpha)\rceil,\quad \epsilon=(1000 e k)^{-1}.
\]
Lemma \ref{lem:almostperiods}, applied to $f=\one_{-B}$, supplies $X$ of density $\beta\ge\alpha^n/2$, with $n\le Cp\epsilon^{-2}$. Convolution with $\mu_A$ and H\"older give, for $u\in kX-kX$,
\[
 \|F(\cdot+u)-F\|_\infty
 \le\|\mu_A\|_{p'}\|h(\cdot+u)-h\|_p
 \le2k\epsilon\alpha^{-1/p}\le1/100.
\]
Here $\alpha^{-1/p}\le e$ and $\|\one_{-B}\|_p\le1$. Average these translations, setting
\[
 F'=F*\mu_X^{*k}*\mu_{-X}^{*k}.
\]
The averaging measure is supported on $kX-kX$, so $\|F-F'\|_\infty\le1/100$. Finite Fourier identities give
\[
 \widehat{F'}(\xi)=\widehat F(\xi)|\widehat{\mu_X}(\xi)|^{2k},
 \qquad \sum_\xi|\widehat F(\xi)|
 \le\sum_\xi|\widehat{\mu_A}(\xi)|^2=\alpha^{-1}.
\]
The inequality uses $\widehat F=\widehat{\mu_A}^2\widehat{\one_{-B}}$ and $|\widehat{\one_{-B}}|\le1$.

Use Lemma \ref{lem:Chang} with spectral threshold $1/2$. It supplies $S$ of size $d\le16\log(1/\beta)\le C(1+\log(1/\alpha))^4$. Put $\rho=\alpha/(1000(d+1))$. For $u\in B(S,\rho)$, every character in that spectrum satisfies $|e(\xi u/q)-1|\le2\pi d\rho$. The contribution of the spectrum to $|F'(u)-F'(0)|$ is therefore at most $2\pi d\rho/\alpha<1/100$. The other characters contribute at most $2\alpha^{-1}2^{-2k}<1/100$. It follows that $F(u)\ge1-4/100>0$, proving the containment. The displayed radius estimate and the cardinality assertion of Lemma \ref{lem:Chang} give the last conclusion. If $d=0$, the same proof uses the whole group as its Bohr set.
\end{proof}
\bp{dense-Bohr}{THT.Additive.denseBogolyubov}{almostperiods, Chang, finite-fourier}

\begin{lemma}[A derivative graph with additive energy]\label{lem:derivative-graph}
Let $|f|\le1$ on $G=\Z/q\Z$ and $\|f\|_{U^3(G)}\ge\delta>0$. There are $H\subset G$ of density at least $\delta^8/2$ and frequencies $\xi(h)\in G$ such that
\[
 \left|\E_x f(x+h)\overline{f(x)}e(-\xi(h)x/q)\right|\ge\delta^4/\sqrt2\qquad(h\in H).
\]
The graph $\Gamma=\{(h,\xi(h)):h\in H\}$ has additive energy at least $c\delta^{64}|G|^3$.
\end{lemma}
\begin{proof}
By Lemma \ref{lem:finite-fourier},
\[
 \delta^8\le\E_h\sum_\xi|\widehat{\Delta_h f}(\xi)|^4
 \le\E_h\sup_\xi|\widehat{\Delta_h f}(\xi)|^2,
 \qquad \Delta_hf(x)=f(x+h)\overline{f(x)}.
\]
The last step uses $\sum_\xi|\widehat{\Delta_hf}|^2=\|\Delta_h f\|_2^2\le1$. Let $H$ be the set on which the supremum squared is at least $\delta^8/2$, and choose a maximizing frequency. Since the supremum is at most one, $|H|/|G|\ge\delta^8/2$.

Set $v_h(k)=\one_H(h)e(\xi(h)k/q)$ and $A(x,k)=f(x+k)\overline{f(x)}$. Expansion of a squared derivative correlation gives the nonnegative quantity
\[
 T=\E_{h,x,k}A(x,k)\overline{A(x+h,k)}v_h(k)
 =\E_h\one_H(h)|\widehat{\Delta_h f}(\xi(h))|^2
 \ge\delta^{16}/4.
\]
Cauchy--Schwarz in $(x,k)$, using $|A|\le1$, gives
\[
 T^2\le\E_{x,k}\left|\E_h\overline{A(x+h,k)}v_h(k)\right|^2.
\]
Expand the square, write the second $h$ as $h+h_1$, and substitute $y=x+h$. The right side is
\[
 \E_{y,k,h_1} \overline{A(y,k)}A(y+h_1,k)
       \E_h v_h(k)\overline{v_{h+h_1}(k)}.
\]
A second Cauchy--Schwarz, now in $(y,k,h_1)$, removes the bounded $A$ factors and yields
\[
 T^4\le\E_{k,h_1}\left|\E_hv_h(k)\overline{v_{h+h_1}(k)}\right|^2.
\]
Expand again, writing the second $h$ as $h+h_2$, and integrate in $k$. Character orthogonality makes this exactly the normalized count of triples $(h,h_1,h_2)$ for which all four points lie in $H$ and
\[
 \xi(h)-\xi(h+h_1)-\xi(h+h_2)+\xi(h+h_1+h_2)=0.
\]
These are precisely the additive parallelograms of $\Gamma$; the count is its additive energy divided by $|G|^3$. Since $T^4\ge\delta^{64}/256$, the assertion follows.
\end{proof}
\bp{derivative-graph}{THT.Gowers.derivativeGraphEnergy}{finite-fourier, Cauchy--Schwarz twice}

\begin{lemma}[A locally additive derivative frequency]\label{lem:local-frequency}
Suppose $q\ge3$ is odd and $f$ satisfies Lemma \ref{lem:derivative-graph}. There are a set $H'\subset G$ of density $\alpha\ge c\delta^C$, frequencies $\xi(h)$ on $H'$ retaining its derivative correlations, a set $S\subset G$ containing $1$, and $0<\rho_0<1/100$ such that
\[
 |S|+\log(1/\rho_0)\le C(1+\log(1/\delta))^C.
\]
There is a map $M:B(S,\rho_0)\to G$, with $M(0)=0$, such that $2M$ is Freiman additive of order two there, and hence
\[
 M(u+v)=M(u)+M(v)
 \quad\text{whenever }u,v,u+v\in B(S,\rho_0).
\]
Whenever $h,h'\in H'$ and $h-h'\in B(S,\rho_0)$,
\[
 \xi(h)-\xi(h')=2M(h-h').
\]
\end{lemma}
\begin{proof}
Apply Lemma \ref{lem:BSG} to the derivative graph, using its energy lower bound and $|\Gamma|\le|G|$. It gives a subgraph $\Gamma'$ of size at least $c\delta^C|G|$ with $|\Gamma'-\Gamma'|\le K|\Gamma'|$, where $K\le C\delta^{-C}$. Define
\[
 D=\{\chi\in G:(0,\chi)\in8\Gamma'-8\Gamma'\}.
\]
The map $\Gamma'\times D\to9\Gamma'-8\Gamma'$, $(a,\chi)\mapsto a+(0,\chi)$, is injective because $\Gamma'$ is a graph. Lemma \ref{lem:sumset-growth} therefore gives $|D|\le K^{17}$.

For a nonzero $\chi\in G$, at most half the $g\in G$ satisfy $\|\chi g/q\|_{\R/\Z}<1/8$. Indeed these values are uniform on the subgroup of order $m\ge2$ generated by $\chi/q$, and the integers with distance less than $m/8$ from $0$ modulo $m$ number at most $m/2$. The latter inequality follows directly for $2\le m\le7$ and from $m/4+1\le m/2$ for $m\ge4$. Choose independently $d_0=\lceil\log_2(2|D|)\rceil$ such $g$. The expected number of nonzero elements of $D$ passing all tests is less than one, so some choice $g_1,\ldots,g_{d_0}$ admits none.

Partition the torus of values $(\xi(h)g_j/q)_j$ into half-open cubes of side $1/256$. One cube contains a subgraph $\Gamma''$ with $|\Gamma''|\ge256^{-d_0}|\Gamma'|\ge c\delta^{C'}|G|$. If $(0,\chi)\in8\Gamma''-8\Gamma''$, its evaluation at each $g_j$ has distance less than $16/256<1/8$ from zero: subtract the common cube origin before summing. It also belongs to $D$. Thus $\chi=0$, and $4\Gamma''-4\Gamma''$ is a graph.

Let $H'$ be the projection of $\Gamma''$ onto its first coordinate. Lemma \ref{lem:dense-Bohr} supplies a Bohr set in $2H'-2H'$. Adjoin frequency $1$ and shrink its radius below $1/100$; the displayed complexity bound persists. For $u$ in this Bohr set define $F(u)$ to be the unique second coordinate over $u$ in $2\Gamma''-2\Gamma''$. Uniqueness follows from the graph property already proved. The point $(0,0)$ belongs to this set, so $F(0)=0$. If $u_1+u_2=u_3+u_4$, the two frequency sums belong to $4\Gamma''-4\Gamma''$ over the same first coordinate and hence agree. This proves the Freiman property. Since $q$ is odd, multiplication by two is invertible in $G$; define $M=2^{-1}F$. The local additivity follows by taking $(u_1,u_2,u_3,u_4)=(u,v,u+v,0)$. Finally $(h-h',\xi(h)-\xi(h'))$ belongs to $\Gamma''-\Gamma''\subset2\Gamma''-2\Gamma''$, giving the last identity.
\end{proof}
\bp{local-frequency}{THT.Gowers.locallyAdditiveFrequency}{derivative-graph, BSG, sumset-growth, dense-Bohr}

\begin{lemma}[Smooth weights and local Fourier decay]\label{lem:Bohr-weight}
Fix a nonnegative Gevrey function $b$, supported in $(-1,1)$, equal to one on $[-1/2,1/2]$, and at most one. For $S\subset G$, $d=|S|\ge1$, and $0<r<1/8$, define the periodic weight
\[
 w_{S,r}(x)=\prod_{s\in S} b\bigl(\langle sx/q\rangle/r\bigr),\qquad
 \beta_{S,r}=\E_xw_{S,r}(x),\qquad \mu_{S,r}=w_{S,r}/\beta_{S,r},
\]
where $\langle t\rangle\in[-1/2,1/2)$ is the centered representative. Then
\[
 (r/6)^d\le\beta_{S,r}\le1,\qquad
 |w_{S,r}(x+t)-w_{S,r}(x)|\le C d\|t\|_S/r,
 \quad \|t\|_S=\max_{s\in S}\|st/q\|_{\R/\Z}.
\]
Suppose a map $\phi$ into $\R/\Z$ is additive whenever the arguments and their sum lie in a larger Bohr set containing $B(S,3r)$. If
\[
 \left|\E_x\mu_{S,r}(x)e(\phi(x))\right|\ge\tau>0,
\]
then, for $t\in B(S,r)$,
\[
 \|\phi(t)\|_{\R/\Z}\le C d\|t\|_S/(r\beta_{S,r}\tau).
\]
Only values on the specified domain are used in these averages.
\end{lemma}
\begin{proof}
The weight is one on $B(S,r/2)$, so Lemma \ref{lem:Chang} gives its mass lower bound. Each periodically extended factor is globally Lipschitz on the circle with constant $C/r$: it vanishes in a neighborhood of the cut in the centered coordinate. Telescope the difference of the product; all other factors have absolute value at most one. This proves the translation bound.

Put $A=\E_xw_{S,r}(x)e(\phi(x))$. Substitution and local additivity on the union of the two translated supports give
\[
 \E_xw_{S,r}(x-t)e(\phi(x))=e(\phi(t))A.
\]
Subtract the expression for $A$ and use the translation bound. This gives $|e(\phi(t))-1|\beta_{S,r}\tau\le Cd\|t\|_S/r$. The elementary inequality $|e(u)-1|\ge4\|u\|_{\R/\Z}$ proves the conclusion. The domain contains all points in the two supports and all sums used to justify the identity, so no extension of $\phi$ outside its domain has been assumed.
\end{proof}
\bp{Bohr-weight}{THT.Additive.smoothBohrWeight}{gevrey, Chang, character identities}


\section{A finite inverse theorem with a complete proof}\label{sec:finite-inverse}
The next argument supplies the structural theorem rather than importing a nilsequence inverse theorem. The finite sampling, spectrum, graph, and lattice ingredients have all been proved above. Multiplication and addition inside a phase with denominator $q$ are taken in $G=\Z/q\Z$; their images in $\R/\Z$ are independent of integer representatives. In this section $q\ge3$ is odd.

\begin{lemma}[Local symmetry]\label{lem:local-symmetry}
Let $S\subset G$ contain $1$, let $d=|S|$, and let $M:B(S,\rho_0)\to G$ be locally additive, with $M(0)=0$. Take $0<r_1\le\rho_0/100<1/100$ and $w_1=w_{S,r_1}$. Suppose $a,b,c:G\to\C$ are bounded by one and
\[
 \left|\E_{h,x} w_1(h)c(h)a(x+h)b(x)e(-2M(h)x/q)\right|\ge\zeta>0.
 \eqtag{eq:local-derivative-correlation}
\]
There are $S'\supset S$ and $0<r'<r_1$ such that
\[
 |S'|+\log(1/r')\le C\bigl(d+\log(1/r_1)+\log(1/\zeta)+1\bigr)^C
\]
and
\[
 \|\{u,v\}\|_{\R/\Z}\le\zeta/1000
 \quad(u,v\in B(S',r')),\qquad
 \{u,v\}=(M(u)v-M(v)u)/q.
\]
All arguments of $M$ used in the conclusion and its proof belong to its stated domain.
\end{lemma}
\begin{proof}
Choose
\[
 0<r_2\le\min(r_1/100,\zeta^2r_1/(100C d)),
\]
where $C$ exceeds the translation constant in Lemma \ref{lem:Bohr-weight}. Set $w_2=w_{S,r_2}$, $\beta_2=\E w_2$, and $\mu_2=w_2/\beta_2$. Average the substitution $x=y+u$, $u$ distributed with weight $\mu_2$, in \eqref{eq:local-derivative-correlation}. Averaging also in uniform $y$ recovers the original expression. For some $y$ the resulting correlation has magnitude at least $\zeta$. Absorb the translate into $a,b$ and the phase $e(-2M(h)y/q)$ into $c$. Thus
\[
 \left|\E_h w_1(h)c(h)\E_u\mu_2(u)a(u+h)b(u)e(-2M(h)u/q)\right|\ge\zeta.
\]
Cauchy--Schwarz in $h$, using $\E w_1\le1$ and $|c|\le1$, and expansion of the square imply
\[
 \zeta^2\le\E_{h,u,v}w_1(h)\mu_2(u)\mu_2(v)
 b(u)\overline{b(v)}a(u+h)\overline{a(v+h)}
 e(-2M(h)(u-v)/q).
\]
This expression is real and nonnegative. Put $z=h+u+v$. Whenever the weight is nonzero, all of $h,u,v,z,u+v$ lie in $B(S,\rho_0)$, since $r_1+2r_2<\rho_0$. Local additivity gives $M(h)=M(z)-M(u)-M(v)$. Define, on the relevant values,
\[
 A_z(u)=b(u)\overline{a(z-u)}e\bigl((2M(u)u-2M(z)u)/q\bigr).
\]
The expression on the right of the preceding inequality becomes
\[
 \E_{z,u,v}\mu_2(u)\mu_2(v)w_1(z-u-v)
 A_z(u)\overline{A_z(v)}e(-2\{u,v\}).
\]
The definitions may be extended by zero where $z$ is outside the domain, because both weights in the following comparison vanish there. Replace $w_1(z-u-v)$ by $w_1(z)$. Lemma \ref{lem:Bohr-weight} bounds the change by $2Cd r_2/r_1\le\zeta^2/2$, since both $\mu_2$ averages have mass one. Some $z$ in the support of $w_1$ therefore satisfies
\[
 \left|\E_{u,v}\mu_2(u)\mu_2(v) A_z(u)\overline{A_z(v)}e(-2\{u,v\})\right|\ge\zeta^2/2.
\]
A Cauchy--Schwarz inequality in $u$ removes $A_z(u)$. Expand the resulting square, use local additivity on $B(S,3r_2)$, and discard the bounded remaining factors by the triangle inequality. This gives
\[
 \E_{v,v'}\mu_2(v)\mu_2(v')
 \left|\E_u\mu_2(u)e(-2\{u,v-v'\})\right|\ge\zeta^4/4.
\]
Fix $v'$ for which the average in $v$ has at least this value. With $\tau=\zeta^4/8$, the set
\[
 A=\left\{v\in\supp w_2:
 \left|\E_u\mu_2(u)e(-2\{u,v-v'\})\right|\ge\tau\right\}
\]
has $\mu_2$-mass at least $\tau$. Indeed the integrand is at most one, so an average at least $2\tau$ cannot have mass less than $\tau$ above $\tau$. Its ordinary density is at least $\beta_2\tau$, because $w_2\le1$.

For $z=v-v'$ with $v\in A$, apply Lemma \ref{lem:Bohr-weight} to the locally additive phase $x\mapsto-2\{x,z\}$. It yields, for $x\in B(S,r_2)$,
\[
 \|2\{x,z\}\|_{\R/\Z}\le D\|x\|_S,
 \qquad D=Cd/(r_2\beta_2\tau).
\]
Every element of $2A-2A$ is a sum of two differences of elements of $A-v'$. Successive partial sums lie in $B(S,8r_2)\subset B(S,\rho_0)$. Local additivity and the triangle inequality on the circle consequently give
\[
 \|2\{x,z\}\|_{\R/\Z}\le4D\|x\|_S
 \quad(x\in B(S,r_2),\ z\in2A-2A).
\]
Use Lemma \ref{lem:dense-Bohr} to find $B(S_C,\rho_C)\subset2A-2A$. Replace $S_C$ by $S_C\cup S$ and choose
\[
 0<r_C\le\min(\rho_C,r_2,\zeta/(10000D))/100.
\]
For $u,v\in C:=B(S_C,r_C)$ the preceding inequality applies. Because doubling is invertible on $G$,
\[
 2C=B(2^{-1}S_C,r_C).
\]
If $x=2u$ and $y=2v$, local additivity gives $\{x,y\}=4\{u,v\}$, hence
\[
 \|\{x,y\}\|_{\R/\Z}\le2\|2\{u,v\}\|_{\R/\Z}
 \le8Dr_C\le\zeta/1000.
\]
Set $S'=S\cup2^{-1}S_C$ and $r'=r_C$. Its Bohr set is contained in $2C$, giving the conclusion without dividing an inequality on the circle by two.

Finally $\log(1/\beta_2)\le d\log(6/r_2)$ by Lemma \ref{lem:Bohr-weight}. All of $\log(1/r_2)$, $\log(1/\tau)$, and $\log D$ are polynomial in the parameters in the statement. The density of $A$ is at least $\beta_2\tau$. The bounds in Lemma \ref{lem:dense-Bohr}, followed by the displayed choices of $r_C$ and $S'$, prove the asserted complexity bound using Lemma \ref{lem:budget-rules}.
\end{proof}
\bp{local-symmetry}{THT.Gowers.localSymmetry}{Bohr-weight, dense-Bohr, Cauchy--Schwarz}

\begin{lemma}[A locally quadratic phase is a bounded chart on a smaller Bohr set]\label{lem:local-chart}
Let $S\subset G$ contain $1$, let $D=|S|\ge1$, and let $M:B(S,r_0)\to G$ be locally additive, with $0<r_0<1/100$. Put $\phi(x)=M(x)x/q$. For
\[
 0<r\le\epsilon_D r_0/(2\sqrt D),\qquad \epsilon_D=(100D)^{-10D},
\]
there is a bounded real-variable chart atom $\psi$ of the form \eqref{eq:atom}, with rational chart matrix of height one and complexity at most
\[
 \exp\bigl(C(D+\log(1/r)+1)^C\bigr),
\]
such that, at every integer $n$,
\[
 \psi(n)=w_{S,r}(n)e(\phi(n))
\]
when the weight is nonzero, and $\psi(n)=0$ otherwise. In particular the complexity does not depend on $q$ or on any choice of lattice basis.
\end{lemma}
\begin{proof}
List $S=(1,s_2,\ldots,s_D)$ with integer representatives $0\le s_i<q$, and set $\theta=(1/q,s_2/q,\ldots,s_D/q)$. The full lattice
\[
 \Lambda=\Z\theta+\sum_{i=2}^D\Z e_i
\]
contains $\Z^D$: $q\theta=(1,s_2,\ldots,s_D)$, so $e_1$ also belongs to it. Its dual is therefore a subgroup of $\Z^D$. The homomorphism $\pi:\Lambda\to G$ sends $n\theta+\sum_{i\ge2}k_ie_i$ to $n\pmod q$; its kernel is $\Z^D$.

Apply Lemma \ref{lem:lattice-box} to $r_0^{-1}\Lambda$ and rescale. It gives an integer coefficient box $P$ in part of a full basis $b_1,\ldots,b_D$ of $\Lambda$, contained in $B_2(0,r_0/4)$ and containing $\Lambda\cap B_2(0,\epsilon_Dr_0)$. All its partial coordinate sums remain in this ball. The map $\pi$ is injective on $P$: a difference of two such points has norm less than $1/2$, and a nonzero vector of its kernel $\Z^D$ has norm at least one. Every point of $P$, and every partial sum, projects into $B(S,r_0)$, since its coordinates are representatives of the corresponding characters and are less than $r_0$ in absolute value.

Let $v_i=\pi(b_i)$ for the basis vectors occurring in the box. Local additivity, applied successively to the coordinate path and to its negative steps, proves
\[
 M\left(\sum_i n_i v_i\right)=\sum_i n_iM(v_i)
 \quad\text{for every coefficient vector in the box}.
\]
Thus $\phi(\pi(\sum_i n_ib_i))$ is, modulo one, the ordinary quadratic polynomial
\[
 q^{-1}\left(\sum_i n_i M(v_i)\right)\left(\sum_jn_jv_j\right),
\]
using any integer representatives of the displayed group elements.

Let $u_i$ be the dual basis vectors of the full lattice basis. They belong to $\Lambda^*\subset\Z^D$, and $u_i\cdot\theta$ is an integer. At an integer $n$ with $w_{S,r}(n)\ne0$, let $k\in\Z^D$ be the unique vector for which $w=\theta n-k$ has all its coordinates in $(-r,r)$. This lift belongs to $\Lambda$: subtracting $k_1q\theta$ changes its first coordinate by $k_1$ and changes the others by integers. Its norm is at most $\sqrt D r\le\epsilon_Dr_0/2$, so it belongs to $P$. Its basis coefficients are the integers
\[
 n_i=u_i\cdot w=(u_i\cdot\theta)n-\sum_j(u_i)_j k_j.
\]
Substitute these integer linear expressions into the quadratic polynomial above. Every monomial is an integer at integer $(n,k)$. Reducing each of its real coefficients modulo one therefore leaves its exponential unchanged on all these points. After this reduction the polynomial has the form
\[
 P(n,k)=\lambda n^2+\beta n+c+n b^Tk+k^TC_0k+d_0^Tk,
\]
where the entries of $b,C_0,d_0$ have absolute value at most one; an upper triangular $C_0$ suffices. Constant and unused linear terms can be zero. All coefficients before this reduction may have been arbitrarily large; none is placed in a smooth amplitude before reducing it.

Substituting $k=\theta n-w$ rewrites this polynomial as
\[
 \lambda' n^2+\beta'n+c+n v^Tw+w^TC_0w-d_0^Tw,
 \qquad v=-b-(C_0+C_0^T)\theta.
\]
Use the enlarged horizontal vector $(\theta,v)$ and the rational block matrix
\[
 C=\begin{pmatrix}0&0\\ I_D&0\end{pmatrix}.
\]
Its chart phase is $n v^Tw$. With $\rho$ from Lemma \ref{lem:gevrey}, take amplitude
\[
 \chi(w,w')=\left(\prod_{j=1}^D b(w_j/r)\right)
 e(w^TC_0w-d_0^Tw)\prod_{j=1}^D\rho(w'_j).
\]
The lattice sum in the $w'$ variables equals one. In the $w$ variables there is at most one nonzero term at every real argument, because $r<1/2$. Consequently the resulting atom, with outside phase $\lambda'y^2+\beta'y+c$, is bounded by one on the entire real line. The calculation proves its required values at integers.

There are $2D$ horizontal coordinates, the chart matrix is integral of height one, the support is bounded, and the only coefficients in the amplitude have absolute value at most one. The product rule, the polynomial phase derivative formula, and Lemma \ref{lem:gevrey} bound derivatives of order $m$ by
\[
 (C(D+1)/r)^{C(m+1)}(m+1)^{C(m+1)}.
\]
This gives the displayed complexity after increasing the fixed Gevrey order once. Neither the lattice basis nor $q$ occurs in this bound.
\end{proof}
\bp{local-chart}{THT.Gowers.locallyQuadraticChart}{lattice-box, integer-basis, gevrey, chart-operations}

\begin{theorem}[Finite bounded chart inverse theorem]\label{thm:finite-inverse}
There is an absolute $C$ such that, for every odd $q\ge3$, every $0<\delta<1/2$, and every $1$-bounded function $f:G\to\C$ with $\|f\|_{U^3(G)}\ge\delta$, there is a real-variable bounded quadratic chart atom $a$ of complexity at most
\[
 H_\delta=\exp\bigl(C(1+\log(1/\delta))^C\bigr)
\]
for which
\[
 \left|\E_{n\in G}f(n)\overline{a(n)}\right|\ge H_\delta^{-1}.
\]
The atom is evaluated at integer representatives $0,\ldots,q-1$. Its real frequencies are unrestricted. There is no bound on $q$ in either conclusion.
\end{theorem}
\begin{proof}
Obtain $H',S,\rho_0,M$ from Lemma \ref{lem:local-frequency}, and let $\alpha=|H'|/|G|$. Put $r_1=\rho_0/100$, $w_1=w_{S,r_1}$, and $\beta_1=\E w_1$. Averaging in $x_0$ gives
\[
 \E_{x_0,h}w_1(h)\one_{H'}(x_0+h)=\alpha\beta_1.
\]
Choose $x_0$ for which the inner average is at least this number. Select one $h_*$ of positive weight with $x_0+h_*\in H'$. For every other such $h$, both $h-h_*$ and $h,h_*$ lie in the domain of $M$. Lemma \ref{lem:local-frequency} and local additivity imply
\[
 \xi(x_0+h)=\xi_0+2M(h),\qquad
 \xi_0=\xi(x_0+h_*)-2M(h_*).
\]
Choose $c(h)$ of modulus one on these $h$ to make the associated derivative correlations nonnegative, and set it zero on the other $h$. With
\[
 a(x)=f(x+x_0),\quad b(x)=\overline{f(x)}e(-\xi_0x/q),\quad
 \zeta=\delta^4\alpha\beta_1/\sqrt2,
\]
the hypothesis \eqref{eq:local-derivative-correlation} holds. The inverse of $\zeta$ is quasipolynomial in $\delta^{-1}$, since $\beta_1\ge(r_1/6)^{|S|}$.

Use Lemma \ref{lem:local-symmetry} to obtain $S'\supset S$ and $r'$, with skew phase bounded by $\zeta/1000$ on $B(S',r')$. Write $D=|S'|$. The map $M$ is locally additive on $B(S',r_0)$ with $r_0=\rho_0/10$, because $S'\supset S$. Choose
\[
 0<r_4\le\min\left(r'/100,\ \epsilon_Dr_0/(2\sqrt D),\ \zeta r_1/(100C|S|)\right),
\]
where $C$ dominates the translation bound. Put $w_4=w_{S',r_4}$, $\beta_4=\E w_4$, and $\mu_4=w_4/\beta_4$. All ranks and inverse radii chosen so far have quasipolynomial bounds in $\delta^{-1}$; in particular $\beta_4^{-1}$ does also. Set $\phi(x)=M(x)x/q$ on the domain of $M$.

Average \eqref{eq:local-derivative-correlation} after substituting $h+t$ for $h$, with $t$ distributed according to $\mu_4$. Replacing $w_1(h+t)$ by $w_1(h)$ changes the result by at most $C|S|r_4/r_1\le\zeta/100$. On the relevant supports, $M(h+t)=M(h)+M(t)$. Thus some $h_0$ in the support of $w_1$ satisfies
\[
 \left|\E_{t,x}\mu_4(t)c(h_0+t)a(x+h_0+t)
 \overline{f(x)}e(-\xi'x/q)e(-2M(t)x/q)\right|\ge\zeta/2,
 \qquad \xi'=\xi_0+2M(h_0).
\]
The existence follows because the outer weight $w_1$ has mass at most one. Next substitute $x=y+u$ and average $u$ according to $\mu_4$ and $y$ uniformly. Some $y$ gives the same lower bound. Absorb the translate and $e(-2M(t)y/q)$ into bounded functions $a_y,c_y$, and discard an overall constant phase. The resulting expression is
\[
 \left|\E_{u,t}\mu_4(u)\mu_4(t)c_y(t)a_y(u+t)
 \overline{f(y+u)}e(-\xi'u/q)e(-2M(t)u/q)\right|\ge\zeta/2.
\]
For $u,t$ in the two supports the skew bound permits replacing $2M(t)u/q$ by $(M(u)t+M(t)u)/q$ with absolute error at most $2\pi\zeta/1000$. Local additivity gives the exact polarization identity
\[
 \phi(u+t)-\phi(u)-\phi(t)=(M(u)t+M(t)u)/q\pmod1.
\]
It follows that the last expression, after this replacement, has magnitude at least $\zeta/4$ and equals $\E_{u,t}F(u)B(t)A(u+t)$, where
\[
 F(u)=\mu_4(u)\overline{f(y+u)}e(\phi(u)-\xi'u/q),\qquad
 B(t)=\mu_4(t)c_y(t)e(\phi(t)),
\]
\[
 A(v)=a_y(v)e(-\phi(v))\quad(v\in B(S',2r_4)),
 \qquad A(v)=0\quad\text{otherwise}.
\]
No value of $\phi$ outside its domain is required: $u+t\in B(S',2r_4)\subset B(S,\rho_0)$.

Finite Fourier inversion and Cauchy--Schwarz give
\[
 \zeta/4\le\left|\sum_\xi\widehat F(-\xi)\widehat B(-\xi)\widehat A(\xi)\right|
 \le\sup_\xi|\widehat F(\xi)|\,\|B\|_2\|A\|_2
 \le\beta_4^{-1/2}\sup_\xi|\widehat F(\xi)|.
\]
The last inequality uses $|A|\le1$ and $\E\mu_4^2\le\beta_4^{-1}$, since $0\le w_4\le1$. For some $\ell\in G$, multiplying this conclusion by $\beta_4$ and conjugating proves
\[
 \left|\E_u w_4(u)f(y+u)e(-\phi(u)+\ell u/q)\right|
 \ge\zeta\beta_4^{3/2}/4.
\]
Lemma \ref{lem:local-chart}, followed by multiplication by a linear phase and translation of the real argument, gives a bounded chart $a(n)$ representing $w_4(n-y)e(\phi(n-y)-\ell(n-y)/q)$ at all integers. The expression just obtained is its correlation with $f$. Both its complexity and the reciprocal correlation are quasipolynomial in $\delta^{-1}$. There have been a fixed number of operations in Lemma \ref{lem:budget-rules}; their degrees are absolute. Increasing one absolute constant gives the stated common $H_\delta$.
\end{proof}
\bp{finite-inverse}{THT.Gowers.finiteChartInverse}{local-frequency, local-symmetry, local-chart, Bohr-weight, finite-fourier, chart-operations}


\section{Continuous witnesses and a bounded decomposition}\label{sec:continuous}
For a function on $[0,1)$, extend it by zero to $\T_{16}=\R/(16\Z)$ with Haar probability and define
\[
 \|g\|_U^8=\int_{\T_{16}^4}\prod_{\omega\in\{0,1\}^3}
 \mathcal C^{|\omega|}g(x+\omega\cdot h)\,d\mu(x)\,d\mu(h).
\]
Here $\mathcal C$ denotes complex conjugation. For a general interval $I$, first normalize affinely to $[0,1)$ to define $U(I)$. Lemma \ref{lem:cube-foundation} proves that this is a seminorm, proves all mixed cube bounds, and gives $\|g\|_{U(I)}\le C\|g\|_{2,I}$ with a fixed absolute normalization constant.

\begin{lemma}[Continuous inverse and decomposition]\label{lem:inverse}
A $1$-bounded $g$ on $I$ with $\|g\|_{U(I)}\ge\delta$, $0<\delta<1/2$, has a normalized moment of magnitude at least $H_\delta^{-1}$ against a bounded atom of complexity at most $H_\delta$, where $H_\delta$ is quasipolynomial in $\delta^{-1}$. For $0<\tau<1/2$, a $1$-bounded $f$ on $I$ consequently has a finite decomposition
\[
 f=\sum_\nu c_\nu\phi_\nu+u+v,\qquad
 \sum_\nu|c_\nu|\le D_\tau,\qquad
 \|u\|_{U(I)}\le\tau,\qquad \|v\|_{1,I}\le\tau,
 \eqtag{eq:decomp}
\]
where $u,v\in L^2(I)$ and the $\phi_\nu$ are bounded atoms of complexity at most $D_\tau$, quasipolynomial in $\tau^{-1}$.
\end{lemma}
\begin{proof}
Normalize $I=[0,1)$. For an integer $L\ge1$, set
\[
 g_L(j)=L\int_{j/L}^{(j+1)/L}g(y)\,dy\quad(0\le j<L)
\]
and extend by zero to the odd cyclic group $G_L=\Z/(16L+1)\Z$. Its raw order-three cube average converges to the continuous cube average above. Here are the details of this discretization. In a nonzero supported discrete cube, choose the representative $x\in\{0,\ldots,L-1\}$ and represent each increment by the difference of $x+h_i\in\{0,\ldots,L-1\}$ and $x$. Then $|h_i|\le L-1$, and every real integer vertex lies in $[-3L,4L]$. On this range it is congruent to an integer in $[0,L)$ modulo $16L+1$ exactly when it already lies there. Thus there are no wraparound cubes. Dividing all coordinates by $L$, the count is a Riemann sum over a fixed bounded box, with normalizing factor $(L/(16L+1))^4\to16^{-4}$.

For a continuous function supported in $(0,1)$, its cell averages differ from endpoint values by its modulus of continuity, so this Riemann sum converges to the corresponding real cube integral. The same no-wraparound argument on $\T_{16}$ identifies that real integral divided by $16^4$ with the required Haar average. For a general $1$-bounded $g$, approximate in $L^1([0,1))$ by continuous functions supported in $(0,1)$ and bounded by one. Such approximation follows from the imported $L^1$ density, multiplication by interior cutoffs, and radial projection onto the unit disk, which does not increase distance from a point of that disk. Telescoping the eight bounded factors changes a cube average by at most eight times the $L^1$ distance on its underlying probability group. Cell averaging contracts the $L^1$ distance and introduces only the bounded factor $L/(16L+1)$. Therefore the approximation error is uniform in $L$, proving convergence for $g$.

Choose $L$ large enough that $\|g_L\|_{U^3(G_L)}\ge\delta/2$. No quantitative bound on $L$ is needed. Theorem \ref{thm:finite-inverse} gives a bounded real chart $a$ with controlled complexity and
\[
 \left|L^{-1}\sum_{j=0}^{L-1}g_L(j)\overline{a(j)}\right|\ge\kappa,
\]
where $\kappa^{-1}$ is quasipolynomial in $\delta^{-1}$: changing the cyclic normalization to $L^{-1}$ multiplies the correlation by $(16L+1)/L\ge1$.

An exact interpolation is required, since the sampled chart may oscillate arbitrarily rapidly. Write its parameters as $(\lambda,\beta,\theta,C,\chi)$ and clear the denominators of $C$ by $P$. At integer $n$, write $\theta=\theta_0+Pz$ with $z\in\Z^d$ and $0\le(\theta_0)_i<P$, and reindex the chart sum by $Pzn$. This replaces $\lambda$ by $\lambda+(Pz)^TC\theta_0$; the remaining extra term is integral. Reduce this new coefficient and $\beta$ modulo one. Hence the same sampled atom has a representation with $0\le\lambda,\beta<1$ and $0\le\theta_i<P$. This representation need only have the same integer values; those values remain bounded by one.

Set $\alpha=2\lambda+\theta^TC\theta$, enlarge the horizontal vector to $\Theta=(\theta,1,\alpha)$, extend $C$ by two zero rows and columns, and set its new entry in row $d+2$, column $d+1$, to $-1$. Choose a Gevrey function $\eta$, supported in $(0,1)$, between zero and one, equal to one outside endpoint strips of total length at most $2\epsilon_1$, and a partition $\sum_{\ell\in\Z}\rho_0(r-\ell)=1$ as in Lemma \ref{lem:gevrey}. Use the amplitude
\[
 \widetilde\chi(w,h,r)=\chi(w-\theta h)\eta(h)\rho_0(r)
 e\bigl(-h\theta^TCw+(\lambda+\theta^TC\theta)h^2-\beta h\bigr).
\]
The resulting atom $A(v)$, with outside phase $\lambda v^2+\beta v$, satisfies exactly
\[
 A(v)=\eta(\{v\})a(\lfloor v\rfloor).
 \eqtag{eq:floor}
\]
Indeed support forces the lift $h=v-m$ to have $m=\lfloor v\rfloor$, the $r$ partition sums to one, and the remaining phase is
\[
 \lambda v^2+\beta v+(\theta v)^TCw-(2\lambda+\theta^TC\theta)vh
 -h\theta^TCw+(\lambda+\theta^TC\theta)h^2-\beta h.
\]
Substituting $m=v-h$ and $w=\theta v-n$ makes this $\lambda m^2+\beta m+(\theta m)^TC(\theta m-n)$. At integer $v$, both sides of \eqref{eq:floor} are zero. Every parameter in the new amplitude is bounded by the old complexity and the cutoff scale $\epsilon_1^{-1}$, because $\lambda,\beta,\theta$ were reduced first. Its complexity is therefore quasipolynomial in those quantities. Identity \eqref{eq:floor} also proves $|A|\le1$ on the whole real line.

Exact cell averaging identifies the displayed discrete correlation with the integral against $a(\lfloor Ly\rfloor)$. Replacing this function by $A(Ly)$ changes the integral by at most $2\epsilon_1$. Take $\epsilon_1\le\kappa/8$. Scaling by $L$ changes only unrestricted real frequencies, so the resulting continuous witness has the asserted complexity. Lemma \ref{lem:chart-operations} treats a general interval $I$.

For the decomposition, all norms in the rest of the proof use probability measure on $I$. Put $\sigma=\tau/2$. Let $\mathcal A$ be the bounded atoms given by the continuous inverse at threshold $\sigma^2$, with correlation lower bound $\kappa>0$, and choose $M\ge\max(1,2\sigma/\kappa)$. Consider the balanced convex set in complex $L^2(I)$
\[
 \mathcal B=M\operatorname{absconv}(\mathcal A)
 +\{u:\|u\|_U\le\sigma\}
 +\{v:\|v\|_1\le\sigma\}.
\]
The absolute convex hull consists of finite complex combinations with sum of coefficient magnitudes at most one. The $U$ ball contains an $L^2$ neighborhood of zero by Lemma \ref{lem:cube-foundation}.

If $f$ were outside the $L^2$ closure of $\mathcal B$, real Hahn--Banach separation and Hilbert-space representation of the separating functional would give $h\in L^2$ which, after positive normalization, satisfies
\[
 \operatorname{Re}\langle f,h\rangle>1,
 \qquad \sup_{z\in\mathcal B}\operatorname{Re}\langle z,h\rangle\le1.
\]
The normalizing supremum is finite by separation and positive because $\mathcal B$ contains a neighborhood of zero. Testing the three summands, using their balancedness, gives
\[
 \|h\|_\infty\le\sigma^{-1},\qquad
 |\langle u,h\rangle|\le\sigma^{-1}\|u\|_U\ (u\in L^2),\qquad
 \sup_{\phi\in\mathcal A}|\langle h,\phi\rangle|\le M^{-1}.
\]
For the first assertion bounded simple functions in the $L^1$ ball suffice: a positive-measure set where $|h|>\sigma^{-1}$, with its phase approximated by finitely many values, would contradict the pairing bound. For the second, rescale $u$ when its $U$ norm is positive; when that norm is zero all its scalar multiples belong to the $U$ ball, forcing its pairing to be zero. The third follows by using a single atom with coefficient of magnitude $M$.

Since $\|f\|_2\le1$, the strict inequality implies $\|h\|_2>1$. Test the second bound with $u=h$ to obtain $\|h\|_U\ge\sigma\|h\|_2^2>\sigma$. Thus $\sigma h$ is $1$-bounded and has $U$ norm greater than $\sigma^2$. Its inverse witness gives $|\langle h,\phi\rangle|\ge\kappa/\sigma>M^{-1}$, a contradiction. Consequently $f$ belongs to the closure of $\mathcal B$. Approximate it by one member of $\mathcal B$ within $\tau/2$ in $L^2$, and add this error to $v$. Its $L^1$ cost is at most $\tau/2$, so the total $L^1$ error is at most $\tau$. All terms remain in $L^2$, the structured sum is finite, and the quantitative assertions follow from the inverse bounds and Lemma \ref{lem:budget-rules}.
\end{proof}
\bp{inverse}{THT.Chart.continuousInverseAndDecomposition}{finite-inverse, cube-foundation, gevrey, chart-operations, library Hahn--Banach and Riesz}


\section{Weyl splitting and progression separation}\label{sec:charts}
\begin{lemma}[A short-return criterion for a rational splitting]\label{lem:return-splitting}
Let $1\le k\le3$, $d\le\log H$, $H\ge e^2$, let $C$ be a rational $d$ by $d$ matrix of height at most $H$, and let $V\in\R^{d\times k}$ and symmetric $\Lambda\in\R^{k\times k}$ be unrestricted. Put $A=C-C^T$, $C_{\rm sy}=(C+C^T)/2$, and $C_{\rm sk}=A/2$. Let $P\ge1$ be an integer with $PC$ integral. Suppose $0<a<1$, $M\ge1$, and
\[
 \|(2\Lambda+V^TCV)h-V^TAm\|\le M
 \quad\text{whenever }m\in P\Z^d,\quad \|h\|+\|Vh-m\|\le a.
\]
There is a rational subspace $U\subset\R^d$ such that, with its orthogonal projection $P_f$, $P_s=1-P_f$, $V_f=P_fV$, and $V_s=P_sV$,
\[
 A|_{U\times U}=0,\qquad \|V_s\|\le H',
\]
\[
 \left\|\Lambda+\tfrac12V^TC_{\rm sy}V+\sym(V_f^TC_{\rm sk}V_s)\right\|\le H'.
\]
Both projection heights and $H'$ are at most
\[
 \exp\!\left(C(1+\log H+\log P+\log(1/a)+\log M)^C\right).
\]
The constant and degree are absolute.
\end{lemma}
\begin{proof}
Write $\Omega(h,m)=(2\Lambda+V^TCV)h-V^TAm$ and $d_h=Vh-m$. For two returns, direct multiplication and use of $A^T=-A$ give
\[
 h'^T\Omega(h,m)-h^T\Omega(h',m')=d_{h'}^TAd_h-m'^TAm.
\]
Replace $a$ by $a_1=\min(a,[100(M+\|A\|+1)]^{-1})$. This changes its reciprocal by at most a quasipolynomial in the stated parameters. The absolute value of the first two terms on the right after solving for $m'^TAm$ is at most $a_1^2\|A\|+2a_1M<1$. That pairing is an integer because $m,m'\in P\Z^d$ and $PC$ is integral. Therefore it vanishes for all returns of size at most $a_1$.

For $d=0$, every $h$ is a return and $\Omega=2\Lambda h$. Take $h=a_1e_j/2$ for each coordinate to bound the columns of $\Lambda$, and use the zero subspace. Now assume $d\ge1$. The set
\[
 K=P^{-1}V[-1,1]^k+[-1,1]^d
\]
is compact, symmetric and convex, and has nonempty interior. Let $C_d=(100d)^{10d}$ and $L=16C_d$. Choose $b_{\max}=\min(1/100,a_1/[100P(d+k+1)])$. Of the $d+1$ disjoint intervals
\[
 (b_{\max}L^{-j-1},b_{\max}L^{-j}],\qquad0\le j\le d,
\]
at least one contains no successive minimum of $K$ for $\Z^d$, because there are only $d$ minima. Let $b_0$ be its left endpoint and set $U=\operatorname{span}_{\R}(\Z^d\cap b_0K)$. Every $n$ in this intersection has a representation $n=P^{-1}Vh+w$ with $\|h\|_\infty,\|w\|_\infty\le b_0$. Thus $m=Pn$ and $d_h=-Pw$ form a return of size below $a_1$. The preceding integer identity makes the pairings of these generators zero, so $U$ is isotropic for $A$.

Put $r=\dim U$. If $r<d$, the empty-minimum interval implies $\lambda_{r+1}>Lb_0$. Lemma \ref{lem:transference} supplies $d-r$ independent integer vectors $q_i$ in $(C_d/(Lb_0))K^\circ$. Their pairings with the generating vectors of $U$ have absolute value at most $C_d/L<1$ and are integral. They are zero, so the $q_i$ span $U^\perp$. Since $[-1,1]^d\subset K$, each coordinate of $q_i$ is bounded by $C_d/(Lb_0)$. Since $P^{-1}V[-1,1]^k\subset K$, every coordinate of $q_i^TV$ is bounded by $PC_d/(Lb_0)$. Let $S$ have the $q_i^T$ as rows. Then
\[
 P_s=S^T(SS^T)^{-1}S.
\]
Lemma \ref{lem:rational-linear} bounds its rational height, its coefficients, and $P_sV$, using the bounds for $S$ and $SV$. All their logarithms are polynomial in $d,\log H,\log P,\log(1/a),\log M$, because $b_0^{-1}\le b_{\max}^{-1}L^{d+1}$. If $r=d$, put $P_s=0$. In either case put $P_f=1-P_s$.

For each coordinate $e_j\in\R^k$, take more than $(1+4/b_0)^d$ equally spaced samples of $P^{-1}V(te_j)$ modulo $\Z^d$, with $t\in[0,b_0]$. Divide the torus into cubes of side at most $b_0/4$. Two samples in the same cube give a return
\[
 h=t_je_j,\quad b_0/M_0\le t_j\le b_0,\quad
 \|d_h\|_\infty\le Pb_0/4,\quad M_0\le2(1+4/b_0)^d.
\]
Its lattice vector $P^{-1}m$ lies in $b_0K$, hence $m\in U$. Thus $\Delta=m-V_fh=-P_fd_h\in U$. Isotropy gives $V^TA\Delta=V_s^TA\Delta$, so this term is bounded by quantities already controlled, without a bound on $V_f$. Put
\[
 \Lambda_* =\Lambda+\tfrac12V^TC_{\rm sy}V+\sym(V_f^TC_{\rm sk}V_s).
\]
Expansion with $V=V_f+V_s$ and $V_f^TC_{\rm sk}V_f=0$ gives the identity
\[
 2\Lambda+V^TCV-V^TAV_f=2\Lambda_*+V_s^TC_{\rm sk}V_s.
\]
Insert $m=V_fh+\Delta$ into $\Omega$. Its assumed bound now bounds each column $\Lambda_*t_je_j$ in terms of the controlled $V_s,C,d_h$ and $M$. Finally $t_j^{-1}\le M_0/b_0$ is quasipolynomial in the same parameters. This proves the effective quadratic bound and all stated quantitative assertions by Lemma \ref{lem:budget-rules}.
\end{proof}
\bp{return-splitting}{THT.Chart.shortReturnSplitting}{transference, rational-linear, chart-definition, finite torus pigeonhole, budget-rules}

\begin{lemma}[Weyl alternative]\label{lem:weyl}
Let $\vartheta$ be supported in a fixed box with supremum and Lipschitz norms at most $H$. If
\[
 \sup_{\xi\in\R^k}\left|\int\vartheta(z)Q(z)e(-\xi^Tz)\,dz\right|\ge\delta>0,
\]
there is a rational subspace $U\subseteq\R^d$, with orthogonal projections $P_f$ onto $U$ and $P_s=1-P_f$, such that
\[
 A|_{U\times U}=0,\qquad \|V_s\|\le H',\qquad \|\Lambda_*\|\le H',
 \eqtag{eq:weyl-splitting}
\]
where $V_f=P_fV$, $V_s=P_sV$ and
\[
 \Lambda_*=\Lambda+\tfrac12V^TC_{\rm sy}V+\sym(V_f^TC_{\rm sk}V_s).
 \eqtag{eq:effective}
\]
The projection heights and $H'$ are quasipolynomial in $H,\delta^{-1}$.
\end{lemma}
\begin{proof}
Put $F=\vartheta Q$ and let $D$ be a quasipolynomial bound for its support size, supremum, $L^2$ norm, and first absolute moment. For a return $Vh=m+d_h$, $m\in P\Z^d$, where $P$ clears $C$, direct reindexing gives
\[
 F(z+h)=e(c(h,m))e(\Omega(h,m)^Tz)F(z)+E_{h,m}(z),
\]
\[
 \Omega(h,m)=(2\Lambda+V^TCV)h-V^TAm,
 \qquad \|E_{h,m}\|_1+\|E_{h,m}\|_\infty\le D(\|h\|+\|d_h\|).
 \eqtag{eq:return}
\]
Indeed, with $u=Vz$, $w=u-n$, $d=d_h$, the nonconstant chart increment is $u^TCd+m^TCu+d^TCw$. The first two terms give the displayed modulation, since $m^TCn$ is integral; the last is bounded by $D\|d\|$ on the amplitude support. Translation of the amplitudes costs $D(\|h\|+\|d\|)$.

The Fourier transform of $F$ is $D$-Lipschitz. Choose a point where its magnitude is at least $3\delta/4$, put $\rho=\delta/(100D)$, and choose an integer $M\ge100D^2\delta^{-2}\rho^{-k}$. Take $a\le\min\{\delta/(100DM),\,[100(1+\|A\|)]^{-1}\}$. For a return of size at most $a$, its first $M$ multiples in \eqref{eq:return} produce Fourier values at least $\delta/2$ at points spaced by $\Omega(h,m)$. If $\|\Omega\|\ge1$, the radius-$\rho$ balls about these points are disjoint and have Fourier magnitude at least $\delta/3$. Their $L^2$ mass contradicts Plancherel and the choice of $M$. Thus $\|\Omega(h,m)\|<1$ on these returns. Apply Lemma \ref{lem:return-splitting} with this return threshold and $M=1$. The clearing period can be chosen $P\le H^{d^2}$ by taking the product of the denominators of $C$. Thus $\log P$ and $\log(1/a)$ are polynomial in $\log H+\log(1/\delta)$. That lemma gives \eqref{eq:weyl-splitting} and \eqref{eq:effective} with the asserted quantitative bounds, including $d=0$.
\end{proof}
\bp{weyl}{THT.Chart.weylSplitting}{chart-operations, fourier-foundation, return-splitting}


\begin{lemma}[Fourier expansion with uniform envelopes]\label{lem:chartfourier}
If a rational splitting obeys \eqref{eq:weyl-splitting} with height and size bound $B$, then on a fixed box
\[
 Q(z)=e(z^T\Lambda_*z+b^Tz)\sum_\ell a_\ell e(\sigma_\ell^TVz)
 \eqtag{eq:chartfourier}
\]
exactly and absolutely. There are quasipolynomially many families with at most $2d$ integer indices; their rational frequencies have height at most $D(1+\|\ell\|)$, and $\sum|a_\ell|\le D$, with $D$ quasipolynomial in $H,B$. For every fixed $J$, the absolute coefficient sum in a dyadic shell $T\le1+\|\ell\|<2T$ is at most
\[
 D_JT^{-J(d+1)^2}.
 \eqtag{eq:shell}
\]
This remains true for the coefficient envelope pooled over all splittings of height at most $B$, separately for each resulting frequency and quadratic part, uniformly over the allowed amplitudes and bounded slow values. A uniform tail $\epsilon$ therefore requires only a quasipolynomial truncation in $H,B,\epsilon^{-1}$.
\end{lemma}
\begin{proof}
Write $u=Vz$, $u_f=P_fu$, $u_s=P_su$, $w=u-n$. The identities
\[
 u^TC_{\rm sy}(u-n)=\tfrac12u^TC_{\rm sy}u+\tfrac12w^TC_{\rm sy}w-\tfrac12n^TC_{\rm sy}n,
\]
\[
 u^TC_{\rm sk}w=u_f^TC_{\rm sk}u_s-u_f^TC_{\rm sk}P_sn+u_s^TC_{\rm sk}w
\]
extract exactly the quadratic in \eqref{eq:effective}. Choose an even period $P_0$ clearing $C_{\rm sy}$ and a denominator $q_0$ for $P_s$. Since $u_s$ and $w$ are bounded, split into the quasipolynomially many possibilities $P_sn=N\in q_0^{-1}\Z^d$ and $n\equiv a\pmod{P_0\Z^d}$. The residue phase is fixed and the remaining linear phase is $-u_f^TC_{\rm sk}N$.

Let $\eta(s)$ equal one on the relevant slow values and let $E(0)=1$ be a smooth cutoff supported in the cube of radius $1/(3q_0)$. The remaining amplitude is the restriction at $s=P_su$ of
\[
 \eta(s)\sum_{n\in a+P_0\Z^d}\chi(u-n)E\bigl(s-P_s(u-n)-N\bigr)
 e\bigl(\tfrac12(u-n)^TC_{\rm sy}(u-n)+s^TC_{\rm sk}(u-n)\bigr).
\]
The equality cutoff is exact by lattice spacing. This function is $P_0$-periodic in $u$ and, after periodizing the compact slow cutoff in a larger box, $P_1$-periodic in $s$, with controlled integer $P_1$. Its Fourier characters are $e(p^Ts/P_1+\ell^Tu/P_0)$. Substitution gives
\[
 \sigma=\sigma_N+P_s^Tp/P_1+\ell/P_0,
 \qquad \sigma_N^Tu=-(P_fu)^TC_{\rm sk}N,
\]
proving the identity and frequency heights.

At derivative order $m$, the rescaled periodic amplitudes are bounded by $D_0^{C(m+d+1)}(m+1)^{C(m+1)}$, uniformly in the stated parameters. This follows by the product rule: differentiating a quadratic exponential produces a polynomial of degree at most $m$, and all its amplitude coefficients, periods and cutoff scales are bounded by $D_0$. For a Fourier index $v\ne0$, differentiate $m$ times in a coordinate of magnitude at least $\|v\|/\sqrt{2d}$. Its coefficient is at most
\[
 D_0^{C(m+d+1)}(m+1)^{C(m+1)}(2d)^{m/2}(1+\|v\|)^{-m}.
\]
A shell has at most $(CT)^{2d}$ indices. Taking $m\ge J(d+1)^2+C(d+1)^2$ proves \eqref{eq:shell}; its prefactor remains quasipolynomial because $d\le\log H$. There are at most $B^{C(d+1)^2}$ rational splittings and quasipolynomially many $(N,a)$. Pool their uniform derivative envelopes and, for coinciding frequencies, sum envelopes. Increasing $m$ by $C(d+1)^2$ absorbs any bounded-height frequency count as well. Geometric summation proves the absolute norm and truncation assertions. When $d=0$ the expansion has one term.
\end{proof}
\bp{chart-fourier}{THT.Chart.splitFourierExpansion}{weyl, chart-operations, fourier-foundation, rational-linear}


\begin{lemma}[Progression separation]\label{lem:separation}
Two bounded atoms in distinct slots $i,k$, of complexity $H$, have an exact expansion
\[
 \phi_i(y_i)\phi_k(y_k)=\sum c\,\psi_a(y_a)\psi_b(y_b)Q(t),
 \eqtag{eq:separation}
\]
where $a,b$ are the other slots. All factors are bounded atoms fixed independently of scale. The terms have shells $T\ge1$ with absolute coefficient mass at most $DT^{-100}$, dimension at most $\log D$, and complexity at most $(DT)^C$, for $D$ quasipolynomial in $H$.

For one term, let $\zeta$ be a fixed nonnegative smooth cutoff equal to one on the allowed $t/r$ range, and put
\[
 A_r=\sup_{\xi\in\R}\left|\int Q(ru)\zeta(u)e(-\xi u)\,du\right|.
\]
Its form norm in the remaining two normalized $L^2(I)$ slots is at most $CA_r$. If $A_r\ge\delta>0$, the two-variable kernel on $I\times I$ has a uniform $\epsilon$-approximation by a finite sum $\sum c_\nu\varphi_{\nu,a}(y_a)\varphi_{\nu,b}(y_b)$ of bounded atoms, with quasipolynomial number, complexity and absolute coefficient norm in $H,\delta^{-1},\epsilon^{-1}$. Before the smooth kernel is separated the phases are
\[
 \lambda_*t^2+\gamma t,\qquad |\lambda_*|r^2\le D,
 \eqtag{eq:separated-phase}
\]
where $\lambda_*$ is a bounded-height rational combination of $\lambda$ and $\theta\otimes\theta$, and $\gamma-\beta$ of $\theta$.
\end{lemma}
\begin{proof}
First separate $\phi_i$ alone. Choose fixed rationals $\alpha_i=1$, $\alpha_k=0$, $\sum\alpha_j=\sum j\alpha_j=0$. Insert a nonnegative Gevrey partition of unity at $\theta y_a$, $\theta y_b$, $\theta t$, with lifts $d_a,d_b,d$. For the original lift $d_i$, the carries $k_j=d_j-d_a-(j-a)d$ are integers. Their supports bound the possible $k_i,k_b$, and specifying them determines the original index; a smooth equality cutoff in $d_b-d_a-(b-a)d-k_b$, equal to one at zero and zero at all other integers, imposes the second relation exactly. With $S=\sum j^2\alpha_j$,
\[
 \sum_j\alpha_j\bigl(\lambda y_j^2+(\theta y_j)^TCd_j\bigr)
 =S\bigl(\lambda t^2+(\theta t)^TCd\bigr)+\sum_j\alpha_j(\theta y_j)^TCk_j.
\]
The last sum, and the original linear phase, distribute as linear phases in $y_a,y_b$. The remaining coupled smooth amplitude in $(d_a,d_b,d)$ is compactly supported with controlled Gevrey bounds. Periodize inside a larger product cutoff and Fourier-expand. Each term factors into three chart amplitudes, with matrices $-\alpha_aC,-\alpha_bC,SC$. Apply the same procedure to $\phi_k$ with coefficient of slot $i$ zero, then multiply the resulting triples. The integration-by-parts estimate in Lemma \ref{lem:chartfourier}, with order exceeding the Fourier dimension by at least $110$, gives the asserted shell bound. Normalizing factors to be bounded only alters the quasipolynomial coefficient budget.

For the single-scale norm use coordinates $(y_a,y_b)$, with Jacobian $|b-a|$ and $t=(y_b-y_a)/(b-a)$. The rescaled smooth kernel has an absolutely summable two-variable Fourier expansion, whose terms only modulate the inputs. The remaining convolution kernel $Q(t)\zeta(t/r)$ has Fourier multiplier bounded by $CrA_r$. Plancherel and the normalization $r^{-2}$ prove the norm bound.

If $A_r\ge\delta$, apply Lemmas \ref{lem:weyl} and \ref{lem:chartfourier} to $Q(r\cdot)$. This gives \eqref{eq:separated-phase} with the indicated rational dependencies. Distribute $e(\gamma t)$ between $y_a,y_b$, and Fourier-expand $w_I(x,t)e(\lambda_*t^2)$ with product cutoffs. Its normalized derivatives satisfy $D^{C(m+1)}(m+1)^{C(m+1)}$, so a quasipolynomial truncation gives the required error and coefficient norm. Each separated factor is $\psi_a$ or $\psi_b$ times a linear phase and a bounded smooth bump at scale $r$. Such a bump on $I$ is itself a zero-matrix chart with one extra slow coordinate: map $I$ strictly inside a lattice cell and support its amplitude away from the other translates. Outside $I$ use its bounded periodic extension. Divide each successive error by preceding coefficient norms. Only a fixed number of quasipolynomial operations is involved.
\end{proof}
\bp{separation}{THT.Chart.progressionSeparation}{chart-fourier, chart-operations, gevrey, absolute}




\section{Detection from a local form}\label{sec:detection}
\begin{lemma}[A generalized von Neumann inequality]\label{lem:von-neumann}
For any ordering of four distinct slots, with $R,g\in L^2(I)$ and $b,c\in L^\infty(I)$,
\[
 |H_I(R,g,b,c)|\le C\|R\|_{2,I}\|g\|_{U(I)}\|b\|_\infty\|c\|_\infty.
 \eqtag{eq:von-neumann}
\]
For $R,g\in L^1(I)$ the bound with $\|R\|_{1,I}\|g\|_{1,I}$ in place of the first two factors also holds.
\end{lemma}
\begin{proof}
The $L^1$ assertion is the two-coordinate case of Lemma \ref{lem:absolute}. For the first, normalize $I=[0,1)$ and extend the smooth kernel periodically on a sufficiently large fixed box; the interior support margin allows an extension equal to zero near its boundary. Lemma \ref{lem:fourier-foundation} gives a uniformly absolutely summable Fourier series. Every character of $(x,t)$ can be distributed between slots $0,1$, since $x=y_0$ and $t=y_1-y_0$. Character multiplication preserves the stated $L^p$ norms and the cube norm, by cancellation of its third multiplicative difference. It suffices to handle the unweighted progression average on a fixed real torus. Zero extension and the support margin identify its relevant real values without wraparound.

Let slots $j,i,k,l$ contain $R,g,b,c$ respectively. In coordinates $u=x+jt,t$, Cauchy--Schwarz first gives
\[
 |\Lambda|^2\le\|R\|_2^2
 \E_u\left|\E_tg(u+(i-j)t)b(u+(k-j)t)c(u+(l-j)t)\right|^2.
\]
Expand the square with the second $t$ equal to $t+h_1$. In the original coordinates the remaining expression is the average of
\[
 g(y_i)\overline{g(y_i+(i-j)h_1)}\,
 b(y_k)\overline{b(y_k+(k-j)h_1)}\,
 c(y_l)\overline{c(y_l+(l-j)h_1)}.
\]
Disintegrate now along $(h_1,y_k)$, with $t$ as free coordinate. The $b$ factors are independent of that coordinate and bounded by $\|b\|_\infty^2$. Cauchy--Schwarz in $(h_1,y_k)$ and expansion of the inner square introduce $h_2$ and replace each remaining factor by its product with its conjugate translate in direction $(-k,1)h_2$. The $g$ factors now occupy the four vertices with increments $(i-j)h_1,(i-k)h_2$, and the $c$ factors the corresponding vertices with increments $(l-j)h_1,(l-k)h_2$. Repeat in $(h_1,h_2,y_l)$, bounding the four $c$ factors by $\|c\|_\infty^4$, and expand with free-direction increment $(-l,1)h_3$. Combining the three inequalities gives
\[
 |\Lambda|^8\le\|R\|_2^8\|b\|_\infty^8\|c\|_\infty^8
 \E_{x,h_1,h_2,h_3}\prod_{\omega\in\{0,1\}^3}
 \mathcal C^{|\omega|}g\bigl(x+\omega_1(i-j)h_1+\omega_2(i-k)h_2+\omega_3(i-l)h_3\bigr).
\]
Each displayed nonzero integer multiplication preserves Haar probability, and the three increments remain independent. The last factor is exactly $\|g\|_{U^3}^8$. This proves the assertion for bounded inputs. Lemma \ref{lem:cube-foundation} and the $L^2$ local form bound give passage to arbitrary $L^2$ inputs by truncation. Sum the kernel Fourier coefficients and restore the fixed interval normalization.
\end{proof}
\bp{von-neumann}{THT.LocalForm.vonNeumannOneL2}{absolute, cube-foundation, fourier-foundation, torus-annihilator}

\begin{lemma}[Local detection]\label{lem:detection}
Suppose four inputs in any order satisfy
\[
 \|R\|_{2,I}\le C_1,\quad \|h\|_{q,I}\le1,\quad
 \|u\|_{2,I},\|v\|_{2,I}\le1,\quad |H_I(R,h,u,v)|>\delta.
\]
For $q>2$, $C_1\ge1$ and $0<\delta<1/2$, $R$ has a normalized moment of magnitude at least $D^{-1}$ against a bounded atom of complexity $D$, where
\[
 D\le\exp\bigl(C_q(1+\log C_1+\log(1/\delta))^C\bigr)
\]
with absolute degree $C$.
\end{lemma}
\begin{proof}
Put $q'=q/(q-1)$. Clip the three opponents at height $B$. Their tails satisfy
\[
 \|h\one_{|h|>B}\|_2\le B^{-(q-2)/2},\qquad
 \|u\one_{|u|>B}\|_{q'},\|v\one_{|v|>B}\|_{q'}\le B^{-(q-2)/q}.
\]
Use \eqref{eq:absolute} with exponents $(2,2,2,2)$ for the first replacement and $(2,q,q',2)$ for the others. Taking $B=\max(2,(100CC_1/\delta)^{q/(q-2)})$ leaves a clipped form above $3\delta/4$.

Decompose $h_B/B$ by \eqref{eq:decomp} at $\tau_1=\delta/(100CC_1B^3)$. Its structured part after restoring $B$ has coefficient norm $A_1=BD_{\tau_1}$. Both errors cost at most $CC_1B^3\tau_1$ by \eqref{eq:von-neumann}, so the form with this structured part is above $\delta/2$. Decompose $u_B/B$ at $\tau_2=\delta/(100CC_1A_1B^2)$, giving coefficient norm $A_2=BD_{\tau_2}$. Summing over the first structured list, each new error costs $CC_1A_1B^2\tau_2$. Some bounded atom pair consequently has form magnitude above $\delta_0=\delta/(4A_1A_2)$ against $R,v_B$, where $\|v_B\|_2\le1$.

Expand this pair by Lemma \ref{lem:separation}. Each coefficient-free term has absolute form at most $CC_1$. Truncate shells so $CC_1D_0T_0^{-100}<\delta_0/4$. Some retained term has magnitude at least $\delta_1=\delta_0/(2D_0)$; its multiplier norm forces $A_r\ge\delta_1/(CC_1)$. Separate its kernel to uniform error $\delta_1/(10CC_1)$. The remaining form is a sum of products of normalized moments of $R$ and $v_B$, with coefficient norm $D_1$ quasipolynomially bounded. Every $v_B$ moment is at most one. Hence one $R$ moment is at least $\delta_1/(2D_1)$. All operations form a fixed sequence of quasipolynomial constructions. The only preliminary $q$ dependence is the clipping power, which multiplies the accuracy logarithm by a $q$-dependent constant and does not change its absolute polynomial degree. Conjugate the witness when a bilinear moment convention is needed.
\end{proof}
\bp{detection}{THT.LocalForm.detectResidual}{inverse, separation, von-neumann, absolute}




\section{The base reduction}\label{sec:base}
\subsection{Two structured inputs and delayed moments}
On a subtree rooted at $J$, call a sequence in a remaining slot admissible if it is a fixed $F\in L^2(J)$, a current inherited projection $P_IF$ with column budget $D_h$, or a sequence $Z_I$ with
\[
 \sum_{I\subseteq J}|I|\|Z_I\|_{2,I}^2\le L^2\|\one_JF\|_2^2.
\]
A fixed number of these sequences may be added. The indices can be any subcollection; complete projection families on missing cells by inheritance.

\begin{lemma}[Scale sums and predictors]\label{lem:predictors}
For one fixed separated term $\psi_a(y_a)\psi_b(y_b)Q(t)$ of complexity $H$ from \eqref{eq:separation}, and admissible opponents,
\[
 \sum_{I:e^{-m-1}<A_{|I|}\le e^{-m}}|I|\,|H_I(\psi_az_{a,I},\psi_bz_{b,I};Q)|
\]
is at most
\[
 e^{-m}\poly(1+m,\log(2H),\log(2D_h),1+L)\|\one_JF_a\|_2\|\one_JF_b\|_2.
 \eqtag{eq:band-scale-sum}
\]
The term vanishes when $A_{|I|}=0$. For two fixed original atoms of complexity $H_{\rm low}$ and $B\ge1$, their product in the form can therefore be replaced by a finite list
\[
 \psi_a(y_a)\psi_b(y_b)w_I(x,t)e(\lambda_*t^2),\qquad |\lambda_*||I|^2\le D_{\rm pred},
 \eqtag{eq:predictors}
\]
with absolute length-weighted error at most
\[
 e^{-B}\poly(\log(2D_h),1+L)\|\one_JF_a\|_2\|\one_JF_b\|_2.
\]
The triples $(\psi_a,\psi_b,\lambda_*)$ are fixed on $J$, the multipliers are bounded atoms, and the list size, coefficient norm and complexities obey $\log D_{\rm pred}\le\poly(1+B,\log(2H_{\rm low}))$. Scalar coefficients may depend on scale.
\end{lemma}
\begin{proof}
Put $K=1+m+\log(2H)$. All intermediate quantities $D_i$ below have logarithms polynomial in $K$. The original bilinear norm is $CA_r$. For $A_r>e^{-m-1}$ the chart expansion gives phases $e(\lambda_*t^2+\gamma t)$, $|\lambda_*|r^2\le D_1$, whose unscaled parameters are bounded-height rational evaluations of one fixed real vector.

Truncate pooled Fourier tails with total coefficient envelope below $e^{-2m}(1+D_1)^{-C}$, with fixed reserve. This is valid in the absolute scale-sum norm, not just at one scale. For each tail pair write $e(\lambda_*t^2)=1+(e(\lambda_*t^2)-1)$. The first term has the mean-zero square-function bound established below at its fixed linear center. The second has smooth-kernel cost $|\lambda_*|r^2\poly(D_1)$ and
\[
 \sum_{r:|\lambda_*|r^2\le D_1}|\lambda_*|r^2\le CD_1.
\]
At each depth use contractions and Cauchy--Schwarz; square-energy sequences cost at most $1+L$. The chosen tail reserve absorbs these factors and the corresponding polynomial in $\log D_h$.

The finite main catalog has coefficient norm and height bounded by $D_2$. Choose a tiny $\epsilon$ after it. Declare exceptional scales where any nonzero $|\lambda_*|r^2$ or $r|\gamma-\gamma'|$ falls in an intermediate annulus between sufficiently small and sufficiently large fixed powers of $\epsilon,D_2$. Lemma \ref{lem:height} bounds their number by $\poly(K)$. At them use the original norm $CA_r\le Ce^{-m}$ before expanding coefficients. On the other scales every retained quadratic is tiny, and the linear centers form clumps of diameter at most $\epsilon/r$ separated by at least $\epsilon^{-1}/r$. To obtain the diameter take connected components at a tiny threshold divided by the catalog count. Replacing quadratics by zero and centers by clump representatives has total error at most the catalog norm times
\[
 \sum_r\max_{\rm tiny}\{r^2|\lambda_*|\}
 +\sum_r\max_{\rm tiny}\{r|\gamma-\gamma'|\},
\]
which is $\epsilon\poly(K)$ after the reserved factors, by \eqref{eq:tiny-sum}. Choose it below the desired error. Each combined clump coefficient is $O(e^{-m})$: test the original expansion against $\zeta(u)e(-r\gamma_Cu)$. Its own coefficient is multiplied by $\int\zeta>0$; other clumps give rapidly decreasing Fourier tails, and all replacements and truncations are already smaller than $e^{-m}$.

Here is the square-function estimate used for these linear terms. In coordinates $(y_a,y_b)$ the kernel has zero integral in either variable. Its double primitive is compactly supported in an interior rectangle. Fourier-expand that primitive inside slightly larger product cutoffs and differentiate once in each variable. This represents the kernel by tensor products of smooth interior mean-zero bumps, with rapidly decreasing coefficients and only polynomial derivative growth in their indices. For one such bump the rows are
\[
 W_{I,\nu}F=|I|^{-1/2}\int_Ie(\nu y)\omega_I(y)F(y)\,dy,
 \qquad \nu=\pm\gamma/(b-a).
\]
They are jointly Bessel over the retained scales, cells and clump representatives, with squared bound $\poly(K)$. In detail, restrict the input to shells $|\xi+\nu|\asymp2^\ell/|I|$. At a fixed scale and center, mean zero and Poincaré give row norm $O(2^\ell)$ for $\ell<0$; two frequency antiderivatives and integration by parts give $O(2^{-2\ell})$ for $\ell\ge0$. Sum on disjoint cells at that scale. For fixed $\xi,\ell$, Lemma \ref{lem:height}, adjoining $\xi$ and the fixed linear modulation to the parameter vector, allows only $\poly(K)$ scales. Center separation allows at most $C(1+2^\ell)$ centers in the shell at one scale. The square root of this frequency multiplicity times the row bounds is summable over $\ell$. Exact center frequencies are a null set. A fixed bounded multiplier $\psi_a$ or $\psi_b$ preserves the estimate.

For a current projection insert it by Lemma \ref{lem:compression}, paying a polynomial in $\log(2D_h)$. For a square-energy sequence use the per-cell Bessel bound over separated centers and then its assumed energy. Cauchy--Schwarz between the two row families and the $O(e^{-m})$ clump coefficients prove \eqref{eq:band-scale-sum}. If $A_r=0$, the cutoff function has zero Fourier transform and the relevant kernel is zero.

Finally expand the low pair by \eqref{eq:separation}. Sum \eqref{eq:band-scale-sum} over its rapidly decreasing shells and the Fourier-supremum bands. Discard large shells and small-$A_r$ bands at logarithmic thresholds polynomial in $B,\log H_{\rm low}$. For the retained terms truncate the pooled chart tails in the scale-sum norm just proved. Distribute each $e(\gamma t)$ between the two remaining slots. The finite choices give exactly \eqref{eq:predictors} with the asserted bounds.
\end{proof}
\bp{predictors}{THT.Base.finitePredictorList}{separation, height, compression, fourier-foundation, absolute}


\begin{lemma}[Delayed predictor replacement]\label{lem:delay}
For a fixed term of \eqref{eq:predictors}, suppose the position-$a$ projection family contains the predictor $(\overline{\psi_a},\sqrt\epsilon\,e_p)$ on every use cell and its descendants. Replace $F_a$ at a cell of depth $s$ by its projection on the partition at depth $s+K_0$. Against an admissible opponent the absolute scale-sum error is at most
\[
 C\bigl(2^{-K_0}+(1+K_0)\sqrt\epsilon\bigr)
 \poly(D_{\rm pred},\log(2D_h),1+L)\|\one_JF_a\|_2\|\one_JF_b\|_2.
 \eqtag{eq:delay-error}
\]
Extra column insertions are permitted, and both slots may be replaced successively.
\end{lemma}
\begin{proof}
Let $E_h$ be conditional expectation at depth $h$ and $P_h$ the direct-sum projection. Put $D_h=(f_a-(P_hF_a)_{\phys})\psi_a$ and $\mu_h=E_hD_h$. By \eqref{eq:predictor-moment}, $\mu_h=\epsilon a_{h,p}$, so $\mu_{h+1}-\mu_h$ is $\sqrt\epsilon$ times the change of the private $p$ coordinate. This remains true after other insertions. For a smooth bump $\omega_I$ and $H=s+K_0$, exact telescoping gives
\[
 \int_I D_H\omega_I=\int_I\mu_H\omega_I
 +\sum_{h\ge H}\int_I(D_h-D_{h+1})(\omega_I-E_h\omega_I)
 +\sum_{h\ge H}\int_I(\mu_{h+1}-\mu_h)(E_{h+1}\omega_I-E_h\omega_I).
 \eqtag{eq:delay-telescope}
\]
To verify it telescope $\int D_h(\omega_I-E_h\omega_I)$ and use $E_{h+1}D_{h+1}=\mu_{h+1}$; the term with $\mu_h(E_{h+1}-E_h)\omega_I$ vanishes. The terminal remainder tends to zero by the uniform $L^2$ bound and smooth conditional-expectation approximation; convergence of $P_hF_a$ to $F_a$ is not required.

Each infinite-sum row has the factor $O(2^{s-h})$. At fixed offset $h-s$, sum over disjoint $s$-cells and then over $s$ using orthogonality of $P_{h+1}-P_h$. Minkowski over the offsets gives $C2^{-K_0}\|F_a\|_2$, with an additional $\sqrt\epsilon$ for the private term. If $\omega_I$ has mean zero, replace the first term by $\mu_H-\mu_s$ and sum its $K_0$ private increments, giving $CK_0\sqrt\epsilon\|F_a\|_2$. Apply the mean-zero tensor decomposition from Lemma \ref{lem:predictors} to $w_I$ and pair these rows with the opposing Bessel rows, using compression or the square-energy bound as appropriate.

For $w_I(e(\lambda_*t^2)-1)$ use an ordinary smooth tensor expansion. Its derivatives contain $|\lambda_*|r^2\poly(D_{\rm pred})$, which sums geometrically. The expectation-free terms are treated by \eqref{eq:delay-telescope}; the remaining $\mu_H$ term has per-depth norm $C\sqrt\epsilon\|F_a\|_2$. This proves \eqref{eq:delay-error}. A previously delayed opponent is its current projection plus a rolling difference. For a nested family,
\[
 \sum_s\|(P_{s+K_0}-P_s)F\|_2^2\le K_0\|F\|_2^2,
\]
because each orthogonal increment is counted at most $K_0$ times. Restricting to use cells only decreases the sum. Thus the rolling difference is admissible (even the weaker $L\le C(1+K_0)$ suffices), proving sequential replacement.
\end{proof}
\bp{delay}{THT.Base.delayedPredictorReplacement}{predictors, coeff, prefix, conditional expectation}


\subsection{Adaptive fits and the two earliest increments}
Let $d_{j,I}(r)$ be the supremum of the form with $r$ in slot $j$ and, in the other slots, one normalized $L^q(I)$ input and two normalized $L^2(I)$ inputs, over all choices of the $L^q$ position.

For each accuracy $e^{-k}$ there are inherited projections $g_I=P_IF_j$ with $\|g_I\|_{2,I}\le1$ and
\[
 d_{j,I}(f_j-(g_I)_{\phys})\le e^{-k},
 \eqtag{eq:adaptive}
\]
on segments with bounded total root length. Their path counts, atom complexities and physical supremum norms are at most $\exp(k^{O(1)})$. To construct them, start with no columns and on a failing cell use Lemma \ref{lem:detection} to find a bounded atom of residual correlation $\rho$, with $\log(1/\rho)\le k^{O(1)}$. Insert $(\phi,e_v)$ with fresh coordinate and refit. Its squared energy gain is at least $c\rho^2|I|$, since the old residual has zero fresh coordinate and the new column has norm at most $\sqrt2$. Fits and refinements form nested subspaces, so total energy is at most $\|F_j\|_{L^2(J)}^2\le|J|$. Stop before the insertion count along a path exceeds $C\rho^{-2}$. Integrated energy makes the first stopped roots cost at most $|J|/4$ for sufficiently large $C$. Restart, using the convention of Section \ref{sec:reduction}. Private weights one bound the coefficient $\ell^2$ norm by one, and Cauchy--Schwarz gives the physical supremum bound. All degrees are absolute by detection.

\bp{adaptive-fit}{THT.Base.inheritedAdaptiveFit}{detection, coeff, prefix, path-packing}

Update the four positions cyclically at levels $t$, with $k_{t+1}=k_t^A$ and sufficiently large $k_1$. Denote the updated projection by $g_{t,I}$; set its previous value $g_{t-4,I}=0$ before the first update. For two earliest increment levels $u<t$ in distinct positions $i,j$, write
\[
 X_i=g_u-g_{u-4},\qquad X_j=g_t-g_{t-4},\qquad
 X_a=F_a-g_{a,\mathrm{prev}},\qquad X_b=F_b-g_{b,\mathrm{prev}},
 \eqtag{eq:base-tuple}
\]
where $a,b$ are the remaining slots and ``prev'' is their last update before $t$. At a finite cutoff, expand each input into its increments and terminal residual. Terms with at least two increments group exactly as \eqref{eq:base-tuple}; the remaining terms have either four terminal residuals or one whole terminal projection and three terminal residuals.

For updates $h>\ell$ in distinct slots, replacing their pair of projections by the actual inputs has bilinear norm error at most
\[
 C_q\bigl(e^{-k_h}\exp(C_qk_\ell^d)+e^{-k_\ell}\bigr)
 \eqtag{eq:stagger}
\]
against normalized $L^2$ spectators, for an absolute $d$. Replace the more accurate projection first against the other's bounded physical supremum, then the less accurate one against the actual normalized $L^q$ input. Choose $A>d$. This also makes any pair of terminal residuals tend to zero and proves convergence of the expansion on a finite interval collection.

For $t\ge12$, $k=k_t$, every pair of full $X$ arguments has bilinear norm at most $e^{-2k^{\theta_0}}$ for some absolute $\theta_0>0$. Here are the index cases. If $u<t-4$, a pair involving the early slot is tested directly against each of its two old projections: the later difference or residual cancels its actual-input terms, and its earliest accuracy index is at least $t-4>u$. Other pairs involve only indices at least $t-4$. If $u\ge t-4$, distinct positions imply $u\ge t-3$ and all nonzero projections involved have indices at least $t-7$. Expand a selected pair and apply \eqref{eq:stagger}; its two-actual terms cancel, since at least one full factor has coefficient sum zero after replacement. Thus the error is at most $C_qe^{-c_qk_{t-7}}$. Since $k_{t-7}=k_t^{A^{-7}}$, any fixed $\theta_0<A^{-7}$ works for large $k_1$. All tuple factors have bounded normalized $L^2$ norm. Intersections of the bounded number of segment families involved still have bounded total root length: every new common root is a root of one constituent family.

\bp{stagger}{THT.Base.twoEarliestIncrements}{adaptive-fit, absolute, budget-rules}

\begin{proposition}[Finite base data]\label{prop:base}
On further bounded-packing normalized segments there are $K_0\le k^{O(1)}$ and two new nested high projection families such that, at depth $s$,
\[
 Y_i=X_i,\quad Y_j=X_j,\qquad
 Y_e=P^{\rm new}_{e,s+K_0}F_e-g_{e,\mathrm{prev}}\quad(e=a,b),
 \eqtag{eq:Y}
\]
replace the tuple with absolute scale-sum error $Ce^{-k}$ times root length. Counts, complexities and inverse private weights are $\exp(k^{O(1)})$. Their normalized $L^2$ sizes are bounded, and every full pair has norm at most $\eta=e^{-k^\theta}$ for a fixed absolute $0<\theta<1$.

Writing $Y_l=B_l+U_l$, with $U_i=U_j=0$ and
\[
 U_{e,I}=(P^{\rm new}_{e,s+K_0}-P^{\rm new}_{e,s})F_e,
\]
each $B_l$ is a difference of at most two current basic projections, and on any normalized root $R$,
\[
 \sum_{I\subseteq R}|I|\|U_{e,I}\|_{2,I}^2\le C(1+K_0)^2|R|.
 \eqtag{eq:basic-side}
\]
Up to error $e^{-k}|R|$ and the already summable two-$U$ term, it suffices to estimate, for finitely many patterns with at least three $B$ factors,
\[
 \sum_{I\text{ active}}|I|\min\{\eta,|H_I(z)|\}.
 \eqtag{eq:active-min}
\]
The active collection has path multiplicity at most $D_{\rm base}^C$, where $e^k\le D_{\rm base}\le\exp(k^C)$ covers all the basic data. The whole complementary $B$ arguments restore to the full pair in \eqref{eq:Y} by adding basic sides $U$.
\end{proposition}
\begin{proof}
On a common older segment the two low lists have a budget $D_{\rm old}=\exp(k^{O(1)})$. Apply Lemma \ref{lem:predictors} to each original pair of low labels at its first joint availability, with $B$ a sufficiently large fixed power of $k$. Insert the conjugate predictors in the two high positions, with private weight $\epsilon$. Choose $K_0$ and $\log(1/\epsilon)$ as sufficiently large fixed powers of $k$ after $D_{\rm pred}$ and the low bounds. Lemma \ref{lem:delay} then gives error smaller than $e^{-2k}$ after all original-label sums. Only a polynomial in the logarithm of any subsequently enlarged high catalog occurs, so these choices survive later fixed-polynomial growth of that logarithm by increasing $k_1$.

Explicitly, a low pair is available below its deeper birth cell, or below the segment root if inherited. Its coefficient product and path multiplicity are bounded by a fixed power of $D_{\rm old}$. For its birth roots $J$, Cauchy--Schwarz and the integrated path bound give
\[
 \sum_J\|\one_JF_a\|_2\|\one_JF_b\|_2\le D_{\rm old}^C|R|.
\]
Apply the predictor approximation on both sides of a replacement, and then the delayed replacement to every retained term. A delayed opponent is current plus a rolling square-energy sequence, as already proved.

Add tuple-specific accurate high fits to obtain pair control. Proceed chronologically in projection depth $h$, first admitting predetermined predictors. On its coarse cell $I$ of depth $s=h-K_0$, test the residual $f_e-(P^{\rm new}_{e,h}F_e)_{\phys}$ in $d_{e,I}$. Its $L^2(I)$ norm is bounded, since the fine-partition projection is a contraction on $I$. A failed test at a precision with logarithm a fixed power of $k$ gives a bounded atom with coarse moment $\rho'$, $\log(1/\rho')\le k^{O(1)}$. Insert it with one fresh private label on all depth-$h$ subcells of $I$. If $m_Q$ are its residual moments there, the joint energy gain is at least
\[
 c\sum_Q|Q||m_Q|^2\ge c|I|\left|\avint_I\text{residual}\,\overline\phi\right|^2
 \ge c(\rho')^2|I|.
\]
Cap the number of these requests along a coarse path at a large multiple of $(\rho')^{-2}$. The orthogonal energy of the whole chronological fitting sequence bounds the first-excess coarse roots by a small fraction of root length. At a cut all ancestor delayed evaluations have already been computed; extend that old estimator on the cut subtree by retaining predictors and refitting, while assigning descendant summands to a fresh attempt. This gives the finer partitions required in \eqref{eq:delay-telescope} without changing old summands.

Choose the first high precision to dominate earlier costs, then the second to dominate the first new catalog's supremum bound. On a fine cell $Q\subset I$ at depth $s+K_0$, $\|F_e\|_{2,Q}\le2^{K_0/2}\|F_e\|_{2,I}$. Private-coordinate recovery bounds the first new projection on the whole coarse cell by $2^{K_0}$ times a fixed power of its own count and conditioning; it does not involve the second high precision. Thus these are successive, noncircular fixed-polynomial choices. Apply the pair comparison \eqref{eq:stagger}, treating the early low position directly when necessary as above. Actual-input terms cancel, and decreasing $\theta$ gives the asserted $\eta$. These tuple-specific fits do not enter the global recurrence for $k_t$.

Contraction gives the local sizes. Orthogonal rolling increments give \eqref{eq:basic-side}, including on normalized subroots using the inherited extensions. The term with both $U$ factors has the full complementary pair $Y_i,Y_j$, so its absolute sum is at most $C\eta(1+K_0)^2|R|$ by Cauchy--Schwarz; for the finitely many initial levels omit $\eta$.

We give the active-scale argument to finish the proof. For three fixed bounded basic columns let $A_I$ be their transpose norm against a normalized $L^2$ input in the fourth slot. On a path through $y$, the squared transpose norm is an integral in $(h,t,t')/r$, with the triple at anchor $y+h$ and its conjugate copy. These six atoms form a chart in three variables. After recentering at this fixed $y$, all unscaled oscillation parameters are fixed along the path. If $A_I>\delta$, apply Lemmas \ref{lem:weyl} and \ref{lem:chartfourier} at Fourier frequency zero, to precision small compared with $\delta^2$. Effective quadratic entries and linear $t$ frequencies are $r^2a\cdot\chi_*$ and $ra\cdot\chi_*$ for a fixed parameter vector, with heights quasipolynomial in $D_{\rm base},\delta^{-1}$ and polynomial dimension.

Choose tiny/large thresholds after the coefficient norm. Outside polynomially many intermediate scales, Lemma \ref{lem:height} makes every bounded effective quadratic entry tiny and each retained $t$ frequency tiny or large. Remove the tiny quadratic. Large $t$ frequencies are negligible by integration by parts; replace tiny ones by constants. The remaining $t$ integral is exactly zero by the marginal in the fourth slot, with $h,t'$ fixed. All approximation errors are below $\delta^2$, a contradiction. Hence the number of scales with $A_I>\delta$ is $\poly(\log D_{\rm base},\log(1/\delta))$. Applying this count to $2^{-n-1}\delta<A_I\le2^{-n}\delta$ and summing shows
\[
 \sum_{I\text{ on path}:A_I\le\delta}A_I
 \le\delta\,\poly(\log D_{\rm base},\log(1/\delta)).
 \eqtag{eq:triple-tail}
\]
Only uniform derivatives and the exact marginal were used, so kernels may vary with $I$.

Declare a cell active when some available triple has norm greater than $D_{\rm base}^{-C'}$. There are at most $D_{\rm base}^3$ triples on a path, yielding the stated multiplicity. On inactive cells expand any three $B$ arguments; their coefficient norms, triple count and remaining local size cost fixed powers of $D_{\rm base}$. Taking $C'$ larger than these powers, \eqref{eq:triple-tail} and \eqref{eq:path} give error $e^{-k}|R|$. The finitely many initial levels have fixed budgets and so a depth-independent bound. For later levels pair-smallness gives $|H_I(Y)|\le\eta$ after decreasing $\theta$ to absorb constants. Expand $Y=B+U$ and use
\[
 \min\{\eta,\sum_j a_j\}\le\sum_j\min\{\eta,a_j\}\quad(a_j\ge0).
 \eqtag{eq:min-subadd}
\]
Take absolute errors and the two-$U$ term outside this minimum. This leaves \eqref{eq:active-min} separately for the remaining patterns, with the specified whole $B$ complements unchanged. All basic hypotheses and rolling estimates hold on each normalized restart root.
\end{proof}
\bp{base}{THT.Base.activeCellReduction}{adaptive-fit, stagger, predictors, delay, detection, height, cuts, path-packing}




\section{Curvature localization}\label{sec:localization}
Fix one pattern in \eqref{eq:active-min} on a normalized root $R$. Write $D=D_{\rm base}$. In the following estimates $0<\alpha<1$ is arbitrary but fixed. Thresholds for $k$ may depend on this fixed $\alpha$, but all polynomial degrees are independent of it and of $q$. Lemma \ref{lem:budget-rules} allows fixed powers of $k$ to be absorbed in a fixed power of $P$ above such a threshold. Put $P=\lceil\eta^{-\alpha}\rceil$. For sufficiently large $k$, $2\le P\le e^k$. Reserve the notation
\[
 G_0=\exp\bigl(C_q(1+\log(2k))^C\bigr)
 \eqtag{eq:small-overhead}
\]
for losses of this size, allowing fixed powers and increases of $C$. For this fixed $\alpha$ they are $P^{o(1)}$, and $P$ dominates every fixed power of $k$. Hence all later bounds polynomial in $P,k$ will be written $P^C$. Thresholds may depend on $q,\alpha$, but these degrees do not. An exponential $\exp(P^C)$ is not a small overhead.

\subsection{Tags and the ambient graph}
Put $c=(-1,3,-3,1)$. A vertex is an original basic column occurrence. In slot $i$ write its quadratic phase as
\[
 c_i\bigl(\lambda_vy^2+(\theta_vy)^TC_v(\theta_vy-n)\bigr).
\]
Its normalized tag is $(\theta_v,C_v,\lambda_v)$; linear modulation is excluded. The original vertex lists increase along a path and have at most $D$ members, ordered by admission. All new columns will have one of these tags, with the appropriate slot multiplier, a fixed-complexity amplitude, and an arbitrary linear modulation. Thus new columns do not add graph vertices.

At depth $s$, length $r$, use the graph on all hypothetical tags of the basic complexity. Equal tags have distance zero. Join $v,w$ when the difference data
\[
 V=r(\theta_v,\theta_w),\qquad C=\operatorname{diag}(C_v,-C_w),
 \qquad \Lambda=r^2(\lambda_v-\lambda_w)
\]
admit \eqref{eq:weyl-splitting} with heights and sizes at most $H_{\rm ed}$. Here concatenation denotes the vertical concatenation of the horizontal vectors. Distances $d_s$ may be infinite. They depend only on tags and absolute depth, not on availability or position, and decrease with depth: refinement halves $V$ and quarters $\Lambda_*$ while retaining the same splitting.

\begin{lemma}[Graph separation and short paths]\label{lem:graph}
Given single-column complexity $\exp(k^{O(1)})$ and $\log(1/\tau_2)\le k^{O(1)}$, one can choose $H_{\rm ed}=\exp(k^{O(1)})$ so that $d_s(v,w)>1$ implies both of the following. Same-slot pure columns have normalized inner product at most $\tau_2$ on every depth-$s$ interval. Distinct-slot pure columns have form norm at most $\tau_2$ against arbitrary normalized $L^2$ spectators.

If two tags are joined at one scale by a path of length at most $P^C$, their difference satisfies \eqref{eq:weyl-splitting} in its own endpoint coordinates with logarithmic height and size bound $P^C$.
\end{lemma}
\begin{proof}
For same-slot columns, smooth the interval indicator with $L^1$ error much smaller than $\tau_2$. A larger inner product gives a Fourier coefficient of the difference chart above a fixed multiple of $\tau_2$; Lemma \ref{lem:weyl} gives the edge splitting. The fixed multiplier $c_i$ only changes rational heights by a fixed power.

For distinct source slots $i,j$, use the two spectator coordinates $(z,w)$ after centering and scaling. The source coordinates have linear parts $r(\xi_lz+\zeta_lw)$, where $\xi_l=(b-l)/(b-a)$ and $\zeta_l=(l-a)/(b-a)$ for the spectator slots $a,b$. They are nonzero fixed rationals, and $c_i\xi_i\zeta_i=-c_j\xi_j\zeta_j=\kappa\ne0$, by $\sum_lc_ll^u=0$ for $u=0,1,2$. Clip spectators at an inverse-polynomial height in $\tau_2$. Their $L^1$ tails are bounded by the inverse height; \eqref{eq:absolute} with exponents $(1,2,\infty,\infty)$ controls the error. For a small return of the difference data, translate $z$ by $h$ and source lift indices by $\xi_i m_v,\xi_jm_w$, choosing a period that clears these rationals. The return calculation \eqref{eq:return} produces the modulation $\kappa\Omega(h,m)$ in $w$. All errors have size $|h|+\|d_h\|$ times a quasipolynomial; the modulation in $z$ is absorbed in its spectator.

A correlation above $\tau_2$ would, by the first $M$ multiples of a sufficiently small return, give correlations of comparable size of one fixed bounded kernel $K_1(z,w)$ against $f_\ell(z)g(w)e(\ell\kappa\Omega w)$. The $f_\ell$ have bounded $L^2$ norms. If $|\Omega|$ exceeds a fixed constant, the exponentials are separated and, multiplied by the clipped bounded $g$, are Bessel on the fixed box. Hence
\[
 \sum_\ell\left|\int K_1(z,w)f_\ell(z)g_\ell(w)\,dz\,dw\right|^2
 \lesssim\int\sum_\ell\left|\int K_1(z,w)g_\ell(w)\,dw\right|^2dz
 \lesssim\|K_1\|_2^2,
\]
with a controlled constant. Choose $M$ larger than this bound divided by $\tau_2^2$, then the return threshold so all $M$ errors are small. This contradiction bounds $\Omega$ on short returns. Lemma \ref{lem:return-splitting}, with the obtained fixed return bound and its reduced threshold, now gives the edge splitting. This proves both separation statements.

For a path $v_0,\ldots,v_L$, concatenate all its horizontal coordinates and impose the rational conditions that each consecutive pair belong to its edge fast space. The resulting subspace $U$ has projection height and a right-inverse norm with logarithms polynomial in $L$, dimensions and $\log H_{\rm ed}$. Its projection $V^F$ of the concatenated horizontal vector differs from that vector by a controlled amount, because each defining equation has bounded error. On an edge the difference between $V^F$ and the chosen fast component is a bounded vector $\Delta$ in the isotropic edge space; the effective correction changes by $\sym(\Delta^TC_{{\rm sk},e}V_e)=\sym(\Delta^TC_{{\rm sk},e}V_{s,e})$, which is bounded. Sum the effective edge quadratics with these common fast and slow components. All intermediate diagonal, skew-mixed and ordinary quadratic terms telescope. The endpoint image $U'$ of $U$ is rational and isotropic: summing the zero edge skew pairings cancels every intermediate coordinate. The endpoint restriction of $V^F$ lies in $U'$ and is a bounded distance from the endpoint vector. Replace it by that vector's orthogonal projection onto $U'$; the difference is bounded and lies in $U'$, so the same isotropy identity controls its effective correction. This proves the endpoint assertion without retaining the large intermediate dimension.
\end{proof}
\bp{graph}{THT.Localization.curvatureGraph}{weyl, return-splitting, chart-operations, rational-linear, fourier-foundation}


\subsection{Decreasing colors and historical masks}
The following explicit construction supplies the nested masks and their whole-cell separation. Set
\[
 d_0=\lceil1000\log(2D)\rceil,\qquad L=d_0^{30},\qquad H_{\max}=100P(L+1),
\]
\[
 b_2=100(H_{\max}+1),\qquad b_1=d_0^{10}b_2(H_{\max}+1).
\]
Give each original vertex, independently in each of $d_0$ primary rounds, a geometric delay $D_z$ with untruncated tail $\Pr(D_z\ge n)=2^{-n}$, truncated to $[0,d_0]$. At a present vertex $v$, take the top score
\[
 g=\max_{z\in V(I)}\{D_z-\lfloor d_s(z,v)/b_1\rfloor\}.
\]
A round is marked if it has a unique winner with gap at least five. Use the first marked round and its winner $z$. If its admission index is $i$ and the present list size is $j=i$, use the equality cohort with prefix through $i$. If $j>i$, take the smallest dyadic index block whose lower half contains $i$ and upper half contains $j$, and freeze $W$ to be the prefix through that lower half. This is its upper-half cohort.

Within this context define, at all absolute depths and hypothetical tags,
\[
 f_a(s')=\max_{w\in W\setminus\{z\}}(b_1D_w-d_{s'}(w,a)),
\]
with empty maximum $-\infty$. These profiles increase with depth. For $1\le l\le d_0$ put
\[
 T_l=(g-3)b_1+\frac{lb_1}{2d_0},\quad U=(g-2)b_1,\quad
 C_l(x)=\min(U,\max(T_l,x)),
\]
\[
 d^{(l)}(a,b)=\sup_{s'\in\Z}|C_l(f_a(s'))-C_l(f_b(s'))|,
 \qquad R_l=10l(d_0+1)b_2.
\]
These are fixed pseudometrics for the context, with $d^{(l+1)}\le d^{(l)}$. Choose the least $l<d_0$ for which the balls $B_{d^{(l)}}(v,R_l)$ and $B_{d^{(l+1)}}(v,R_{l+1})$, counted on present vertices, have the same floor logarithm to base $3/2$, called $n$. Such an $l$ exists because the balls are nested, nonempty, and have at most $D$ members, while $d_0$ exceeds the possible log range.

Use $d_0$ further independent high rounds of geometric delays, shared if desired among contexts, now with score $D_a-\lfloor d^{(l+1)}(a,v)/b_2\rfloor$. Take the first marked high round, its winner $a$ and score $g_2$. The color metadata are the primary round and $g$, the equality case or separating length level, $l,n$, and the high round and $g_2$. Use fixed placeholders on failures. Initialize a change counter at $d_0^8$ and decrement whenever these data or either winner change. Order counter and metadata lexicographically, counter first, and embed them in $\{1,\ldots,L\}$. A successful group additionally remembers $z,a,W$ and the exact cohort block, but not the changing list size or central vertex.

The color is nonincreasing and there are at most $d_0^8$ changes per vertex. Indeed each competitor's primary score only increases. At a fixed top score a marked winner has at most one episode: once a rival destroys its gap it cannot recover that gap at the same top value. Across rounds there are at most $5d_0^2$ episodes. Within each there are at most $d_0$ cohorts. In a frozen cohort the vector of ball floor-log sizes has at most $2d_0^2$ epochs, since each coordinate increases. With the selected $l,n$ fixed, high scores increase only by admissions; there are at most $5d_0^2$ high episodes. Thus at most $50d_0^7\le d_0^8$ changes occur. The seven metadata coordinates and counter require fewer than $256d_0^{15}<L$ states.

The delays can be fixed so that first marking failures have an arbitrarily small fixed total root length. For one untruncated score situation, increase the canonically chosen top winner's delay by five. The map is injective and probabilities decrease by $2^{-5}$, proving marking probability at least $2^{-5}$. Truncation changes it by at most $D2^{-d_0}$, so it remains at least $1/64$. On a path, capped distances $\lfloor d_s(z,v)/b_1\rfloor\wedge(d_0+5)$ change only $O(D^2(d_0+1))$ times; values beyond the cap do not affect gaps. There are $O(D^3(d_0+1))$ primary situations to test. Conditional on the primary arrays, there are $O(Dd_0^5)$ primary/cohort/ball epochs and at most $D$ present prefixes in each, hence $O(D^2d_0^5)$ fixed high situations. The independent high arrays fail all rounds at any such situation with probability at most $(1-1/64)^{d_0}$. The union bound tends to zero as $D$ grows. Integrating the pathwise bound over $R$ gives a realization with small failed-path measure; its first failure cells are disjoint and have at most that length. Restart there. In what follows all required markings exist.

For a successful group $G$ and integer pad $1\le h\le H_{\max}$, put a present vertex $b$ in $M_{G,h}$ if a historical central vertex $v_0$ of this color and group occurred at or above the current cell and
\[
 b_1D_z-d_s(z,b)\ge(g-1)b_1-h,\qquad
 d^{(l+1)}(b,v_0)\le h,\qquad d^{(l)}(b,v_0)\le10R_lh.
 \eqtag{eq:ball-mask}
\]
The mask increases inside its cohort; afterwards retain its eligible vertices by inheritance when a nested operator is needed. A group has an actual use only at a current central assignment. A \emph{usable endpoint} is a cell with nonempty mask and an actual-use descendant. Such a cell is still in the group's cohort.

\bp{coloring}{THT.Localization.monotoneColoring}{graph, finite geometric probabilities, finite dyadic cohorts, path-packing}

\begin{lemma}[Mask properties]\label{lem:masks}
At an actual use centered at $v$, all present vertices within $h$ graph edges of $v$ belong to $M_{G,h}$. Every present neighbor of a vertex in $M_{G,h}$ belongs to $M_{G,h+1}$ when $h<H_{\max}$. Each admitted mask vertex has already at admission a path to $z$ of length $P^C$. For distinct groups of one color, any vertices in their pad masks on overlapping usable endpoints are at distance greater than one at the finer endpoint.
\end{lemma}
\begin{proof}
At a marked primary center, every rival's true score is at most $(g-5)b_1$, while the winner's is greater than $(g-1)b_1$. If $d_s(b,v)\le h$, its winner score loses at most $h$. Its rival profile and that of $v$ are below $T_l$ at every earlier depth; at later depths their clipped profiles differ by at most $h$ by persistent distance. This proves all predicates with witness $v$.

Compare a historical witness $v_0$ and the current center $v$. Their common high winner gives $d^{(l+1)}(a,v_0),d^{(l+1)}(a,v)<(d_0+1)b_2$. The earlier low ball at $v_0$ is contained in the current high ball at $v$, since $R_l+2(d_0+1)b_2<R_{l+1}$. Both low balls contain at least $(3/2)^n$ vertices, while the high ball has fewer than $(3/2)^{n+1}$; they intersect. Hence $d^{(l)}(v_0,v)\le2R_l$. A pad-$h$ vertex has low distance to $v$ at most $A_h=10R_lh+2R_l$. The chosen constants ensure $A_h<T_{l+1}-T_l-10$, so at this use
\[
 f_b(s)\le T_l+A_h<T_{l+1}-10.
 \eqtag{eq:rival-bound}
\]
For a graph neighbor $b'$, the winner score loses at most one. Both profiles are below $T_{l+1}$ before this depth and differ by at most one afterwards, giving $d^{(l+1)}(b,b')\le1$. For the low metric, after the depth its distance to $v_0$ is at most $10R_lh+1$. Before it, its clipped profile is between $T_l$ and $T_l+A_h+1$, and that of $v_0$ between $T_l$ and $T_l+2R_l$. Their difference is at most $A_h+1\le10R_l(h+1)$. This proves neighbor closure. At admission the winner-score predicate also gives $d_s(z,b)\le(d_0+1)b_1+H_{\max}$, proving the short path.

For separation take overlapping endpoints $I,J$, $J$ finer, and a later actual-use certificate $u$ for the group $G'$ on $J$. Suppose the two mask vertices $b,b'$ are adjacent at $J$; they stay so at $u$. If primary winners $z,z'$ differ, the earlier winner lies in the later frozen prefix $W'$. At a common separating length, cohorts on an increasing admission path either have the same dyadic block or ordered disjoint blocks; the assertion follows in either case, and also for equality cohorts. The earlier winning score persists, so adjacency gives a rival score at $b'$ at least $(g-1)b_1-h-1$. This contradicts \eqref{eq:rival-bound}, since $T_{l+1}\le(g-5/2)b_1$ and $b_1\gg H_{\max}$.

If $z=z'$, its admission index and common separating length give the same block and metrics. The high winners must differ. Order the historical witnesses $v_0,v'_0$ on the common ancestor path. At the later one the earlier high winner is a rival. Its score dropped from greater than $(g_2-1)b_2$ to at most $(g_2-5)b_2$, implying $d^{(l+1)}(v_0,v'_0)>4b_2>2H_{\max}+1$. At $u$, however, \eqref{eq:rival-bound} and adjacency make both profiles below the high clipping range at all earlier depths and within one at all later depths. Hence $d^{(l+1)}(b,b')\le1$. The two mask predicates give $d^{(l+1)}(v_0,v'_0)\le h+1+h'\le2H_{\max}+1$, a contradiction. No later use of $G$ on the branch of $u$ was assumed, so the proof also covers certificates on different branches.
\end{proof}
\bp{masks}{THT.Localization.historicalMaskSeparation}{graph, coloring}


\begin{lemma}[Whole endpoint ranges]\label{lem:whole-ranges}
Fix a slot and color. Let $Z_G$ be finite sums of available pure main columns, with coefficients constant on entire usable endpoint cells whose masks contain their tags. Repeated labels and arbitrary endpoint times are allowed. If group and label counts on paths and inverse private weights are $\exp(k^{O(1)})$, choosing $\tau_2$ sufficiently small after these bounds gives
\[
 \left\|\sum_G Z_G\right\|_2^2\lesssim\sum_G\|Z_G\|_2^2.
 \eqtag{eq:whole-ranges}
\]
\end{lemma}
\begin{proof}
Combine occurrences of label $v$ in group $G$ into $a_{v,G}(y)=\sum_Ia_{v,G,I}\one_I(y)$. Its private coordinate is $\sqrt{\epsilon_v}a_{v,G}$, so $\sum_v|a_{v,G}|^2\le W_0^2\|Z_G(y)\|^2$ if $W_0$ bounds inverse square-root weights. For labels in different groups let $s_*$ be their last separating depth, capped at the last depth under consideration. Lemma \ref{lem:masks} says every nonempty intersection of supporting endpoint cells has depth at most $s_*$. Thus the combined coefficient product $a_{v,G}\overline{a_{w,G'}}$ is measurable on the $s_*$ grid; finer terms have no overlap. On each grid cell Lemma \ref{lem:graph} bounds the physical pairing by $\tau_2$. Private labels cannot coincide on an overlap. Therefore
\[
 |\langle Z_G,Z_{G'}\rangle|\le\tau_2\sum_{v,w}\int|a_{v,G}a_{w,G'}|.
\]
Pointwise Cauchy--Schwarz bounds all cross terms by $\tau_2$ times a fixed power of the path counts and $W_0$, times $\sum_G\|Z_G\|_2^2$. Choose $\tau_2$ below its reciprocal. Combining occurrences first retains cancellations and avoids any factor for repeated endpoint times.
\end{proof}
\bp{whole-ranges}{THT.Localization.usableRangeBessel}{masks, graph, coeff}


In applications to adjoints, a restriction to a smaller use cell is on the \emph{input} side of its endpoint projection. The adjoint projection still produces constant coefficients on the \emph{whole} endpoint cell. The downstream certificate makes this a usable range, so \eqref{eq:whole-ranges} applies; arbitrary restrictions of its output ranges are not asserted to satisfy the lemma.

\subsection{Finite same-tag catalogs}
For each original tag there is a fixed family of bounded templates in every slot with this property: a pure source column of that tag having form norm $t_0$ against a residual in another slot and two normalized $L^2$ spectators forces a residual moment at least $t_0/H_0$ against a same-tag template. Here $H_0=\exp(k^{O(1)})$ is independent of insertion rank and linear modulation, and the source may itself be a template.

To prove it, choose a fixed nonnegative Gevrey partition $\rho$ and an integer period $P_0$ covering both original and partition supports. The templates have amplitude $\rho(d)e(\ell^Td/P_0)$ with $\ell$ taken modulo $P_0$, in addition to arbitrary external linear modulation. Insert this partition in the other slots and extract the matching phases $c_l(\lambda_vy_l^2+(\theta_vy_l)^TC_vd_l)$. The lifts obey $d_l=(1-l)d_0+ld_1+k_l$ with bounded integer carries. The identities $\sum c_ll^u=0$, $u=0,1,2$, leave only the linear carry phase $\sum_lc_l(\theta_vy_l)^TC_vk_l$. Impose each carry by an exact smooth equality cutoff and Fourier-expand the coupled lift amplitudes. Reduce every index as $\ell=P_0q+r$; since $e(\ell^Td/P_0)=e(q^T\theta_vy)e(r^Td/P_0)$, only the residue enters the template. The demodulated source is a smooth periodization and is expanded in characters as well. The derivative argument of Lemma \ref{lem:chartfourier} bounds all absolute coefficient norms by $H_0$, with no iterative nonlinear complexity growth.

For one term, use the spectator coordinates $y_a,y_b$, separate the smooth kernel, and absorb source linear modulations into spectators. The target coordinate is $y_j=\alpha y_a+\beta y_b$ with nonzero fixed rational $\alpha,\beta$. Fourier-expand the zero-extended demodulated residual on a fixed larger interval. If $M$ bounds its coefficients, the form is bounded by $M\sum_m|A_mB_m|\le M(\sum|A_m|^2)^{1/2}(\sum|B_m|^2)^{1/2}\lesssim M$: these spectator coefficients have fixed nonzero frequency spacings and are Bessel. Bounded templates preserve their $L^2$ norms. Finite expansions followed by $L^2$ convergence justify the calculation. Summing all separated terms proves the claim, with a harmless enlargement of $H_0$ to avoid an unattained supremum.

In addition to ball masks use the cumulative masks
\[
 U_{m,h}(I)=\{v\in V(I):d_s(v,\{w:\operatorname{color}(w)\le m\})\le h\}.
\]
They increase in depth and have neighbor closure at pad $h+1$. Original labels have rank zero. A rank-$r$ projection uses all available columns of ranks at most $r$ in its vertex mask. All new labels have private weight one. A positive smeared filter averages $P$ such projections with pads $h+p$, $1\le p\le P$. They are nested both in $p$ and in depth, so their averages and complements are contractions. At a cumulative split $m$ use rank $L-m+1$ and base pad $10P(L-m+1)$. The four final ball insertions use ranks $L+j$, $1\le j\le4$, and base pads $20P(L+1)+10(j-1)P$. All needed neighbor enlargements lie below $H_{\max}$.

Close the catalogs under the following tests. For each target rank $r$, individual mask, and required history of length $O(\log L)$ on a fixed basic $B$ or $U$, apply the target residual projection and require its form norm against every lower-rank pure source whose tag belongs to that mask to be at most
\[
 \frac{\tau_3}{D\sqrt{n_{r-1}}}\,\|\text{base of history}\|_{2,I},
 \qquad \tau_3=D^{-C_\tau}.
 \eqtag{eq:catalog-test}
\]
The spectators are arbitrary normalized $L^2$ functions. Histories include every individual pad choice, so the tests also hold for averages. Reserve one additional helper rank $L+5$ for the same tests on any family of fixed whole vector-valued bases, with main sources only. The helper assertion is a conditional universal statement for such fixed bases; it is not instantiated during the main construction. No individual mode, height, or amplitude selection will be a helper base.

\bp{same-tag}{THT.Localization.sameTagTemplateTest}{chart-operations, fourier-foundation, absolute, gevrey}

\begin{lemma}[Catalog closure]\label{lem:closure}
These main tests terminate on every active cell with path counts $n_r$ satisfying $\log n_{L+5}\le k^{O(1)}$. The same bound holds for the reserved helper closure whenever its fixed bases number at most $\exp(k^{O(1)})$ per active cell and have private coordinates disjoint from fresh helper labels. Degrees are independent of $H_{\rm ed}$ and all later mode precision.
\end{lemma}
\begin{proof}
For a fixed original graph, masks and named histories are fixed. Their number is $\exp(k^{O(1)})$: names involve original vertices, cohorts, ranks, pads, and $O(\log L)$ factors, not the numeric graph distances or inserted source identities. On failure the same-tag test supplies a fresh column with moment at least \eqref{eq:catalog-test} divided by $H_0$. Its private direction is absent from the old state. Thus, for old and new target projections and the old history value $X$,
\[
 \|(P_{\rm new}-P_{\rm old})X\|_{2,I}
 \gtrsim\frac{\tau_3}{D\sqrt{n_{r-1}}H_0}\,\|\text{base}\|_{2,I}.
\]
A zero base cannot fail. Although $X$ changes with other insertions, every individual positive projection in its named history is nested in the insertion index. In a trial of $T$ insertions, expand complements and apply \eqref{eq:prefix} with target increments as orthogonal rows and right histories at the old index. Their squared actions sum to at most $(C\log(2T))^{O(\log L)}\|\text{base}\|_{2,I}^2$. Therefore rank-$r$ insertions at one cell number at most $n_{r-1}\exp(k^{C_1})(C\log(2T))^{O(\log L)}$. This counts histories and targets, not triggering source labels: any source charges the same target-increment sequence. Absorb the active path multiplicity and ranks into $\exp(k^{C_1})$, and choose $n_r/n_{r-1}=\exp(k^{C_2})$ with fixed $C_2$ sufficiently large. There are $L+5=k^{O(1)}$ ranks, so the resulting global trial cap $T=n_{L+5}$ has polynomial logarithm. Its binary-prefix factor is smaller than the chosen ratio, so no cap is attained. The argument accommodates simultaneous insertions at other ranks because all positive families stay nested. The helper proof is identical on its fixed bases.

All bounds use only template complexity and the number of named original-vertex contexts. Thus first fix this count, then choose $\tau_2$ for \eqref{eq:whole-ranges} and all mismatch errors, and finally choose $H_{\rm ed}$ by Lemma \ref{lem:graph}. The graph-dependent construction then obeys the precomputed bound, so this order is not circular.
\end{proof}
\bp{closure}{THT.Localization.catalogClosure}{same-tag, prefix, base, coloring, masks, coeff, budget-rules}


Write $D_m$ for a common main catalog, coefficient, inverse-weight and use-count budget, with $\log D_m\le k^{O(1)}$. Because only original columns have small private weights, expanding a rank-$r$ outer primal costs at most $D\sqrt{n_r}$ times its input norm. This explains the denominator in \eqref{eq:catalog-test}; a full coefficient sum will not be charged in retained main estimates.

\subsection{Binary localization and variation}
\begin{proposition}[Current-group decomposition]\label{prop:localize}
Up to total error $Ce^{-k}|R|$, each of $\poly(k)$ fixed search leaves contributes
\[
 \sum_{G\text{ current}}H_I(Z_{0,G,s},Z_{1,G,s},Z_{2,G,s},Z_{3,G,s}),
 \qquad Z_{l,G,s}=L_{l,G,s}\mathcal B_{l,s}z_l.
 \eqtag{eq:localized}
\]
Here $L_{l,G,s}$ is one positive smeared ball filter and $\mathcal B_{l,s}$ is a shared product of cumulative filters or complements of length $O(\log L)$. Slot and product choices depend only on the search leaf. A group is used only at an actual current central assignment. The cell minimum may be distributed over leaves.
\end{proposition}
\begin{proof}
Maintain a color interval $[\ell,h]$, an upper primal guard, and for $\ell>1$ a lower residual guard. Initially use any basic $B$ slot as upper guard and $[1,L]$. At $m=\lfloor(\ell+h)/2\rfloor$ choose a third slot, split by its cumulative filter and complement, and in the positive branch make it the upper guard with range $[\ell,m]$; in the residual branch make it the lower guard with range $[m+1,h]$. Release the former guard of that kind, but never postcompose an unreleased guard. A deterministic slot rule gives at most $L$ leaves and $O(\log L)$ factors per history.

At color $k'$ the upper guard is basic or has rank $L-k'+1$. Expand only this outer primal. Call a source old if it is within its upper pad of a vertex of color at most $k'-1$, taking pad zero in the basic case. When $k'>1$, the lower residual guard has rank $L-k'+2$ and pad larger by at least $10P$, so \eqref{eq:catalog-test} kills every old source. For $k'=1$ there are none. Every other source is near a current vertex of color $k'$, hence lies in the first ball mask of its group by Lemma \ref{lem:masks}. It cannot match two groups, and is separated from every other current group's ball sources.

In a third slot insert the sum of first-rank ball projections over current groups. The matching difference from identity is controlled by \eqref{eq:catalog-test}; other groups are controlled by pure pair separation in Lemma \ref{lem:graph}. Old sources are still controlled by the untouched lower guard. Keep this third slot as pure group source and postcompose the other three slots by the successive larger-rank, larger-pad ball filters. Each residual insertion is tested against that fixed lower-rank source in its mask. This gives \eqref{eq:localized}.

A target-rank test expands a lower-rank primal with coefficient norm at most $D\sqrt{n_{r-1}}$, canceled by \eqref{eq:catalog-test}. Local sizes and same-time contractions are bounded. The remaining counts of uses, groups, leaves and insertion steps are fixed powers of $D$, so $C_\tau$ makes their total error below $e^{-k}|R|$. Pair-mismatch terms may spend powers of the full main count; $\tau_2$ was chosen after that count to absorb them. Averages remain averages and do not pay their number of pad choices. Finally use \eqref{eq:min-subadd}.
\end{proof}
\bp{localize}{THT.Localization.binaryGroupExpansion}{closure, masks, graph, base, coeff}


For a basic $B$ slot, $\mathcal B_{l,s}B_l$ is a fixed signed sum of averaged products of $O(\log L)$ main projections acting on $F_l$, shared across groups. The compression loss is $G_0$. Before an attempt set ball outputs to zero; between actual uses retain the inherited projections on all needed endpoint cells.

\begin{lemma}[Certified variation]\label{lem:variation}
For a basic coarse output $Z_{G,s}$ in \eqref{eq:localized}, let $[r,t]$ be common disjoint or boundedly overlapping depth intervals. At each right endpoint restrict to a union $E_{G,r,t}$ of $t$-cells with downstream actual-use certificates. Then
\[
 \sum_{G,[r,t]}\|\one_{E_{G,r,t}}(Z_{G,t}-Z_{G,r})\|_2^2
 \le G_0\|\one_RF\|_2^2.
 \eqtag{eq:variation}
\]
Without an outer ball the same assertion holds for a shared history, omitting $G$. Bounded overlap multiplies the bound by its multiplicity.
\end{lemma}
\begin{proof}
Expand complements and couple individual pad choices across times; Jensen reduces to $Z_{G,s}=L_{G,s}B_s'F$ with individual ball projection and a shared product of main projections. Split
\[
 Z_{G,t}-Z_{G,r}=(L_{G,t}-L_{G,r})B_r'F+L_{G,t}(B_t'-B_r')F.
\]
For fixed $G$, the rows $\one_E(L_{G,t}-L_{G,r})$ are Bessel by orthogonal increments. Their adjoints lie in sums of whole ball ranges at the two endpoints; a right certificate certifies every needed earlier endpoint with nonempty mask. Lemma \ref{lem:whole-ranges} makes the rows jointly Bessel across groups. They are local to $r$-cells, so compression inserts $B_r'$ at cost $G_0$.

For the second term, sum the outer ball rows at each common right endpoint using \eqref{eq:whole-ranges}. Telescope $B_t'-B_r'$, discard left contractions, and use the changed middle factors as orthogonal rows on $[r,t]$. Compression inserts their right factors, evaluated at $r$. The $O(\log L)$ terms, signed expansions and averages are absorbed in $G_0$. This also proves the shared-history assertion. Color the depth intervals to handle bounded overlap. Intervals crossing the attempt start have bounded overlap and are controlled by the one-time projection estimate with the zero convention. All adjoint ranges have been kept on entire endpoint cells.
\end{proof}
\bp{variation}{THT.Localization.certifiedVariation}{localize, whole-ranges, compression, prefix}




\section{Fixed modes, the linear side, and hereditary representatives}\label{sec:modes}
This section proves the strong linear-side energy estimate before it is used. Fix a search leaf and one of its basic $B$ roles. Every group output is an outer averaged main ball projection followed by shared main histories. The main budget is $D_m$, with $\log D_m\le k^{O(1)}$. Let $d_{\max}\ge2$ bound the dimension of one tag or a pair of endpoint tags; thus $d_{\max}\le k^{O(1)}$, independently of $P$.

\subsection{One expansion at first match}
First prove the expansion and lift assertions for an arbitrary fixed branch offset, on the entire finite attempt tree, before selecting any depth blocks. All bounds in this construction are uniform in that offset. A \emph{primary} is an individual main column available to the maximal pad and rank of this role's outer ball. At its first match on $I_0$, choose one expansion and keep its modes and labels on every descendant. Separate first-match branches may have separate mode labels. The mode lists, including primaries whose coefficients are zero, are nested.

For group center $z$, choose an integer $Q_0$ in the main budget clearing $C_z$, and put $\vartheta=\theta_z/Q_0$. For a uniform offset $\rho\in[0,1)^{\dim\theta_z}$ set
\[
 d_\rho(y)=\rho+\{\theta_zy-\rho\},\qquad
 p_\rho(y)=\lambda_zy^2+(\theta_zy)^TC_zd_\rho(y).
\]
Choose a product periodic Gevrey cutoff $\zeta^\rho$, between zero and one, vanishing within $\delta_c$ of each branch cut and equal to one outside distance $2\delta_c$, where $\delta_c=\exp(-k^{C_c})$ with $C_c$ large after the main budgets. Translation invariance gives pointwise
\[
 \E_\rho\zeta^\rho=c_\zeta\ge\tfrac12,
 \qquad \E_\rho|1-(\zeta^\rho)^2|^2\le Cd_{\max}\delta_c.
 \eqtag{eq:offset-mean}
\]
For zero horizontal dimension take $\zeta^\rho=1$. Each role/group uses one offset for all its branches and times, chosen independently of main selections.

At first match the short-path part of Lemma \ref{lem:graph} and Lemma \ref{lem:chartfourier}, applied to the smoothly cut off difference chart, give a finite expansion
\[
 \zeta^\rho\phi_\Phi e(-c_lp_\rho)
 =\sum_m a_{\Phi,m}e(\lambda_my^2+\gamma_my)+O(\delta_c)
 \quad\text{on }I_0.
 \eqtag{eq:firstmatch}
\]
All modes of one primary have the same $\lambda_m$, since one splitting was used. With $\log D_1\le P^C$, their number, absolute coefficient norm and $|\lambda_m||I_0|^2$ are at most $D_1$. Moreover
\[
 \gamma_m-\gamma_{m'}=a\cdot(\theta_v,\vartheta),
 \qquad a\in\Q^{\dim(\theta_v,\vartheta)}\text{ of height at most }D_1.
 \eqtag{eq:mode-differences}
\]
Recentring adds only a common unrestricted linear phase; Lemma \ref{lem:chartfourier} gives these rational differences.

Introduce augmented modes
\[
 \psi_m(y)=\bigl(\zeta^\rho(y)e(c_lp_\rho(y))e(\lambda_my^2+\gamma_my),\sqrt{\epsilon_m}e_m\bigr),
 \eqtag{eq:modes}
\]
with private indices distinct across primaries and first-match branches. There is a single pointwise bounded linear map $\mathcal R_{l,G}$ carrying every main ball output to exactly its corresponding mode synthesis. To see this, put
\[
 q_\Phi=\sum_ma_{\Phi,m}(\psi_m)_{\phys}-(\zeta^\rho)^2\phi_\Phi,
 \qquad |q_\Phi|\le\delta_c.
\]
For a main augmented input $F$, define its output physical coordinate to be
\[
 (\zeta^\rho)^2F_{\phys}+\sum_{\Phi\text{ matched at }y}\frac{F_\Phi}{\sqrt{\epsilon_\Phi}}q_\Phi,
\]
and its mode coordinate $m$ belonging to $\Phi$ to be $\sqrt{\epsilon_m}a_{\Phi,m}F_\Phi/\sqrt{\epsilon_\Phi}$, zero before the match. The physical correction has pointwise norm at most $\delta_cD_m^C$. Choose the mode weights small enough that the squared norm of the private map is at most one; this requires only $\log(1/\epsilon_m)\le P^C$. Thus $\mathcal R_{l,G}$ has an absolute norm bound. A main output has constant coefficients $\beta_\Phi$ on its cell and private coordinates $\sqrt{\epsilon_\Phi}\beta_\Phi$; substitution gives exactly $\sum_{\Phi,m}\beta_\Phi a_{\Phi,m}\psi_m$. No projection was commuted with this map.

Choose an a priori budget $D_*\ge e^k$, $\log D_*\le P^C$, covering all mode counts, rational heights in \eqref{eq:mode-differences}, inverse private weights, $Q_0$, and all modes and actual main uses along a path across groups. It also bounds the absolute coefficient sums of $\mathcal R_{l,G}Z_{l,G,s}$ on every normalized endpoint cell, at any depth $s$. These bounds follow from private-coordinate recovery on that cell and the uniform first-match estimates; they do not depend on a block length or on a later freezing depth. Thus the numerical value of $D_*$ may be chosen uniformly in the branch offsets, before selecting their realizations.

\bp{first-match}{THT.Modes.firstMatchExpansionAndLift}{graph, chart-fourier, coeff, localize, gevrey}

There is a common progression multiplier needed later. Allow four different offsets. With $d_j=d_{\rho_j}(y_j)$, the bounded integers $k_j=d_j-(1-j)d_0-jd_1$ give
\[
 \prod_je(c_jp_{\rho_j}(y_j))
 =e\!\left(\sum_jc_j(\theta_zy_j)^TC_zk_j\right),
 \eqtag{eq:gauge-cancel}
\]
using the three identities for $c$. For fixed carries the right side is an integer character on $(\vartheta y_{l'},\vartheta t)$ for any anchor $l'$, because $Q_0$ clears $C_z$. Slightly enlarge the cutoffs inside each branch region, making them one on the required supports, and extend by zero at its cuts. The disjoint carry pieces produce a smooth periodic multiplier bounded by one. Its Gevrey and absolute Fourier bounds are uniform in offsets. Enlarge $D_*$ to cover these full Fourier norms, including $(\zeta^\rho)^2$, and only then truncate the Fourier series to tails at most $D_*^{-300}$. Let $B_*$ be a common cutoff, with $\log B_*\le P^C$.

Fix $\Delta=10$ and a sufficiently small absolute $\delta_*>0$. Set
\[
 C_*=(10d_{\max}D_*)^{100d_{\max}^4},\qquad
 R_*^{\rm ht}=\lceil\log_2C_*\rceil+2\Delta+10.
\]
Expressing a height-$D_*$ vector in a basis of such vectors in dimension at most $d_{\max}$ costs coefficients and common denominator at most $C_*$, by determinant bounds. Choose an integer
\[
 J_*\ge C\delta_*^{-2}d_{\max}R_*^{\rm ht}\log(2D_m).
\]
Take $A_0$ to be a sufficiently large fixed power of $C_*6^{R_*^{\rm ht}}B_*D_*$, and put $L_h=6^h$, $E_h=A_0^{h+1}$ for $0\le h\le J_*$. Finally set
\[
 \Xi=D_*^{1000},\qquad
 H>(d_{\max}C_*E_{J_*}B_*L_{J_*}Q_0\Xi)^{10},\qquad
 T>(d_{\max}H)^{Cd_{\max}^4}\Xi.
 \eqtag{eq:mode-parameters}
\]
All logarithms are polynomial in $P,k$. Choose $K$ last, a dyadic integer polynomial in $P,k$, at least the logarithmic packet width in Lemma \ref{lem:height} at height $H$ plus $2\log_2T$, and sufficiently large that $D_*^C2^{-2K}$ is smaller than every fixed inverse power requested below. All such requests have fixed exponents. Now partition the depths into blocks $[b,b+K)$ and set
\[
 u_b=b-4K,\qquad v_b=b-2K,\qquad t_b=b+K.
 \eqtag{eq:lags}
\]
Freeze $Z_{l,G,u_b}$, setting it to zero before the attempt. Retain a $v_b$-cell $J$ only if a descendant actual use of $G$ occurs in that block, and write $\widetilde Z_{l,G,b}=\mathcal R_{l,G}Z_{l,G,u_b}$ on the retained branches. The preceding uniform budget applies to these particular endpoint evaluations, and their coefficients are constant on their $u_b$-cells. A mode available at $u_b$ has first match no later than $u_b$, so on a $v_b$-cell
\[
 r_{v_b}\osc_J\omega_m\le CD_*2^{-4K},\qquad \omega_m(x)=2\lambda_mx+\gamma_m.
 \eqtag{eq:slope-variation}
\]
This applies also to earlier retained modes on later pre-cells.

\paragraph{Frozen-mode data interface.}
For each primary and its first-match cell, store the immutable finite list of modes, its one effective curvature, the coefficient vector, the bound \eqref{eq:mode-differences}, and its new private indices. The pointwise lift is the explicit coordinate map just defined; no commutation law with projections is part of the interface. Store $D_*,B_*,J_*,L_h,E_h,H,T,K$ in the displayed order, and the block depths \eqref{eq:lags} only after these values have been chosen. A frozen evaluation carries its full decision cell, its pre-cell, and a downstream-use certificate. The list at a descendant is required to contain the inherited list with the same numerical mode coefficients.
\bp{mode-schedule}{THT.Modes.chooseSchedule}{first-match, height, rational-linear, budget-rules}

\subsection{Two graphs and height stabilization}
On a retained $v_b$-cell $J$, use the common list frozen at $u_b$. At height $h$ the PRE and CURRENT edges are, respectively,
\[
 \sup_{x\in J}|\omega_m(x)-\omega_{m'}(x)+n\cdot\vartheta/L_h|\le\Xi/r_{v_b},
 \eqtag{eq:PRE}
\]
\[
 \sup_{x\in J}|\omega_m(x)-\omega_{m'}(x)+n\cdot\vartheta/L_h|\le\Xi/r_{t_b},
 \qquad n\in\Z^{\dim\theta_z},\quad |n|_\infty\le E_h.
 \eqtag{eq:CURRENT}
\]
Denote components by $E_h(m)$ and $F_h(m)$. PRE refines CURRENT. Both graphs increase with height (replace a witness by $6n$) and with successive blocks on a path (lists increase, cells shrink, thresholds grow). Write $\pi(S)$ for the primary labels in a mode set $S$. ``Minimal'' means height zero.

A primary with $\Theta=(\theta_v,\vartheta)$ is nonexceptional at the block when every rational $a$ of height at most $H$ satisfies
\[
 |a\cdot\Theta|\le T^{-1}/r_{v_b}\quad\text{or}\quad |a\cdot\Theta|\ge T/r_{t_b}.
 \eqtag{eq:gap}
\]
Each primary is exceptional in at most $Cd_{\max}$ blocks on a path. Indeed $x(a)=\log_2(|a\cdot\Theta|r_0)$ lies in at most $d_{\max}$ packets of width $W$. Exception means $b-2K-\log_2T<x(a)<b+K+\log_2T$. A packet affects starts $b$ in an interval of length at most $W+3K+2\log_2T\le4K$, hence at most five starts. Zero evaluations are always tiny.

At a nonexceptional block let $K_{\rm rat}$ be the rational span of height-$H$ tiny vectors. Every height-$H$ vector in this span is itself tiny. Express it in a basis of tiny vectors; its coefficient sum is at most $(d_{\max}H)^{Cd_{\max}^3}$. Its evaluation is therefore below the large threshold by \eqref{eq:mode-parameters}, and \eqref{eq:gap} forces the tiny alternative. Put $K_B=\{w:(0,w)\in K_{\rm rat}\}$.

\bp{tiny-space}{THT.Modes.exceptionalBlocksAndTinySpan}{mode-schedule, height, rational-linear}

\begin{lemma}[Stabilization]\label{lem:stabilize}
For a nonexceptional primary, all but $Cd_{\max}R_*^{\rm ht}$ heights have this property: two of its modes in one CURRENT component at height $h$ are directly adjacent in PRE at $h-\Delta$. Call those heights good for the primary.
\end{lemma}
\begin{proof}
Write $R'=R_*^{\rm ht}$ and discard $h<R'$. A simple CURRENT path has at most $D_*$ edges. Between two modes of the same primary their quadratic parts cancel, so \eqref{eq:mode-differences} gives an integer $m_h$, $|m_h|_\infty\le D_*E_h$, with
\[
 |a\cdot\Theta+m_h\cdot\vartheta/L_h|\le D_*\Xi/r_{t_b}.
\]
The coefficient vector has height at most $H$, so the gap and $D_*\Xi<T$ imply $a+(0,m_h/L_h)\in K_{\rm rat}$.

Let $V_i$ be the rational span of height-$D_*$ vectors $q'$ admitting $q'+(0,m/L_i)\in K_{\rm rat}$ with integer $|m|_\infty\le D_*E_i$. Let $U_i$ be the rational span of integer vectors in $K_B$ of size at most $C_*^3D_*E_i$. Both families increase. Discard heights where either rank changes between $h-R'$ and $h$; there are at most $2d_{\max}R'$ such heights. At a remaining height, express $a=\sum_jb_ja_j$ in a defining basis of $V_{h-R'}$, with witnesses $m_j$. Coefficients and their common denominator are bounded by $C_*$. Subtracting the relations gives
\[
 m_h-6^{R'}\sum_jb_jm_j\in K_B.
\]
After clearing the denominator this is an integer vector of size at most $C_*^3D_*E_h$. It therefore lies in $U_h=U_{h-R'}=:U$, and so does the uncleared difference.

By Lemma \ref{lem:integer-basis}, the lattice $U\cap\Z^{\dim\theta_z}$ is saturated and its integer quotient is free. In an integer basis of that quotient,
\[
 [m_h]=6^{R'}\sum_jb_j[m_j].
\]
The coordinates on the right are integral. Since the common denominator is at most $C_*$ and $R'>\log_2C_*+2\Delta$, their valuations at both $2$ and $3$ are at least $\Delta$. Thus $[m_h]$ is divisible by $6^\Delta$. Choose an integer lift $m'$ of $[m_h]/6^\Delta$. Then $m'-6^{R'-\Delta}\sum_jb_jm_j\in U$. Select linearly independent integer generators of $U_{h-R'}$ of size at most $C_*^3D_*E_{h-R'}$ and round the coefficients of this difference, subtracting integer multiples of those generators from $m'$. It remains a lift and now obeys
\[
 |m'|_\infty\le6^{R'-\Delta}\sum_j|b_j|\,|m_j|_\infty
 +d_{\max}C_*^3D_*E_{h-R'}\le E_{h-\Delta},
\]
by the fixed-power choice of $A_0$. No lattice-basis or covolume bound was used. Since $m_h-6^\Delta m'\in U\subseteq K_B$, the height-$H$ vector $a+(0,m'/L_{h-\Delta})$ belongs to $K_{\rm rat}$ and is tiny. This gives a direct PRE edge at height $h-\Delta$.
\end{proof}
\bp{stabilize}{THT.Modes.heightStabilization}{mode-schedule, tiny-space, rational-linear, integer-basis}


For $h\ge\Delta$, call a mode good when
\[
 |\pi(E_{h-\Delta}(m))|\ge\tfrac9{10}|\pi(F_h(m))|
 \eqtag{eq:good}
\]
and at least nine tenths of $\pi(F_h(m))$ are nonexceptional and good for the primary at $h$. Heights below $\Delta$ are bad. A mode is core if at most $\delta_*(J_*+1)$ heights are bad; otherwise it is a linear-side mode.

Two good modes in one CURRENT component have the same PRE predecessor at $h-\Delta$: their predecessor primary sets each contain at least nine tenths of $\pi(F_h)$, so their intersection contains at least four fifths. It includes a nonexceptional height-good primary, whose two witnessing modes are joined by Lemma \ref{lem:stabilize}. Modes in the same minimal PRE component have identical predicates, so the core/side decision is constant there. Finally two core modes in one minimal CURRENT component have a common good height below $4\delta_*J_*$: their combined bad count is at most $2\delta_*(J_*+1)$. Their PRE predecessors agree there and, by nesting, at every higher height. Hence on
\[
 \mathcal H_{\rm up}=\{h:\lceil4\delta_*J_*\rceil\le h\le J_*\}
 \eqtag{eq:upper-heights}
\]
their good predicates are identical. Each core mode has bad proportion at most $2\delta_*$ in this range.

\subsection{Bessel rows and the linear-side energy}
\bp{coherence}{THT.Modes.goodPredicateCoherence}{stabilize, mode-schedule, finite intersections}

\begin{lemma}[Separated mode rows]\label{lem:mode-rows}
For a fixed role and group, let $P_e$ project on a nonempty mode set $M_e$ on a dyadic cell $J_e$. Suppose these sets are disjoint along every path, with total cardinality at most $D_*$, and every cross pair for overlapping cells (or distinct sets on the same cell) has normalized inner product at most $D_*^{-200}$ on the smaller cell. Then $\sum_e\|P_eF\|_2^2\lesssim\|F\|_2^2$.

The cross-pair condition holds on pre-cells when the sets have no shared or adjacent modes in the later PRE graph; it holds on actual-use cells with later CURRENT nonadjacency.
\end{lemma}
\begin{proof}
The common gauge in \eqref{eq:modes} cancels. Expand $(\zeta^\rho)^2$ on the $\vartheta$ torus with coefficient sum at most $D_*$ and tail $D_*^{-300}$. All retained shifts are allowed: $E_0>B_*$ and $E_h\ge B_*L_h$. Nonadjacency gives a point with phase derivative above the appropriate threshold. By \eqref{eq:slope-variation} its absolute value is at least half that threshold everywhere on the pre-cell. This affine derivative has constant sign. The first-derivative integration estimate gives normalized integral at most $C/\Xi$, both on a pre-cell and on an actual cell $I$ because $|I|\ge r_{t_b}$. The Fourier sum and tail give $CD_*/\Xi+D_*^{-300}\le D_*^{-200}$. Private coordinates of different labels are orthogonal.

By duality it suffices to bound synthesis of $f_e=\one_{J_e}\sum_{m\in M_e}\beta_{e,m}\psi_m$. Put $b_e=(\sum_m|\beta_{e,m}|^2)^{1/2}$. Private coordinates give $|J_e|b_e^2\le D_*^2\|f_e\|_2^2$. A cross pairing costs at most $D_*^{-199}|J_e\cap J_f|b_eb_f$. At a path point there are at most $D_*$ participating sets. Summing all cross terms by Cauchy--Schwarz gives at most $D_*^{-198}\sum_e|J_e|b_e^2\le D_*^{-196}\sum_e\|f_e\|_2^2$. Adding the diagonal proves the result.
\end{proof}
\bp{mode-rows}{THT.Modes.separatedProjectionRows}{mode-schedule, coeff, fourier-foundation}


\begin{lemma}[Composition using only the main budget]\label{lem:lift-composition}
Suppose, for each group, mode projection rows $P_e$ on $J_e$ have squared Bessel bound $G_0$. Let $J_e\subseteq u_e$, where $u_e$ is a usable main endpoint with a downstream use certificate, and draw the row's modes from the list available at $u_e$. Then the outputs
\[
 P_e\mathcal R_{l,G}L_{l,G,u_e}\mathcal B_{l,u_e}B_{l,u_e}
\]
have joint squared sum at most $G_0|R|$ over all groups and rows.
\end{lemma}
\begin{proof}
For one group $P_e\mathcal R_{l,G}$ is Bessel because the lift is one fixed bounded operator. The rows are local to $J_e$ and hence to $u_e$. Apply Lemma \ref{lem:compression} to the outer main ball projection at $u_e$, then restore its average by the row-space triangle inequality. The cost is a fixed power of $\log(2D_m)$. If $A_G$ is this row operator, every term of $A_G^*$ has outer factor $L_{l,G,u_e}$, so its range is a sum of whole usable main endpoint ranges. Lemma \ref{lem:whole-ranges} implies
\[
 \left\|\sum_GA_G^*w_G\right\|_2^2\lesssim\sum_G\|A_G^*w_G\|_2^2
 \le G_0\sum_G\|w_G\|_2^2.
\]
Thus the rows are jointly Bessel before the shared histories. Apply compression again, with decision cells $u_e$, to insert those histories and the constituent basic projections acting on $F_l$. Expand complements and the fixed signed basic sum, and retain averages as averages. Product lengths are $O(\log L)$, so the loss is $G_0$. Finally $\|F_l\|_2^2\le|R|$. Only main projections have been compressed; the larger mode budget has not entered this loss.
\end{proof}
\bp{lift-composition}{THT.Modes.composeLiftWithMainHistories}{first-match, compression, whole-ranges}


\begin{proposition}[Linear-side energy]\label{prop:linear-side}
Let $V_{l,G,b}$ be the exact subtotal of linear-side modes in $\widetilde Z_{l,G,b}$, including private coordinates, on retained pre-cells, and zero on branches without a needed output or certificate. Then
\[
 \sum_{G,b}\|V_{l,G,b}\|_2^2\le G_0|R|.
 \eqtag{eq:linear-energy}
\]
This is uniform in offsets and on fresh normalized attempt roots.
\end{proposition}
\begin{proof}
An event is $(b,J,E)$ with $J$ a retained pre-cell and $E$ a minimal PRE component of side modes; it is entirely side by the predicate coherence. Split block indices modulo three. Link an event to a later one on a descendant path when some modes agree or are adjacent in the later minimal PRE graph. Since starts in one class differ by at least $3K$, $v_{\rm later}\ge t_{\rm earlier}$. Thus earlier CURRENT edges persist in later PRE, at every height. Along a linked chain,
\[
 \pi(E_h^{(i)})\subseteq\pi(F_h^{(i)})\subseteq\pi(E_h^{(i+1)}).
\]
We bound the length $N$ of such a chain directly. Write $P_{i,h}=\pi(E_h^{(i)})$ and $C_{i,h}=\pi(F_h^{(i)})$. Their sizes are at most $D_m$. A failed predecessor-density test contributes at most
\[
 10\log\frac{|C_{i,h}|}{|P_{i,h-\Delta}|}
 =10\log\frac{|P_{i,h}|}{|P_{i,h-\Delta}|}
  +10\log\frac{|C_{i,h}|}{|P_{i,h}|}.
\]
The first term sums in $h$ to at most $10\Delta\log D_m$ per event; the second sums in $i$ to at most $10\log D_m$ per height. For the other failures use this elementary fact: if nonempty sets $S_j$ increase, $|S_j|\le M$, and each label belongs to at most $A$ subsets $X_j\subseteq S_j$, then $\sum_j|X_j|/|S_j|\le A(1+\log M)$. Enumerate labels by first entry; the label of rank $r$ contributes at most $A/r$, which proves the fact.

At one event, a nonexceptional primary is height-bad at most $Cd_{\max}R_*^{\rm ht}$ times by stabilization. At a fixed height, a primary is exceptional in at most $Cd_{\max}$ blocks. Apply the preceding fact in these two directions; a bad-fraction indicator costs at most ten times the fraction. Including the $N\Delta$ initially bad incidences gives
\[
 \delta_*N(J_*+1)
 <C N d_{\max}R_*^{\rm ht}\log(2D_m)
  +C(J_*+1)d_{\max}\log(2D_m).
\]
The choice of $J_*$ absorbs the first term into half the left side. Therefore
\[
 N\le C_{\delta_*}d_{\max}\log(2D_m).
 \eqtag{eq:bad-density}
\]
In particular the large height parameter does not occur in the chain-length bound.

Give events their longest-chain level. Events of one level on overlapping cells have neither shared nor adjacent modes in later PRE; distinct components on the same cell have this separation as well. Their label sets are disjoint along paths with total count at most $D_*$. Lemma \ref{lem:mode-rows} gives Bessel rows for each level. Summing the three classes and the bounded number of levels costs $G_0$. Their decision cells are the containing $u_b$-cells; Lemma \ref{lem:lift-composition} yields
\[
 \sum_{G,b,J,E}\|P_{E,J}\widetilde Z_{l,G,b}\|_2^2\le G_0|R|.
\]
For the exact coefficient subtotal $T_{E,J}$, all complementary modes are PRE-nonadjacent. Pair them against a unit vector of the augmented range. Its coefficient norm is at most $D_*^{3/2}$, the complementary coefficient norm is at most $D_*$, and cross pairings are at most $D_*^{-200}$. Hence $\|P_{E,J}\widetilde Z_{l,G,b}-T_{E,J}\|_{2,J}\le D_*^{-190}$. The number of events along a path in one level is at most $D_*$, and all groups are included in the path budget. These squared additive errors are negligible. On each $J$ the different minimal PRE subtotals are almost orthogonal by the same private-coordinate synthesis calculation. Their sum is $V_{l,G,b}$, proving \eqref{eq:linear-energy}. 
\end{proof}
\bp{linear-side}{THT.Modes.linearSideSquareEnergy}{stabilize, coherence, tiny-space, mode-rows, lift-composition}


\subsection{Side states and fixed offsets}
Let $A_{l,G,b}$ be the full core subtotal. In a basic $B$ role define
\[
 S_{l,G,s}=Z_{l,G,s}-Z_{l,G,u_b}+V_{l,G,b}\quad(b\le s<b+K).
 \eqtag{eq:sides}
\]
In a basic $U$ role take the whole $Z_{l,G,s}$ as a mandatory side, with no core. In physical coordinates,
\[
 Z_{l,G,s}=A_{l,G,b}+S_{l,G,s}+E_{l,G,b},
 \qquad E_{l,G,b}=(Z_{l,G,u_b})_{\phys}-(\widetilde Z_{l,G,b})_{\phys}.
 \eqtag{eq:core-side-error}
\]
All expansion and energy assertions so far are uniform statements for an arbitrary fixed offset. We now choose the offsets without requiring any measurable selection of expansions. On a normalized $u_b$-cell the already fixed main output $Z_{l,G,u_b}=\sum_\Phi\beta_\Phi\Phi$ has coefficient sum and physical supremum at most $D_m^C$, with no pre-lag loss. Every admissible choice of the first-match expansions satisfies
\[
 E_{l,G,b}=(1-(\zeta^\rho)^2)(Z_{l,G,u_b})_{\phys}-\sum_\Phi\beta_\Phi q_\Phi,
 \qquad |q_\Phi|\le\delta_c.
\]
Consequently its squared absolute value is bounded by the explicit majorant
\[
 B_{l,G,b}(\rho,y)=2|1-(\zeta^\rho(y))^2|^2|(Z_{l,G,u_b})_{\phys}(y)|^2
 +2\delta_c^2D_m^{C'}.
\]
The exponent $C'$ covers the squared main coefficient bound. This majorant is jointly measurable in $(\rho,y)$ and depends only on the main selections, the uniform budgets, and the explicitly defined cutoff. Equation \eqref{eq:offset-mean} and the main path multiplicity give
\[
 \E_\rho\sum_{l,G,I\text{ used}}\int_I B_{l,G,b(I)}(\rho,y)\,dy
 \le D_m^{C''}d_{\max}\delta_c|R|.
\]
All the actual main uses in this sum were fixed before the offsets. Select offsets for which the displayed nonnegative sum is at most twice its expectation; such a choice follows from Markov's inequality. Then instantiate the preceding first-match, lift, and linear-side existence results at these fixed offsets, using classical choice only after the measurable selection. Their numerical parameter bounds were uniform, so the same $D_*,H,T,K$ apply. The pointwise domination now proves
\[
 \sum_{l,G,I\text{ used}}|I|\|E_{l,G,b(I)}\|_{2,I}^2
 \le 2D_m^{C''}d_{\max}\delta_c|R|.
 \eqtag{eq:physical-error}
\]
Taking $C_c$ sufficiently large makes this smaller than any required fixed inverse main-budget power times $|R|$. Later decisions and cuts only remove nonnegative terms. No independence from the selected offsets, or measurability of a chart-selection map, is needed. Proposition \ref{prop:linear-side} applies to the selected objects because it was proved uniformly.
\bp{offset-selection}{THT.Modes.selectOffsetsByMajorant}{first-match, mode-schedule, linear-side, finite probability, Tonelli}

Rolling sides have integrated square budget $KG_0|R|$ by \eqref{eq:variation}, since the intervals $[u_b,s]$ have overlap $O(K)$. A fixed $V_{l,G,b}$ is counted at most $K$ times, so \eqref{eq:linear-energy} gives the same bound. A basic $U$ side is bounded by current group Bessel and \eqref{eq:basic-side}. Lemma \ref{lem:cuts}, applied to these integrated budgets, gives simultaneously
\[
 \sum_{s,G}\|S_{l,G,s}\|_{2,I_s(y)}^2\le P^C
 \eqtag{eq:side-cap}
\]
on retained paths, at an arbitrarily small fixed root-length cost. The lift is pointwise bounded and $Z_{u_b}=Z_s-(Z_s-Z_{u_b})$, so the full core and error sizes on retained cells are also at most $P^C$. Thus \eqref{eq:physical-error}, Cauchy--Schwarz over original uses, and \eqref{eq:absolute} make every term containing a physical error $O(e^{-k}|R|)$, even when the other factors are core or side sums.

\bp{side-caps}{THT.Modes.retainedSideCaps}{offset-selection, linear-side, variation, base, cuts}
\paragraph{Helper instantiation.}
At this point the whole states $S_{l,G,s}$ are fixed. For one active cell there are only the already bounded main groups, roles and prescribed histories; the helper bases are these entire states, not their individual mode subtotals or amplitude and height selections. Apply the universal helper part of Lemma \ref{lem:closure} to this finite family. Its count depends on the number of these names and the relative residual tolerance, not on a representation of a base by modes. Helper ranks exceed every main rank, so their insertion leaves all previously fixed main projections, offsets and side states unchanged. This provides exactly the residual tests needed for transferring filters from a whole current complement to a side.
\bp{helpers}{THT.Localization.instantiateHelperCatalog}{closure, side-caps, localize}

\subsection{Representatives, component sizes, and first tops}
At height $h\in\mathcal H_{\rm up}$ call a CURRENT component successful if it contains a core mode good at $h$. Its good modes have one PRE predecessor $E_{h-\Delta}$. Its type is
\[
 \bigl(\lfloor\log_2|\pi(E_{h-\Delta})|\rfloor,
       \lfloor\log_2|E_{h-\Delta}|\rfloor\bigr).
\]
Fix height, role, group and type. Representatives can be chosen in these predecessors and inherited under predecessor inclusion along paths. To prove this, link an earlier successful component to a later one when some modes are equal or CURRENT-adjacent at the later height-$h$ graph. Persistence gives $F_{\rm early}\subseteq F_{\rm late}$. With $F=\pi(F_{\rm late})$, goodness and equality of the first type coordinate imply
\[
 |\pi(E_{\rm late})|\ge\tfrac9{10}|F|,
 \qquad |\pi(E_{\rm early})|>\tfrac12|\pi(E_{\rm late})|
 \ge\tfrac9{20}|F|.
\]
Their intersection therefore contains more than $7|F|/20$ primaries, hence a later nonexceptional height-good one. Stabilization joins its two witnessing modes in later PRE, so $E_{\rm early}\subseteq E_{\rm late}$. Process pre-cells chronologically: inherit an earlier representative when its predecessor is contained, otherwise pick any mode by a fixed label order. Two earlier predecessors contained in the current one each have more than half its modes by the second type coordinate. They intersect and, lying on one path, are linked and nested by the preceding argument. Their representatives agree by induction. This proves consistency, including all actual descendants of one pre-cell.

Consequently different representatives on comparable actual-use cells have no shared or adjacent cross modes in the later CURRENT graph. At any $x_*\in I\cap I'$ their representative frequencies obey
\[
 |\omega_a(x_*)-\omega_{a'}(x_*)+n\cdot\vartheta/L_h|
 >\frac\Xi2\max(|I|^{-1},|I'|^{-1})
 \quad(|n|_\infty\le E_h).
 \eqtag{eq:representative-separation}
\]
Nonadjacency gives a derivative exceeding $\Xi/r_{t_{b'}}$ somewhere on the later pre-cell; \eqref{eq:slope-variation} makes it exceed half that everywhere, and the later actual cell has length at least $r_{t_{b'}}$.

\bp{representatives}{THT.Modes.hereditaryRepresentatives}{coherence, stabilize, mode-schedule, finite dyadic size classes}

At an actual cell, write $A_D$ for the full core subtotal in a minimal CURRENT component $D$. For a successful $F=F_h$ define
\[
 \sigma_F^2=\sum_{D\subseteq F}
 \left(\|A_D\|_{2,I}+D_*^{-30}\sum_{m\text{ core in }D}|\alpha_m|\right)^2.
 \eqtag{eq:component-size}
\]
Let $P_{F,I}$ project onto all modes of $F$, and let $V_F$ be its exact linear-side subtotal. The single-cell version of the synthesis calculation in Lemma \ref{lem:mode-rows} gives two-sided almost orthogonality for arbitrary subtotals in distinct minimal CURRENT components. Thus, for the full subtotal $T_F$,
\[
 \sigma_F^2\lesssim\|P_{F,I}\widetilde Z_{l,G,b}\|_{2,I}^2+\|V_F\|_{2,I}^2+D_*^{-50},
 \qquad \sum_F\|V_F\|_{2,I}^2\lesssim\|V_{l,G,b}\|_{2,I}^2.
 \eqtag{eq:size-comparison}
\]
Indeed $\sum_{D\subseteq F}\|A_D\|_2^2\lesssim\|T_F-V_F\|_2^2$, the coefficient floors cost at most $CD_*^{-60}(\sum|\alpha_m|)^2\le CD_*^{-58}$, and pairing complementary modes against a unit range vector gives $\|P_{F,I}\widetilde Z-T_F\|_{2,I}\le D_*^{-190}$. The second assertion follows by decomposing $V$ into all height-$h$ CURRENT components and discarding unsuccessful ones. 

\bp{component-sizes}{THT.Modes.componentSizeComparison}{mode-rows, first-match, linear-side}

\begin{proposition}[Pathwise label count]\label{prop:tops}
After further first-crossing cuts of arbitrarily small fixed root fraction, for every retained path, core role $j$, height, type and dyadic bin $\sigma_j\le\sigma_F<2\sigma_j$, the sets $A_j^G$ of distinct used representatives satisfy
\[
 \sum_G|A_j^G|\le\min(D_*,P^C\sigma_j^{-2}).
 \eqtag{eq:top-count}
\]
All component sizes and full local side sizes on retained uses are at most $P^C$, and \eqref{eq:side-cap} holds.
\end{proposition}
\begin{proof}
For a fixed height, type, and dyadic threshold $D_*^{-1}\le\sigma\le D_*^2$, take the first actual-use cells on each branch where a given representative has $\sigma_F\ge\sigma$. Tops for that representative do not overlap. Different representatives on overlapping tops have component-mode nonadjacency by \eqref{eq:representative-separation} and its stronger cross-mode conclusion. Their sets are disjoint on a path, with at most $D_*$ total modes. Lemma \ref{lem:mode-rows} applies to the rows $P_{F,I}$. Their decision cells are the containing $u_b$-cells, with the tops themselves as certificates. Lemma \ref{lem:lift-composition} gives
\[
 \sum_{\rm tops}\|P_{F,I}\widetilde Z_{l,G,b}\|_2^2\le G_0|R|.
\]
For the side part of \eqref{eq:size-comparison}, sum first at a fixed depth of a block: cells are disjoint and component subtotals sum in square to $C\|V_{l,G,b}\|_2^2$. Summing the at most $K$ depths and using \eqref{eq:linear-energy} gives $KG_0|R|$. Therefore
\[
 \sigma^2\sum_{\rm tops}|I|\le C(1+K)G_0|R|+CD_*^{-50}\sum_{\rm tops}|I|.
\]
Absorb the last term because $\sigma^2\ge D_*^{-2}$. There are polynomially many choices in $P,k$ of height, type, and threshold. Markov's inequality and first crossings of their partial top counts give simultaneously $\sum_{\rm tops}\one_I(y)\le P^C\sigma^{-2}$ outside a small root fraction, after increasing the fixed power for their number. Each representative used in the bin has a retained first $\sigma$-top above that use; otherwise the use would have been removed. This proves the second bound in \eqref{eq:top-count}, while the full mode budget proves the first. For $\sigma<D_*^{-1}$ the full budget already implies the minimum; above $D_*^2$ the bins are empty by coefficient and atom bounds, enlarging $D_*$ by a fixed power if necessary. Finally the lift, the identity for $Z_{u_b}$, and the componentwise side caps bound $\widetilde Z,V$ by $P^C$ on uses. Projection contraction and \eqref{eq:size-comparison} give the component bound. All first-top estimates apply anew on normalized restart roots; no restricted global energy bound is being assumed.
\end{proof}
\bp{tops}{THT.Modes.firstTopCount}{representatives, component-sizes, mode-rows, lift-composition, linear-side, side-caps, cuts}




\section{The two-side estimate}\label{sec:sides}
The side states are now fixed by \eqref{eq:sides}; their linear parts have the root-uniform energy bound \eqref{eq:linear-energy}. Construct the helper rank reserved in Lemma \ref{lem:closure}, on these whole states and the specified histories. Their number on a path is within the main budget, and their extra private directions are distinct from the fresh helper directions. Thus the helper count still has logarithm $k^{O(1)}$, independently of the mode count. No individual Fourier mode, height, or amplitude choice is a helper base.

\subsection*{Transfer and its paired defects}
Suppose positions $a,b$ are sides $S_a,S_b$ and the other positions are whole current outputs. Expand complements in the latter two histories as identity minus a positive average. There are only $2^{O(\log L)}$ products. In each product remove the larger rank among the two current outermost positive filters, regarding an unfiltered basic argument as rank zero. The ranks are distinct: a cumulative rank was inserted in just one slot at its binary split, and the four ball ranks are distinct. Therefore the other complementary argument always has a strictly lower-rank outer primal source.

For a filter outgoing from $c$, write $Q_{c,p}$ for its individual pad projections and $Q_{a,p}$ for the helper projections on the identical vertex masks and pads. Set, for $l=a,c$,
\[
 Q_{l,0}=0,\qquad Q_{l,P+1}=1,\qquad A_p^l=Q_{l,p}-Q_{l,p-1},\qquad \overline Q_l=P^{-1}\sum_{p=1}^P Q_{l,p}.
\]
These bands are orthogonal and sum to the identity. On band $p$, $\overline Q_l$ has weight $(P-p+1)/P$, including weight zero at $P+1$. If $x_c$ is the inner complementary history, $x_a$ is the helper history already applied to $S_a$, $x_b=S_b$, and $x_d$ is the other complementary history, the exact difference is
\[
 H_I(x_a,x_b,\overline Q_c x_c,x_d)-H_I(\overline Q_a x_a,x_b,x_c,x_d)
 =\sum_{p,q=1}^{P+1}\frac{p-q}{P}H_I(A_p^a x_a,x_b,A_q^c x_c,x_d).
 \eqtag{eq:transfer}
\]
Only $|p-q|\le2$ must be retained. Here and below formulas written in role order are placed in their original progression slots.

Indeed, if $p>q+2$, expand the lower $c$ band into its main pure sources. Their tags lie in both masks writing the higher helper band as a difference of residuals, and their ranks are lower than the helper rank. The helper test \eqref{eq:catalog-test} applies with arbitrary $L^2$ spectators. If $q>p+2$, expand instead the lower-rank primal source in position $d$, using its basic columns when it is unfiltered. Call its tag near when it is within one edge of a vertex in the outgoing pad-$p$ mask. Neighbor closure puts a near source into pad $p+1$ and hence into both masks of the higher $c$ residual band. The main test applies, keeping the entire lower helper band as an $L^2$ spectator. For a far source, expand the helper primal and use the different-tag pair estimate. Cumulative masks have neighbor closure by the triangle inequality; ball masks have it at actual uses by Lemma \ref{lem:masks}. The two-pad gap includes the needed margin.

The main source coefficient factor is canceled by the denominator in \eqref{eq:catalog-test}. Only pair mismatches expand both primals, and their fixed powers of the main and helper budgets were included before choosing $\tau_2$. The $O(P^2)$ band pairs and the polynomial number of transfer steps are also included. Thus the far terms cost at most $D^{-C_3}\|S_a\|_{2,I}\|S_b\|_{2,I}$ per cell and group, for any required fixed $C_3$. All other filters are contractions. Forward just the transferred main term and retain its defects. After all transfers the main term is
\[
 H_I(\mathcal D_{a,s}S_a,S_b,B_c,B_d),
 \eqtag{eq:transferred-main}
\]
where $\mathcal D_{a,s}$ is a product of $O(\log L)$ helper averages. Restoring $B_c,B_d$ to $Y_c,Y_d$ gives the pair-small main term; each correction has an additional basic side $U$.

For estimates of the defects, retain the sum over $p$ at each fixed offset $q-p$. If the outgoing filter is a ball filter, split its terminal complementary band $1-Q_{c,P}$ into the shared identity and the positive ball term. All other band endpoints have usable ball ranges. If the outgoing filter is cumulative, its outer ball has already been removed, so its complementary band history is shared across groups. The other complement similarly either has a positive ball outermost or is shared. A shared history will not be counted once for every group.

\bp{transfer}{THT.Sides.transferComplementaryFilters}{helpers, localize, masks, graph, coeff}

The only maximal inequality needed for the comparison is the already proved Lemma \ref{lem:vector-maximal}:
\[
 \|M_qf\|_{L_y^pL_\omega^4}\le C_{p,q}\|f\|_{L_y^pL_\omega^4}\qquad(p>q,\ 2<q<3).
 \eqtag{eq:mixed-Mq}
\]
There is one inner norm, on the product of the fast torus and random parameter spaces. No weighted extrapolation or further mixed-norm theorem is used.

\subsection*{Comparison on a short depth block}
\begin{lemma}\label{lem:shortblock}
Let a finite dyadic tree piece have root $R$, internal summation cells in at most $K$ consecutive depths of the original $N$-depth range, and leaves in that block or one depth below. Suppose deterministic Hilbert-valued inputs, supported on $R$, satisfy $\|z_j\|_{2,L}\le s_j$ on every leaf, and scalar randomizations satisfy $(\E_\sigma|z_j^\sigma(y)|^a)^{1/a}\le C_a\|z_j(y)\|$ for each finite moment used. Then, for arbitrary measurable $|\epsilon_I(\sigma)|\le1$ and any subcollection of internal cells,
\[
 \E_\sigma\left|\sum_I|I|\epsilon_I(\sigma)H_I(z^\sigma)\right|
 \le C_q(1+C_N)(1+\log(K+2))^C(K+2)^{2(q-2)/(q-1)}|R|\prod_js_j.
 \eqtag{eq:shortblock}
\]
Two successive independent Rademacher randomizations of a vector stack satisfy the moment hypothesis, by Lemma \ref{lem:finite-fourier}.
\end{lemma}
\begin{proof}
Normalize $s_j=1$; a zero size gives a zero input. Absorb the phase of the whole random sum into the kernel signs. It suffices to bound the absolute expected signed sum, uniformly in those signs.

\emph{Leaf-depth rectangles.} Pad the at most $K+1$ leaf depths to a dyadic integer interval of length at most $4(K+1)$. Partition their four-tuples as follows. If the minimum is repeated, specify its attaining slots and designate two, $a,b$, in fixed order; the others have greater depth. Put $t$ equal to this minimum. If the minimum is unique, designate its slot $a$ and the smallest dyadic interval whose lower and upper halves contain the minimum and second minimum. Specify which other slots lie in the upper half, designate one of them $b$, and put $t$ equal to that half's left endpoint. Remaining slots lie above the upper half. These conditions are rectangles of leaf masks and partition all depth tuples.

There are $O(1+\log(K+2))$ length levels and boundedly many slot assignments. At one level and assignment the positions $\iota$ number $O(K+1)$; their selected leaf unions $E_{a,\iota}$ are disjoint, as are $E_{b,\iota}$. Treat the repeated-minimum cases as one additional singleton level. Fix a rectangle and suppress $\iota$. Every selected leaf has depth at least $t$, except possibly in slot $a$, called heavy; other slots are regular. Only summation depths $s<t$ contribute. A selected leaf of depth at most $s$ cannot meet an internal $s$-cell: dyadic nesting would make that cell a leaf or a descendant of one. The interior support margin preserves this exclusion after the small enlargements below.

\emph{Sampling and randomization.} Put $\delta=\beta r_t$ for a sufficiently small fixed dyadic $\beta$. Write
\[
 y_0=\delta(n_0+u),\quad y_1=\delta(n_1+v),\quad
 n_j=(1-j)n_0+jn_1,\quad w_j=(1-j)u+jv,
\]
initially with $u,v\in[0,1)$. On a fixed circle $\T_0$ of length $16$, set $F_j^\sigma(n,w)=z_j^\sigma(\delta(n+w))\one_{E_j}(\delta(n+w))$, restricted to $[0,1)$ in slots $0,1$ and to fixed lift intervals in $[-4,4]$ containing all needed $w_j$ in slots $2,3$, and zero elsewhere. With fixed Haar constants the fast progression integral reproduces the original integral without wraparound.

Sample the kernel at $(\delta n_0,\delta(n_1-n_0))$. The error is $O(\delta/r_s)$, with all coordinate windows still inside $I$ by the support margin and choice of $\beta$. Let $M_{j,I}$ be deterministic squared mass on these windows, with their bounded overlap counted. The leaf bounds give
\[
 M_{j,I}\le C|I|,\qquad
 \sum_{\operatorname{depth}(I)=s}M_{j,I}\le C|E_j|.
\]
The all-$L^2$ case of \eqref{eq:absolute} on local counting measure times Haar measure, followed by H\"older in $\sigma$, bounds the summed error at depth $s$ by $C(\delta/r_s)|E_a|^{1/2}|E_b|^{1/2}$. Here local counting measure assigns weight $\delta/|I|$ to a relevant integer. Since $\sum_{s<t}\delta/r_s\le C$, designated disjointness bounds all positions' errors by $C|R|$ per level.

For $c=(-1,3,-3,1)$ use the fast Fourier multipliers
\[
 \widehat{U_{j,\gamma}F}(n,m)=e(\gamma m^2/c_j)\widehat F(n,m),
 \qquad \gamma\in\R/(3\Z),
\]
and put $G_j=U_{j,\gamma}F_j^\sigma$. The fast form is unchanged: its Fourier support has $\sum m_j=\sum jm_j=0$, whose nullspace is spanned by $(c_j)$ and $(jc_j)$. Thus $m_j=c_j(h+jh')$ and $\sum m_j^2/c_j=0$. For each $n,\sigma$,
\[
 \E_\gamma\int_{\T_0}|G_j(n,w)|^4\,dw
 \le2\|F_j^\sigma(n,\cdot)\|_2^4.
 \eqtag{eq:fourth-moment}
\]
Indeed integration in $w$ equates the sums of two Fourier indices and integration in $\gamma$ equates their squared sums. The pairs agree as multisets, giving at most the two pairings. Finite Fourier sums prove both statements first; the same fourth-moment bound for differences and the $L^2$ progression estimate pass to limits.

Let $e_j(n)$ be the fast $L^2$ norm of the deterministic stack on the corresponding window. Minkowski and the sign moments give $(\E_\sigma\|F_j^\sigma(n,\cdot)\|_2^a)^{1/a}\le C_a e_j(n)$, and
\[
 \delta\sum_ne_j(n)^2\le C|E_j|,\qquad
 \frac{\delta}{|I|}\sum_{n\text{ relevant to }I}e_j(n)^2\le C.
 \eqtag{eq:window-mass}
\]
In a regular slot $e_j(n)\le C$: disjoint selected leaves meeting a window of length $O(\delta)$ have total length at most $O(\delta)+2r_t$, since all but two are contained in the window and each has length at most $r_t$. Bound their full squared masses by their lengths and divide by $\delta=\beta r_t$. The local bound in \eqref{eq:window-mass} also holds in the heavy slot because the relevant windows lie inside the internal cell $I$.

\emph{Delete large deterministic windows.} Set $A=K+2$. In a heavy slot delete precisely the integers with $e_a(n)>A$; do not clip any randomized function value. By unitarity and Minkowski, the $L^4_{\gamma,\sigma}$ norm of the deleted part's local $L^1$ norm is at most
\[
 C\frac{\delta}{|I|}\sum_{n\text{ relevant},\ e_a(n)>A}e_a(n)
 \le\frac{C}{A}\frac{\delta}{|I|}\sum_{n\text{ relevant}}e_a(n)^2.
\]
For a regular slot its local $L^3$ norm has bounded fourth random moment: normalized norm monotonicity, \eqref{eq:fourth-moment}, and $e_j\le C$ bound it by $C(\delta/|I|)\sum e_j(n)^4\le C$. The absolute estimate with exponents $(1,3,3,3)$ and H\"older in the random variables therefore bound the deletion error at $I$, after multiplication by $|I|$, by $(C/A)\delta\sum_{n\text{ relevant}}e_a(n)^2$.

Sum disjoint cells at a fixed depth using bounded window overlap, then at most $K$ depths: the cost is at most $CK|E_a|/A$. Summing positions uses the disjoint $E_a$ masks and costs at most $CK|R|/A\le C|R|$ per level. After deletion all remaining heavy window sizes are at most $A$. Deletion depends only on the deterministic windows and commutes with the fast Fourier maps. Continue to write $G_j$ for the retained functions.

\emph{Continuous comparison with one random norm.} For fixed $\gamma,\sigma$ and arbitrary $u,v\in\T_0$, define
\[
 g_j(y,w)=\sum_nG_j(n,w)\one_{\{|y-\delta n|<c\delta\}}(y),
 \qquad w=(1-j)u+jv,
\]
where $c<1/20$ is fixed. If four such intervals meet a real progression, their integer centers satisfy $n_j=(1-j)n_0+jn_1$: the discrepancies for $j=2,3$ are integers of absolute value less than one. The offsets form the fixed positive-area polygon $|(1-j)a_0+ja_1|<c$ for $0\le j\le3$. Sampling the continuous kernel at its center therefore reproduces a fixed positive multiple of the sampled lattice form.

The second kernel-sampling error has the earlier bound: unitarity and deletion decrease the local fast $L^2$ masses, and introduce no slow indices, so the same window masses and $\delta/r_s$ estimate apply. We can now use \eqref{eq:CN} with the original admissible kernels and original summation cells. For almost every fast and random parameter the finitely many interval values define bounded, compactly supported inputs. It remains to estimate
\[
 C_N\int_{\R}\E_{\gamma,\sigma}\int_{\T_0^2}
       \prod_jM_qg_j(y,(1-j)u+jv)\,du\,dv\,dy.
\]
Put $p_a=p_b=p_*=2q-2$ and $p_j=2p_*/(p_*-2)$ in the other slots. They exceed $q$ and their reciprocals sum to one. One H\"older inequality of exponent four in $(\gamma,\sigma,u,v)$, then spatial H\"older, bounds the last display by
\[
 C_N\prod_j\|M_qg_j\|_{L_y^{p_j}L_{\gamma,\sigma,w}^4}
 \le C_qC_N\prod_j\|g_j\|_{L_y^{p_j}L_{\gamma,\sigma,w}^4},
\]
using \eqref{eq:mixed-Mq}. Each map $(u,v)\mapsto(1-j)u+jv$ has Haar pushforward, so the inner norm is precisely the indicated single $L^4$ norm. Equation \eqref{eq:fourth-moment} and the sign moments give
\[
 \|g_j\|_{L_y^{p_j}L_{\gamma,\sigma,w}^4}
 \le C\left(\delta\sum_{n\text{ retained}}e_j(n)^{p_j}\right)^{1/p_j}.
\]
For regular slots this is at most $C|E_j|^{1/p_j}$ by \eqref{eq:window-mass}. For a heavy slot it is at most $CA^{1-2/p_*}|E_a|^{1/p_*}$. At a level with $J\le C(K+1)$ positions, designated disjointness and H\"older give
\[
 \sum_{\iota=1}^J|E_{a,\iota}|^{1/p_*}|E_{b,\iota}|^{1/p_*}
 \le J^{1-2/p_*}|R|^{2/p_*}.
\]
The other masks cost $|R|^{1-2/p_*}$. Consequently the block factor is
\[
 A^{1-2/p_*}(K+2)^{1-2/p_*}
 =(K+2)^{2(q-2)/(q-1)}.
\]
There are only logarithmically many levels; both sampling errors and deletion errors already cost $C|R|$ at each. This proves the lemma. All limiting steps follow from the difference version of \eqref{eq:fourth-moment}, the bounded deterministic deletion operators, and \eqref{eq:mixed-Mq}; slow index sets are finite.
\end{proof}
\bp{shortblock}{THT.Sides.shortBlockComparison}{finite-CN, finite-fourier, vector-maximal, absolute, fourier-foundation}


\subsection*{Power costs and frozen complements}
Use the componentwise square caps \eqref{eq:side-cap}. All additional integrated square budgets below are $P^C|R|$; Lemma \ref{lem:cuts} converts their sum into simultaneous path caps at an arbitrarily small fixed root-length cost. Polynomially many such budgets are absorbed by increasing $C$.

For three strong factors write $a_s,b_s$ for the group square sums of the two obligatory group-indexed side sizes. A third factor is shared or bounded by its group square sum $c_s$. The fourth retained factor is bounded by $P^C$. Cauchy--Schwarz in groups bounds the group-summed absolute term by $P^Ca_sb_sc_s$. For $2/3<p<1$,
\[
 \sum_s(a_sb_sc_s)^p
 \le\left(\sum_sa_s^{3p}\right)^{1/3}\left(\sum_sb_s^{3p}\right)^{1/3}\left(\sum_sc_s^{3p}\right)^{1/3}
 \le\prod_{x=a,b,c}\left(\sum_sx_s^2\right)^{p/2}.
 \eqtag{eq:three-sides}
\]
The last step is monotonicity of sequence norms since $3p>2$. Integrating the path caps gives total length-weighted power cost $P^C|R|$. A fourth strong factor uses its retained supremum. The group sum is formed \emph{before} taking this power.

The main term \eqref{eq:transferred-main}, with complements restored to their full $Y$ arguments, costs $\eta P^C|R|$ absolutely, by pair-smallness, contractions, and Cauchy--Schwarz in the two side sizes. Restoration errors have a third basic side and obey \eqref{eq:three-sides}. The numerical transfer errors are negligible by their arbitrary inverse-main-budget precision.

For the remaining paired defects in a block, let $(C_{G,s}^q)_q$ be the complementary band stack and let $D_{G,s}$ be the other complement. Freeze both at $b$. Their current local $L^2$ norms are bounded: pad bands are orthogonal at one cell, and all inner histories are contractions on bounded basic inputs. Differences from their values at $b$ have integrated square budgets $P^C|R|$, by Lemma \ref{lem:variation}, paying the polynomial number of bands and the $O(K)$ overlap. The terminal-band convention above ensures joint group budgets for positive ball ranges and shared budgets for group-independent histories.

A freeze error has a third strong factor. If its other frozen complement is unbounded, write that complement as its bounded current value plus its strong difference; both resulting terms have three or four strong factors. We can also discard a use whenever a frozen complement has local size greater than twice its current uniform bound at that use or at an ancestor in its block. If that uniform bound is $C_0$, the event indicator is bounded by $C_0^{-2}$ times the squared difference size. Lemma \ref{lem:variation} applies at every potential ancestor having a downstream certificate. Propagating the event through its block costs at most $K$ in integrated square energy. Its square-root indicator is therefore another strong factor; multiply by it when estimating the flagged term. Include these energies among the power-cost caps. On unflagged uses both frozen complements have bounded local sizes on their \emph{entire} ancestor hull within the block. The flags remove terms from the critical estimate, charging them to \eqref{eq:three-sides}; a flag itself is not a restart root.

\bp{power-costs}{THT.Sides.threeSidePowerAndFreezing}{side-caps, variation, transfer, base, cuts, absolute}

\subsection*{Fixed stacks for the two varying sides}
Fix a transfer step, search leaf, and near offset $q-p$. Put $m=O(\log L)$. We claim that the two sides can be expanded into a number of configurations absorbed in $G_0$, each using fixed vector stacks for $(G,b)$ with one position selected at a use, and with
\[
 \sum_{G,b}\|\mathrm{stack}_{a,G,b}\|_2^2\le PG_0|R|,
 \qquad \sum_{G,b}\|\mathrm{stack}_{b,G,b}\|_2^2\le G_0|R|.
 \eqtag{eq:stack-energy}
\]
The first stack includes the band index. These functions are fixed before the subsequent cuts.

For any sequence $W_s$, $b\le s<b+K$, expand
\[
 W_s=W_b+\sum_{[r,t]\in\mathcal P(b,s)}(W_t-W_r),
\]
where $\mathcal P(b,s)$ is the binary partition of the prefix $[b,s]$, with at most one interval at each length. At a fixed length these intervals are disjoint as depth intervals. A fixed stack component may be restricted to the union of use cells selecting it, which preserves its value there and decreases its norm. The base position and length choices have logarithmic cost.

For a basic $U$ side use each depth directly as a position. Its square energy controls the sum, and subsequent local helper products are contractions. On the banded side pay at most $P+1$ indices. For a fixed linear side $V_{G,b}$, denote a subsequent product by $T_s$. An outer band is a difference of two such products. Lemma \ref{lem:prefix}, telescoping the products on intervals of a fixed binary length, gives
\[
 \|T_bV_{G,b}\|_2^2+\sum_{[r,t]\text{ at fixed length}}
 \|(T_t-T_r)V_{G,b}\|_2^2
 \le(C\log(K+2))^{O(m)}\|V_{G,b}\|_2^2.
\]
The individual positive projections are nested. Couple pad choices across times and restore their averages by Jensen; expand complements. Sum lengths and use \eqref{eq:linear-energy}. Again only the banded side pays $P+1$.

For a rolling side put $W_s=T_s(Z_{G,s}-Z_{G,u_b})$. Its base term uses Lemma \ref{lem:variation} on $[u_b,b]$, whose overlap over blocks is bounded. Its increments split exactly as
\[
 W_t-W_r=T_t(Z_{G,t}-Z_{G,r})+(T_t-T_r)(Z_{G,r}-Z_{G,u_b}).
\]
The first term is controlled directly by that lemma; when estimating norms retain the entire endpoint cells containing selected use cells. They have downstream certificates and the enlargement is legitimate by locality. For the second term expand $[u_b,r]$ into binary intervals in a common $O(K)$-depth window. Fix an outer length, an inner length, and an inner interval $[a,v]$. For all outer positions $[r,t]$ using it, take the one fixed input $X=\one_\Omega(Z_{G,v}-Z_{G,a})$, where $\Omega$ is the union of the required $r$-cells. This input is independent of the outer position; since $v\le r$, locality gives the required action of $T_t-T_r$. Enlarge its support for a norm estimate to the union of containing $v$-cells, all with certificates. The fixed-input product-jump estimate bounds the outer-position sum by $(C\log(K+2))^{O(m)}\|X\|_2^2$. Now sum the squared inputs over groups and inner intervals at the fixed length by Lemma \ref{lem:variation}. Windows across blocks overlap a bounded number of times. Summing the two length choices completes \eqref{eq:stack-energy}. Search leaves, transfer steps, side-component choices, and length configurations have total cost $G_0$. Individual pads remain inside averages, not as extra configurations.

The two position decompositions need not be synchronized. For each configuration use independent Rademacher families $(\epsilon_i)$, $(\eta_j)$, $(\zeta_p)$ for the two position indices and the paired bands. Randomize the first fixed stack by $\epsilon_i\zeta_p$, the untouched side by $\eta_j$, and the frozen complementary band stack by $\zeta_p$. Multiply the kernel at a use by its selected position signs $\epsilon_{i(s)}\eta_{j(s)}$. Expectation selects exactly those positions and the matched band indices. The remaining transfer coefficient $(p-q)/P$ is $O(P^{-1})$ for the fixed offset. Two applications of Khintchine give the pointwise moment hypothesis in Lemma \ref{lem:shortblock}, with the deterministic stack norms. Arbitrary initial kernel signs are allowed.

\bp{fixed-stacks}{THT.Sides.fixedStackRepresentation}{prefix, variation, linear-side, base, helpers, finite-fourier}

\subsection*{Common cuts, hull maxima, and summation}
For an index $\nu=(G,b,\text{configuration})$, let $\Omega_\nu$ be the union of its use cells. Its maximal cells are disjoint roots; its hull consists of all cells between those roots and the uses. A hull maximal function is defined on \emph{all} of each hull interval, including missing branches. The total $\sum_\nu|\Omega_\nu|\le M_0|R|$ with $M_0\le\exp(k^{O(1)})$, by active multiplicity, the main group count, and the number of configurations.

Apply the hull part of Lemma \ref{lem:cuts} to the two energies in \eqref{eq:stack-energy}. Since $\log(2+M_0)$ is absorbed in $G_0$, their squared hull maxima $m_{a,\nu},m_{b,\nu}$ satisfy the joint caps $PG_0$ and $G_0$, respectively. These caps hold everywhere on the whole retained hull intervals, including missing and stopped branches. Combine these cuts with the power-cost cuts. The fixed stack functions are unchanged.

For each retained hull, perform a sparse decomposition for its two fixed side stacks: from a root $Q$ stop at first descendants where either squared stack average exceeds sixteen times its root average. The next roots have total length at most $|Q|/8$. A zero root size gives no contribution. The remainders $E_Q$ are disjoint within that hull construction and have size at least $7|Q|/8$. Complete pieces by children immediately below their included cells, including missing branches and common-cut boundaries. A child's squared average is at most twice its parent's; hence side leaf norms are bounded multiples of the root sizes $a_Q,b_Q$. Frozen complementary norms are bounded on full retained hull cells, and therefore on these leaves as well. Each piece has at most $K$ internal depths and leaves at most one step below, so Lemma \ref{lem:shortblock} applies to the randomized inputs.

On every $x\in E_Q$, the hull interval $Q$ itself shows $a_Qb_Q\le\sqrt{m_{a,\nu}(x)m_{b,\nu}(x)}$. Thus, including missing or stopped branches,
\[
 \sum_{\nu,Q}|Q|a_Qb_Q
 \le C\int_R\sum_\nu\sqrt{m_{a,\nu}m_{b,\nu}}
 \le C\int_R\sqrt{\sum_\nu m_{a,\nu}}\sqrt{\sum_\nu m_{b,\nu}}
 \le P^{1/2}G_0|R|.
\]
Multiplying by the $O(P^{-1})$ transfer coefficient, and absorbing logarithms and configurations in $G_0$, proves the critical absolute bound
\[
 C_q(1+C_N)P^{-1/2}G_0(K+2)^{2(q-2)/(q-1)}|R|.
 \eqtag{eq:near-band}
\]
Arbitrary initial signs imply the absolute sum of the original cell contributions, with the prescribed near-band sums retained inside each contribution. We have not taken absolute values of individual bands.

We have proved the following unconditional side estimate. Terms with at least three sides, freeze and restoration errors, and large-frozen-complement terms have length-weighted $p$th-power cost $P^C|R|$ for every fixed $2/3<p<1$. Transferred main terms have absolute cost $\eta P^C|R|$. Numerical errors are negligible. The remaining two-side terms obey \eqref{eq:near-band}. In the original active-cell minimum use
\[
 \min(\eta,x)\le\eta^{1-p}x^p\quad(x\ge0)
 \eqtag{eq:min-power}
\]
for the power costs, after the group sum; use \eqref{eq:min-subadd} for absolute errors and the critical absolute terms. All cuts have arbitrarily small fixed total root fraction. Every hypothesis is available anew on a normalized subroot by the base construction and Proposition \ref{prop:linear-side}; the estimates in a parent attempt use their original fixed functions.


\paragraph{Exported side estimate.}
The input consists of the finite base data, a fixed localization leaf, the root-uniform linear-side bound, and the instantiated helper catalog. The outputs are a first-cut family of any prescribed sufficiently small root fraction, the powered bounds for three strong factors and all freeze/restoration terms, the absolute transferred-main bound, and the coupled-band estimate \eqref{eq:near-band}. All use the same original cell summation and the same $C_N$. Stacks are functions fixed before the cut family is selected; restarting constructs new stacks on the new normalized root. These are precisely the assertions proved in this section, so no local-energy conclusion is added by restricting an arbitrary global family.
\bp{side-estimate}{THT.Sides.completeSideEstimate}{transfer, shortblock, power-costs, fixed-stacks, cuts, side-caps}



\section{Common heights and concentration of the core terms}\label{sec:core}
We now estimate the terms with three or four core roles. A basic $U$ role, when present, remains a side. In the scalar identity \eqref{eq:core-side-error}, all terms containing a physical error are negligible: use its integrated square bound \eqref{eq:physical-error}, the $P^C$ retained sizes, the main active multiplicity, and Cauchy--Schwarz. For terms with exactly two sides, replace the two core complements by their whole current $Z$ arguments; corrections have at least three sides. Section \ref{sec:sides} estimates all these terms. It remains to treat at least three cores.

\subsection*{An exact common-height selection}
Fix an actual cell $I=I_s$, group $G$, and an arbitrary point $x_*\in I$. Put $r=|I|$ and $R_s=r^{-1}$. There are $m_0\in\{3,4\}$ core roles. Sample heights independently and uniformly from \eqref{eq:upper-heights}. For each pure tuple of core modes, select its first sampled height at which every role is good. Each core mode has bad proportion at most $2\delta_*$, so success probability is at least $1-2m_0\delta_*>0$. The first-success rule has total weight exactly one for every tuple, hence gives an exact identity in expectation for the finite multilinear expansion. Explicitly, if its one-trial success proportion is $p_\tau$, the sum of the weights of all successful words of length $n+1$ with their first success last is $(1-p_\tau)^np_\tau$. Its sum over $n\ge0$ is one. Thus this statement can be defined entirely through finite words and a convergent geometric series; no countable product probability space is required.

If success is at trial $n+1$, assign each preceding failed trial to its first failing role in a fixed order. There are at most $m_0^n$ assignments. For a fixed assignment, the conditions in a role form a product mask: failure at assigned trials and success at earlier roles' required tests. The upper-height coherence proved in Section \ref{sec:modes} means these masks select unions of entire core subtotals $A_D$ in minimal CURRENT components. Group the selected modes at the final successful components $F_j$. They lie in the unique good predecessor of each $F_j$. Let $a_j$ be its hereditary representative, and set $\xi=\vartheta/L_{h-\Delta},\quad \gamma_j^*=\omega_{a_j}(x_*)$. A PRE path to the representative has at most $D_*$ edges. It supplies an integer shift of size at most $D_*E_{h-\Delta}$, with slope error at most $D_*\Xi/r_{v_b}$. The pre-lag makes this error times $r$ smaller than any prescribed fixed inverse power of $D_*$. The quadratic Taylor error is likewise small, since each selected mode was admitted by $u_b$ and had curvature at most $D_*$ at its admission scale. Absorbing constants at $x_*$ into coefficients, the selected physical argument is, up to uniform error $D_*^{-150}$ times its absolute coefficient sum,
\[
 e(c_jp_{\rho_j}(y_j))e(\gamma_j^*(y_j-x_*))\,g_j(\xi(y_j-x_*)).
 \eqtag{eq:core-torus-form}
\]
Here $g_j$ is the true periodic cutoff times a finite trigonometric polynomial. The cutoff is periodic on the $\vartheta$ torus and therefore on the $\xi$ torus, since $L_{h-\Delta}$ is an integer. Use the true cutoff first when invoking \eqref{eq:gauge-cancel}, and then its Fourier truncation. The modes, cutoff factors, and common progression multiplier, in any anchor coordinates, have all relevant Fourier indices in
\[
 |n|_\infty,|m|_\infty\le B_h',\qquad B_h'=D_*^5E_{h-\Delta}.
 \eqtag{eq:frequency-box}
\]
This includes the cutoff indices $B_*L_{h-\Delta}$. Absolute Fourier costs, apart from the selected mode coefficients themselves, are fixed powers of $D_*$. Dilation of indices does not change coefficient sums. The fixed-power reserves in the choice of $A_0$, the Fourier cutoffs, and $K$ were taken to include these bounds.

With three core roles and side role $l$, insert an independent offset cutoff $\zeta^{\rho_l}$ divided by its constant expectation $c_\zeta\ge1/2$. Put $S'=\zeta^{\rho_l}e(-c_lp_{\rho_l})S_{\phys}\quad\text{on }I$. Its normalized $L^2$ norm is at most that of the side. On all true cutoff supports the four gauge factors have the bounded smooth periodic multiplier in \eqref{eq:gauge-cancel}. In coordinates $u=\xi(y_{l'}-x_*)$, $v=\xi t$, denote its shifted truncated version by $R_z(u,v)$. It has uniformly bounded supremum, including for different offsets and anchors.

\bp{common-height}{THT.Core.firstSuccessfulHeight}{coherence, representatives, mode-schedule, finite probability, geometric series}

\subsection*{Torus norms and Fourier tests}
A successful component contains a nonexceptional primary. Apply its gap \eqref{eq:gap} to pure $\vartheta$ shifts in all the boxes under consideration. Let $K_B$ be its rational tiny-frequency space in the horizontal coordinates and set
\[
 \mathcal H=(K_B\cap\Z^{\dim\theta_z})^\perp\subseteq\T^{\dim\theta_z}.
\]
Lemma \ref{lem:torus-annihilator}, applied to $K_B$, proves that this is a connected torus and gives its exact character integrals. For every character difference needed below, $r n\cdot\xi$ is tiny if $n\in K_B$ and huge otherwise.

Averaging a character along $\xi(I-x_*)$ gives $1$ plus a tiny error in the first case and a tiny error in the second, by direct integration of an exponential. Haar integration on $\mathcal H$ gives exactly $1$ or $0$. Expand $|g_j|^2$ in characters and use the coefficient bounds: its two averages differ by at most $D_*^{-100}$ times the squared selected absolute coefficient sum. Together with \eqref{eq:core-torus-form}, almost orthogonality of minimal CURRENT components, and the floors in \eqref{eq:component-size}, this gives
\[
 a'_j:=\|g_j\|_{L^2(\mathcal H)}
 \le C\left[\sum_{D\text{ selected}}
 \left(\|A_D\|_{2,I}+D_*^{-30}\sum_{m\text{ core in }D}|\alpha_m|\right)^2\right]^{1/2}.
 \eqtag{eq:torus-size}
\]
For example, the squared total coefficient error is bounded by $D_*$ times the sum of squared component coefficient sums; the floor $D_*^{-60}$ dominates this after the stronger error power. The physical norm is bounded by the augmented one. Thus no coefficient norm replaces the principal component amplitude.

Any two distinct progression coordinates on $\mathcal H^2$ are independent Haar variables: their integer matrix has nonzero determinant and is surjective on a connected torus. The pair-coordinate argument of Section \ref{sec:reduction} therefore gives
\[
 \int_{\mathcal H^2}\prod_{j=0}^3|h_j(u+(j-l')v)|\,du\,dv
 \le\prod_{j=0}^3\|h_j\|_{L^2(\mathcal H)}.
 \eqtag{eq:Haar-four}
\]
This follows by the positive H\"older interpolation proved in Lemma \ref{lem:absolute}, applied to the six two-$L^1$, two-$L^\infty$ bounds; a bounded multiplier has its supremum cost. The zero-dimensional torus is included.

Put $U_0=D_*^{250}$. For four cores take anchor $l'=0$ and write
\[
 Q=\sum_j\gamma_j^*,\qquad T=\sum_jj\gamma_j^*.
\]
Seek integer $n,m$ in the full box \eqref{eq:frequency-box} with
\[
 q'=Q+n\cdot\xi,\qquad t'=T+m\cdot\xi,\qquad\max(|q'|,|t'|)\le U_0R_s.
\]
Choose one reference by a fixed rule independent of all masks and sampling histories. Use dyadic bands, with fixed endpoint conventions,
\[
 \mu\asymp\min\left(1,\min_{0\le j\le3}\frac{|t'-jq'|}{R_s}\right),\qquad
 \nu\asymp\max\left(1,\frac{|q'|}{R_s},\frac{|t'|}{R_s}\right).
 \eqtag{eq:four-tests}
\]
Discard a tuple if no reference exists or the inner minimum is below $U_0^{-1}$.

To justify this test and its error, let $W_s$ be the normalized kernel in anchor coordinates and expand
\[
 D(u,v)=R_z(u,v)\prod_jg_j(u+jv).
\]
Its characters sample $\widehat W_s$ at $r(Q+n\cdot\xi,T+m\cdot\xi)$. By the tiny/huge gap, two moderate pairs have the same pair of character cosets modulo $K_B$. Their frequencies differ by a tiny quantity, so replace the Fourier values by the reference value. Other cosets have rapidly decreasing Fourier tails. The retained coset coefficient is
\[
 \int_{\mathcal H^2}D(u,v)e(-n\cdot u-m\cdot v)\,du\,dv,
\]
bounded by $C\prod_ja'_j$ by \eqref{eq:Haar-four}. The four marginals give zeros of $\widehat W_s(q,t)$ on the four lines $t=jq$. Smoothness near these lines and rapid decay give, for every fixed $J$,
\[
 |\widehat W_s(rq',rt')|\le C_J\mu\nu^{-J}.
\]
Indeed use the mean value theorem toward a nearest line and rapid derivative decay; the fixed-angle geometry changes the distance by only an absolute factor. Away from the lines use rapid decay directly. No reference or a distance below $U_0^{-1}$ gives the small absolute errors estimated below.

For three cores use anchor $l$ and put
\[
 Q=\sum_{j\ne l}\gamma_j^*,\qquad T_l=\sum_{j\ne l}(j-l)\gamma_j^*.
\]
Choose a mask-independent reference $m$ in the box with $|T_l+m\cdot\xi|\le U_0R_s$. Put $\beta=r(T_l+m\cdot\xi)$, and discard when no reference exists or $|\beta|<U_0^{-1}$. Set $\mu\asymp\min(1,|\beta|)$ and $\nu\asymp\max(1,|\beta|)$. Fourier integration in $t$ gives a smooth anchor function $W_\beta(h)$ supported in the normalized interior, all of whose fixed-order derivatives are bounded by $C_J\mu\nu^{-J}$. The zero at $\beta=0$ is precisely the anchor marginal; integration by parts gives the decay.

Only the reference $m$ coset contributes appreciably. With
\[
 D(u,v)=R_z(u,v)\prod_{j\ne l}g_j(u+(j-l)v),\qquad
 D_m(u)=\int_{\mathcal H}D(u,v)e(-m\cdot v)\,dv,
\]
duality and \eqref{eq:Haar-four} give $\|D_m\|_2\le C\prod_{j\ne l}a'_j$. Collect horizontal Fourier coefficients by cosets and choose one reference $n$ from the box per coset. Parseval bounds their $\ell^2$ norm by this quantity. Distinct real reference frequencies $rn\cdot\xi$ are more than $1$ apart; replacing the other frequencies of a coset by its reference has tiny error.

Expand $S'$ in normalized Fourier coefficients $\widehat{S'}_I(k)$ for $e(k(y-x_*)/r)$, $k\in\Z$. The matrix of pairings with $W_\beta$ has entries bounded by $C_{J,A}\mu\nu^{-J}(1+|k+r(Q+n\cdot\xi)|)^{-A}$. In a shell of radius $\Lambda$, each row and column has $O(\Lambda)$ entries, since both the integers and the reference frequencies are separated. The Schur bound, with $A$ increased, gives the contribution
\[
 C_J\mu\nu^{-J}\left(\prod_{j\ne l}a'_j\right)
 \sum_{\substack{\Lambda=1,2,4,\ldots\\\Lambda\le U_0}}\Lambda^{-J}b_\Lambda,
 \eqtag{eq:side-windows}
\]
where
\[
 b_\Lambda^2=\E_{\rho_l}\sum_{\substack{k\in\Z:\,|k+r(Q+n\cdot\xi)|\le2\Lambda\\
 \text{for some }n\in\Z^{\dim\theta_z},\ |n|_\infty\le B_h'}}|\widehat{S'}_I(k)|^2.
\]
We averaged the fresh offset and used Cauchy--Schwarz; its expectation normalization costs at most $2$. Windows use the entire box and do not depend on the mode masks. The same fresh offset may be coupled for all tested triples at a cell. Shells above $U_0$ are negligible by the same estimate or ordinary $L^2$ pairing.

Every discarded or replacement error above is bounded by
\[
 CD_*^{-100}\prod_{j\text{ core}}\sum_{m\text{ selected in }j}|\alpha_m|,
 \eqtag{eq:core-numerical}
\]
with the local side size as an extra factor when present. The Fourier tails and tiny frequency errors have stronger precision reserves; at the discarded near-zero or high-frequency thresholds the factor $U_0^{-1}=D_*^{-250}$ covers the fixed Fourier coefficient costs. The finite number of dimensions is already in those budgets. These errors do not pay for failure histories: after expanding absolute coefficient sums, each pure tuple has first-success weight one. The path multiplicities, coefficient sums, and retained size caps make their integrated sum $O(e^{-k}|R|)$. Estimate them absolutely, before any power below one.

For a fixed final height and component tuple, an assignment of the previous failures gives
\[
 a'_j\le C\sigma_{F_j},\qquad \E(a'_j)^2\le C(2\delta_*)^{n_j}\sigma_{F_j}^2,
\]
where $n_j$ is the number assigned to role $j$. Apply the probability bound separately to each minimal-component summand in \eqref{eq:torus-size}; dropping extra success requirements only enlarges its probability. H\"older with $m_0$ factors, using the deterministic bound to pass from second to $m_0$th moments, yields
\[
 \E\prod_ja'_j\le C^{m_0}(2\delta_*)^{n/m_0}\prod_j\sigma_{F_j}.
\]
The sum over $m_0^n$ assignments and all $n$ converges for the chosen small absolute $\delta_*$. Consequently, after fixing the final height, role types, dyadic amplitude and frequency bins, every principal cell contribution is bounded by a sum of quantities
\[
 C_J\mu\nu^{-J}\Lambda^{-J}\left(\prod_{j=0}^3\sigma_j\right)N_s^G.
 \eqtag{eq:principal-core}
\]
For a core role $\sigma_j$ bins \eqref{eq:component-size}; for a side it bins $b_\Lambda$. For four cores take $\Lambda=1$. The integer $N_s^G$ counts the representative tuples passing the nonnegligible tests, at most once per tuple and scale. Zero factors contribute nothing. The proof will keep $N_s^G$ \emph{inside} its power.

\subsection*{The finite sparse-count lemma}
\bp{frequency-tests}{THT.Core.principalFrequencyTests}{common-height, torus-annihilator, component-sizes, absolute, fourier-foundation, mode-schedule}

\begin{lemma}\label{lem:sparse-count}
There are absolute constants $C,C_{\rm res}$ with the following property. Among roles $0,1,2,3$, let $\mathcal J$ consist of either all four roles or three core roles. In each core role let $A_i$ be a finite set of distinct real labels, $|A_i|\le M$, with $M,L\ge1$. Fix $\zeta\in\R^d$, $B\ge1$, and $0<\mu\le1\le\nu$. For finitely many distinct integers $s$, put $R_s=R_02^s$ and let $E_s\subseteq\prod_{i\in\mathcal J}A_i$, $N_s=|E_s|$. Assume:
\begin{enumerate}
\item Two labels in distinct core roles determine at most one full core tuple in all the $E_s$, with at most $L$ scale occurrences.
\item If $l$ is a side role, there are $w_s\ge0$, $\sum_sw_s\le M$, such that each fixed core label is incident to at most $w_s$ tuples at scale $s$.
\item If distinct same-role labels $a,a'$ occur at scales $s,t$, then
\[
 |a-a'+n\cdot\zeta|>\Sigma\max(R_s,R_t)\qquad(n\in\Z^d,\ |n|_\infty\le Q).
\]
\item For every triple $T\subseteq\mathcal J$, with excluded role $l'$ and $\alpha_i=i-l'$, each tuple at scale $s$ has a witness $n\in\Z^d$, $|n|_\infty\le B$, such that
\[
 \mu R_s\le\left|\sum_{i\in T}\alpha_i a_i+36n\cdot\zeta\right|\le\nu R_s.
\]
\item At every occupied scale,
\[
 |n\cdot\zeta|\notin[\mu R_s/1000,1000\nu R_s]\qquad(n\in\Z^d,\ |n|_\infty\le Q).
\]
\end{enumerate}
If $Q\ge C_{\rm res}(d+1)^2B$ and $\Sigma\ge C_{\rm res}(d+1)^2\nu$, then, with $\ell=1+\log_2(2+(d+1)^2\nu/\mu)$,
\[
 \sum_{s:N_s\le\delta M^2}N_s\le C\ell L^{12}\delta^{1/300}M^2
 \qquad(0<\delta\le1).
 \eqtag{eq:sparse-count}
\]
There is no dependence of the constants or exponents on the dimension, frequencies, or number of scales.
\end{lemma}
\begin{proof}
We may perform cardinality arguments on a common fine grid. Indeed there are only finitely many labels, scales, and reserved shifts. Approximate all real coordinates on one grid, taking their approximants on a fixed multiple of its step so divisions by $1,2,3$ stay on the grid. For sufficiently small step, separation holds with $\Sigma/2$, annular bounds with $\mu/2,2\nu$, and the pure-shift gap still contains the narrower interval with endpoints $\mu R_s/100$ and $100\nu R_s$ after these relaxations. Round widths upward by at most one step. Choosing the step after all finite data preserves these fixed margins; the final estimate is independent of it. Increase the absolute reserves and retain the original letters for the relaxed grid parameters.

\emph{Many four-step paths.} In a bipartite graph with parts of size at most $M$, at least $\eta M^2$ edges, and a proper coloring with at most $M$ colors, there is a subset of the first part of size at least $c\eta M$ such that every ordered pair of its vertices has at least $c\eta^5M^3$ four-step paths. Repetitions are allowed. To prove this, pad the second part to $m=\lceil M\rceil\le2M$, pick one of its vertices uniformly, and let $Y$ be its neighborhood. Then $\E|Y|^2\ge(\E|Y|)^2\ge\eta^2M^2/4$. Call an ordered first-part pair bad if its codegree is less than $c_0\eta^2M$. Its chance to be in $Y^2$ is its codegree divided by $m$, so the expected number of bad pairs is at most $c_0\eta^2M^2$. For sufficiently small $c_0$, some $Y$ has $|Y|\ge c\eta M$ and fewer than $|Y|^2/100$ bad pairs, by maximizing $|Y|^2-100(\text{bad pairs})$. Delete vertices with more than $|Y|/4$ bad partners; fewer than $|Y|/25$ are deleted. Two retained vertices have at least $c|Y|$ common good partners in $Y$. Each endpoint and such a partner have at least $c_0\eta^2M$ common neighbors; choosing the two neighbors gives the claimed four-step paths. For fixed endpoints their ordered four-color sequences are distinct, since a vertex and an edge color determine the next edge uniquely.

\emph{Blocks and difference covers.} It suffices to treat sufficiently small absolute $\delta$, since total mass is at most $LM^2$. Put $\tau=\delta^{1/300}$; then $\delta\le\tau^2/2$. Restrict to $N_s\le\delta M^2$, order scales by decreasing $R_s$, and form consecutive counting blocks of mass at most $\tau^2M^2$. Except possibly the last, each has mass greater than $(\tau^2-\delta)M^2$, so $J_0\le3L\tau^{-2}$. These are not the $K$-depth blocks of Section \ref{sec:modes}.

For a block $b$ let $R_b^+$ be its largest scale. With grid step $h_0$, set $N_0=\lceil32B\rceil$, $\epsilon_b=h_0\lceil16\nu R_b^+/h_0\rceil\le17\nu R_b^+$, and
\[
 U_b(t)=\{6n\cdot\zeta+u: n\in\Z^d,\ |n|_\infty\le tN_0,\ u\in h_0\Z,\ |u|\le t\epsilon_b\}.
\]
These finite symmetric sets contain zero and satisfy $U_b(t)+U_b(t')\subseteq U_b(t+t')$.

Fix a core role $i$ in a block. While occurrences with uncovered $i$-label have mass greater than $\tau^3M^2$, project to pairs in $i$ and another core role $j$, retaining one witness scale per pair. There are at least $(\tau^3/L)M^2$ edges. Color by a third core label $k$. The coloring is proper by pair determination. Put $\eta=\tau^3/L$ and apply the graph observation to get $B_0\subseteq A_i$ with $|B_0|\ge c\eta M$. Alternately add and subtract the four annular relations along a four-step path between $a,a'\in B_0$. The intermediate $i$ label and both $j$ labels cancel. Dividing by $\alpha_i\in\{\pm1,\pm2,\pm3\}$ gives
\[
 a-a'\in q(c_1,c_2,c_3,c_4)+U_b(1),
\]
where $q$ depends only on the four colors. The scalar error is at most $4\nu R_b^+$ before division; the shift has index at most $4B$ before division, and $36/\alpha_i$ is a multiple of $6$.

Take a maximal packing in $B_0-B_0$ with pairwise differences outside $U_b(2)$, and choose endpoints for each packed difference. Each pair has at least $c\eta^5M^3$ four-color sequences, and two packed differences cannot share a sequence, since that puts their difference in $U_b(1)-U_b(1)$. There are at most $M^4$ color sequences. Thus the packing has at most $C\eta^{-5}M$ elements, and maximality covers $B_0-B_0$ by that many translates of $U_b(2)$. Remove $B_0$ and repeat. Each step removes at least $c\eta M$ labels, so there are at most $C\eta^{-1}$ subsets. We have proved: apart from at most $\tau^3M^2$ uncovered occurrences in each role and block, the incident labels are covered by at most $C\eta^{-1}$ subsets $B$, each with difference-cover size at most
\[
 K_1M,\qquad K_1=C\eta^{-5}.
 \eqtag{eq:two-cover-counts}
\]
\emph{Nonfuture pairs.} Discard occurrences missing an own-block covered label in any core role, losing at most $4J_0\tau^3M^2$. An allowed core label is future if it belongs to a cover subset of that role in a later block. Nonfuture label sets are disjoint over blocks, in each role: a later covered occurrence would make the earlier one future. If their counts are $m_{b,i}$, then $\sum_bm_{b,i}\le M$. A side role is always nonfuture, with $m_{b,l}=\sum_{s\in b}w_s$ and the same sum bound. An occurrence with fewer than three future roles has two nonfuture roles; assign it one such pair. Its block mass is at most $\min(\tau^2M^2,Lm_{b,i}m_{b,j})$, by pair multiplicity or the side incidence bound. Therefore its total over blocks is at most
\[
 \sqrt L\,\tau\,M\sum_b\sqrt{m_{b,i}m_{b,j}}\le\sqrt L\,\tau\,M^2.
\]
There are at most six pairs.

All remaining occurrences have three future core roles. Assign a triple and one future cover subset containing each label by fixed rules. There are at most $CJ_0^4\eta^{-3}$ combinations. Fix one with mass $\rho M^2>0$, current block, and future subsets $B_1,B_2,B_3$. The subscripts now enumerate these roles without changing their progression coefficients $\alpha_i$.

\emph{Three mass ranges.} Put $\Gamma=C_{\rm gap}(d+1)^2\nu/\mu$. If $\rho\le C\delta\ell$, keep this bound. Otherwise cut the ordered mass into thirds, omitting the at most two scale bins crossing a cut. Each range has mass at least $\rho M^2/4$ when $C$ is large. Let $R_\#$ be the largest scale in the middle range. The first range gives anchor sets $B_i^0\subseteq B_i$ of size at least $c\rho M/L$: a fixed label has at most $LM$ occurrences. Their witnesses have scales at least $R_\#$, so distinct anchors are separated modulo the reserved shifts by more than $\Sigma R_\#$.

The last range gives a set $F$ of distinct triples with $|F|\ge c\rho M^2/L$, taking one witness for each. Let $R_F$ be the largest scale of the last range. Delete from the middle range the scales below $\Gamma R_F$; there are at most $1+\lceil\log_2\Gamma\rceil\le C\ell$ of them. The middle range began with at least $\rho/(4\delta)$ occupied bins, so $\rho>C\delta\ell$ leaves at least $c\rho/\delta$. Greedily select scales with successive ratios at least $\Gamma$. Each selection removes at most $C\ell$ bins, leaving $q\ge c\rho/(\delta\ell)$ selected scales $R_v$. Choose one actual triple $t_v$ at each. Every $R_v$ exceeds every last-range and future-block scale by at least $\Gamma$. These are all the margins needed below; no separate intervening mass range is required. We compare the $q|F|$ indexed outputs of $t_v+F$ with a common slab cover.

\emph{One small-growth scale.} Write $U_i$ for the neighborhood of the future block associated with $B_i$. Subdividing the coefficient and scalar intervals covers $U_i(t)$ by at most $(2t+1)^{d+1}$ translates of $U_i(1)$. Hence some odd $H\le1000(d+1)^2$ satisfies simultaneously
\[
 |U_i(H+2)|\le2|U_i(H)|\quad(i=1,2,3).
\]
Otherwise the product of the three cardinalities would more than double at every odd step, whereas it is at most $(2t+1)^{3(d+1)}\prod_i|U_i(1)|$. At $t=1000(d+1)^2+1$ this contradicts $2^{500(d+1)^2}>(2000(d+1)^2+3)^{3(d+1)}$. The preliminary exponential dimension cost thus locates $H$ but never enters a subsequent estimate.

The sets $a+U_i(H)$ for $a\in B_i^0$ are disjoint. An intersection would express their difference in $U_i(2H)$, with shift index $O((d+1)^2B)$ and scalar width at most $C(d+1)^2\nu R_\#$, contrary to the reserves on $Q,\Sigma$. Put $Z_i=B_i^0+U_i(H)$. Since $B_i^0\subseteq B_i$, the difference cover gives
\[
 |Z_i-B_i|\le K_1M|U_i(H+2)|\le2K_1M|U_i(H)|
 \le C(K_1L/\rho)|Z_i|.
\]
Choose a nonempty $X_i\subseteq Z_i$ minimizing $r_i=|X_i-B_i|/|X_i|$; thus $r_i\le r_0:=CK_1L/\rho$.

Here is the required minimal-growth iteration. For a finite translate set $C_0=\{c_1,\ldots,c_q\}$, let $X_{i,j}$ consist of those $x\in X_i$ whose translate $x+c_j$ is not in an earlier $X_i+c_v$. Then $|X_i+C_0|=\sum_j|X_{i,j}|$. Every point of $(X_i\setminus X_{i,j})-B_i+c_j$ already lies in an earlier translate of $X_i-B_i$. The number of new points at step $j$ is therefore at most
\[
 |X_i-B_i|-|(X_i\setminus X_{i,j})-B_i|
 \le r_i|X_i|-r_i|X_i\setminus X_{i,j}|=r_i|X_{i,j}|.
\]
For an empty complement the same inequality is immediate. Summing and iterating with $C_0=-(m-1)B_i$ proves $|X_i-mB_i|\le r_i^m|X_i|\quad(m\ge1)$, where $mB_i$ is the $m$-fold sumset.

One anchor translate contains a portion $Y_i$ of $X_i$ of size at least $|X_i|/M$, and $Y_i-Y_i\subseteq U_i(2H)$. For a packing $\mathcal P\subseteq mB_i$ with distinct differences outside this neighborhood, the translates $Y_i-p$ are disjoint in $X_i-mB_i$. Hence $|\mathcal P|\le r_i^mM$. A maximal packing is also a cover by translates of $U_i(2H)$. We use this only for $m=2,3$.

\emph{A slab cover with two label-count factors.} Consider outputs in $2B_1\times2B_2\times2B_3$ satisfying
\[
 \left|\sum_i\alpha_i x_i+36n\cdot\zeta\right|\le2\nu R_\#
 \quad\text{for some }n\in\Z^d,\ |n|_\infty\le2B.
\]
Cover the first two sumsets by at most $r_0^2M$ translates each of $U_i(2H)$. Above a fixed pair of boxes the third coordinate lies in $2B_3$ intersected with a set of the form
\[
 c+\{n\cdot\zeta:|n|_\infty\le C(d+1)^2B\}
 +[-C(d+1)^2\nu R_\#,C(d+1)^2\nu R_\#].
\]
Substitute the first two boxes in the slab equation and divide by $\alpha_3\in\{\pm1,\pm2,\pm3\}$. The factors $6$ in boxes and $36$ in the slab ensure that all divided shifts are still integral.

Let $\mathcal P$ pack this local third-coordinate set outside $U_3(2H)$. The set $B_3^0+\mathcal P\subseteq3B_3$ is still such a packing: with equal anchors use the definition of $\mathcal P$; with distinct anchors, a collision would express their difference modulo a reserved shift with magnitude at most $C(d+1)^2\nu R_\#$, contradicting separation. It follows that $|B_3^0|\,|\mathcal P|\le r_0^3M,\quad |\mathcal P|\le Cr_0^3L/\rho$. Maximality gives a local cover. Thus the whole slab is covered by at most
\[
 Cr_0^7(L/\rho)M^2\le CK_1^7(L/\rho)^8M^2
 \eqtag{eq:slab-cover}
\]
translates of $\prod_{i=1}^3U_i(2H)$. The anchors are what remove a third factor $M$.

\emph{Separation of translated triples.} All $t_v+F$ lie in this slab by adding their annular relations. A cover box contains at most one indexed output. Within one translate, different members of $F$ differ in a label coordinate, whose last-range witness scale exceeds every future scale; hereditary separation excludes that difference from $U_i(4H)$. For two different translates let $R_{\rm big}$ be the larger middle-range scale. Its annular residual has magnitude between $\mu R_{\rm big}$ and $\nu R_{\rm big}$, whereas the other middle-range and both last-range residuals have total magnitude at most $\mu R_{\rm big}/10$ by the $\Gamma$ margins. Thus the output difference has progression linear form, after a shift of index $O(B)$, between $9\mu R_{\rm big}/10$ and $2\nu R_{\rm big}$ in magnitude. If the outputs lay in one box, their coordinate differences would belong to $U_i(4H)$; the total scalar error would be at most $C(d+1)^2\nu\max_iR_{{\rm future},i}^+\le\mu R_{\rm big}/10$. All accumulated shift indices are $O((d+1)^2B)\le Q$. We would obtain a pure shift of magnitude between $4\mu R_{\rm big}/5$ and $3\nu R_{\rm big}$, forbidden by the gap. This also excludes equality of two indexed outputs.

Consequently $q|F|$ is bounded by \eqref{eq:slab-cover}. Substitute the lower bounds for $q,|F|$ and cancel $M^2$:
\[
 \rho^{10}\le C\delta\ell K_1^7L^9.
\]
Since $K_1=C\eta^{-5}$, a convenient weaker bound, also covering $\rho\le C\delta\ell$, is $\rho\le C\ell\eta^{-4}L\delta^{1/10}$. There are at most $CJ_0^4\eta^{-3}$ combinations. Add their mass to the uncovered mass and the nonfuture-pair mass. Using $\delta^{1/10}=\tau^{30}$ gives
\[
 M^{-2}\sum_{s:N_s\le\delta M^2}N_s
 \le C\left(J_0\tau^3+\sqrt L\,\tau
              +\ell L J_0^4\eta^{-7}\tau^{30}\right).
 \eqtag{eq:count-balance}
\]
Since $J_0\le3L\tau^{-2}$ and $\eta=\tau^3/L$, the last term is at most $C\ell L^{12}\tau^{-8-21+30}=C\ell L^{12}\tau$. The first two terms obey the same bound. This proves \eqref{eq:sparse-count} for small $\delta$; total mass $LM^2$ covers the remaining $\delta$ after increasing $C$.
\end{proof}
\bp{sparse-count}{THT.Counting.sparseProgressionMass}{four-paths, sumset-growth, finite packing, integer arithmetic}


\subsection*{Verification of the counting hypotheses}
Fix the path point $x_*$, final height, all role types, and the frequency and amplitude bins in \eqref{eq:principal-core}. In a group put
\[
 M_i^G=|A_i^G|\quad(i\text{ core}),\qquad
 M_l^G=\sigma_l^{-2}\sum_{s:G\text{ used}}\|S_{l,G,s}\|_{2,I_s}^2\quad(l\text{ side}),
\]
and $M'=\max_iM_i^G$. Omit zero-count groups. Then $M'\ge1$; in a side role a nonzero use has $b_\Lambda\ge\sigma_l$, so Parseval also gives $M_l^G\ge1$. Identify a representative with its fixed real value $\omega_a(x_*)$. Distinct labels in a fixed role and type have distinct values by \eqref{eq:representative-separation} at zero shift.

Two core labels determine the full tested core tuple across all scales. For four cores, subtract the two reference equations for tuples sharing these labels. The two remaining label differences solve a two-by-two system for $Q,T$, up to bounded shifts and an error $O(\nu\max(R_s,R_t))$. Its determinant has magnitude at most $3$. Separation \eqref{eq:representative-separation} forces both differences to vanish. For three cores the single $T_l$ equation determines the remaining label in the same way. Divisions by $1,2,3$ are permitted because $\xi=6^\Delta\vartheta/L_h$ and the shift reserves include those fixed denominators. The larger reciprocal scale in separation is what permits this argument across different cells.

Moreover, for every triple of core roles, with excluded anchor $l'$,
\[
 c\mu R_s\le\left|\sum_{i\text{ in triple}}(i-l')\omega_{a_i}(x_*)+n\cdot\xi\right|
 \le C\nu R_s,\qquad |n|_\infty\le CB_h'.
 \eqtag{eq:triple-test}
\]
For four cores this is $t'-l'q'$ from \eqref{eq:four-tests}; for three it is the nonzero $T_l$ test. Fixing a pair fixes the entire tuple and these label values. Lemma \ref{lem:height}, applied to them and $\xi$, bounds its scale occurrences by $L_0\le P^C(1+\log(\nu/\mu))^C$, with $L_0\ge1$. The rational height logarithms and dimensions are polynomial in $P,k$, absorbed in $P^C$ for the chosen $P$. This proves the total pair bound for two core roles.

For a core/side pair, fix the core label and a scale. The windows defining $b_\Lambda$ for different incident triples are disjoint. Otherwise their $Q$ values differ by a bounded shift plus $O(\Lambda R_s)$; together with the two $T_l$ equations this determines the two other labels by a two-by-two system, again with divisor at most $3$. Since $\nu,\Lambda\le CU_0\ll\Xi$, separation forces the triples equal. Couple the fresh offset for all windows. Parseval for each offset, then expectation, bounds the number of triples in the fixed side bin by $w_s=\sigma_l^{-2}\|S_{l,G,s}\|_{2,I_s}^2$. This is valid on every subset of scales, and $\sum_sw_s=M_l^G\le M'$. Summing over fixed core labels gives, for all distinct roles $i,j$,
\[
 \sum_sN_s^G\le L_0 M_i^G M_j^G,\qquad N_s^G\le(M')^2.
 \eqtag{eq:total-count}
\]

Apply Lemma \ref{lem:sparse-count} with $M=M'$, $L=L_0$, $d=\dim\theta_z$, and $\zeta=\xi/36$. Absorb the absolute constants in \eqref{eq:triple-test} into $\mu,\nu$. Pair determination, multiplicity, and side incidence have just been verified. For the separation reserve,
\[
 \zeta=6^{\Delta-2}\vartheta/L_h,\qquad Q=\lfloor E_h/6^{\Delta-2}\rfloor,\qquad\Sigma=\Xi/2.
\]
The choice of $A_0$ and $B_h'=D_*^5E_{h-\Delta}$ gives $Q\ge C_{\rm res}(d+1)^2CB_h'$. Also $\Xi=D_*^{1000}$ and $\nu\le CD_*^{250}$ give the required $\Sigma$ reserve. Finally any successful component contains a nonexceptional primary, even if its representative is exceptional. For that primary all pure-horizontal shifts in this finite lemma have rational height below $H$. By \eqref{eq:gap} they are at most $T^{-1}/r_{v_b}$ or at least $T/r_{t_b}$. Since $U_0^{-1}\lesssim\mu\le1\le\nu\lesssim U_0$ and actual scales lie between these depths, the chosen $T$ excludes the entire pure-shift annulus. Every finite-lemma hypothesis is therefore satisfied, and
\[
 \sum_{s:N_s^G\le\delta(M')^2}N_s^G
 \le C\ell L_0^{12}\delta^c(M')^2,\qquad c=1/300.
 \eqtag{eq:applied-sparse}
\]
Keep this coefficient logarithmic in the frequency bands: $\ell=1+\log_2(2+(d+1)^2\nu/\mu)$. There is no need to replace it by a power of $\mu^{-1},\nu$, or $\Lambda$.

\bp{tested-counts}{THT.Core.testedTupleCounts}{frequency-tests, representatives, tops, sparse-count, height, side-caps}

\subsection*{One distribution calculation and the final sums}
\begin{lemma}\label{lem:distribution}
Let $M>0$. Suppose $0\le N_s\le M^2$, $\sum_sN_s\le A$, and
\[
 \sum_{s:N_s\le\delta M^2}N_s\le B\delta^cM^2\quad(0<\delta\le1),
\]
where $A,B>0$ and $0<c<1$. For $0<\varepsilon<c$,
\[
 \sum_sN_s^{1-\varepsilon}
 \le\frac{c}{c-\varepsilon}A^{1-\varepsilon/c}B^{\varepsilon/c}
 M^{2\varepsilon(1/c-1)}.
 \eqtag{eq:distribution}
\]
\end{lemma}
\begin{proof}
Omit zero terms and set $F(t)=\sum_{0<N_s\le t}N_s$. Then $F(t)\le\min(A,BM^{2-2c}t^c)\quad(t>0)$. For $t\le M^2$ use the sparse bound with $\delta=t/M^2$; for $t>M^2$, $F(t)=F(M^2)\le BM^2$. Tonelli and $N^{1-\varepsilon}=\varepsilon N\int_N^\infty t^{-1-\varepsilon}dt$ give $\sum N_s^{1-\varepsilon}=\varepsilon\int_0^\infty t^{-1-\varepsilon}F(t)dt$. Split at $t_0=(A/(BM^{2-2c}))^{1/c}$. The integral is bounded by
\[
 \frac{\varepsilon BM^{2-2c}}{c-\varepsilon}t_0^{c-\varepsilon}
 +At_0^{-\varepsilon}
 =\frac{c}{c-\varepsilon}A^{1-\varepsilon/c}B^{\varepsilon/c}M^{2\varepsilon(1/c-1)}.
\]
Both integrals converge since $0<\varepsilon<c$.
\end{proof}
\bp{distribution}{THT.Counting.distributionPower}{Tonelli, one-variable power integration}


First discard negligible amplitudes in the absolute estimate. At each cell, group, height, type, and frequency choice, each label tuple has one amplitude tuple and there are at most $D_*^4$ label tuples. These choices and the active path multiplicity have fixed-power $D_*$ bounds, so their length-weighted incidence count is at most $D_*^{C_1}|R|$. Retained amplitudes satisfy $\sigma_j\le P^{C_0}\le D_*^{C_0}$. Since $\mu\nu^{-J}\Lambda^{-J}\le1$, terms with some $\sigma_j<D_*^{-B}$ sum to at most $C_JD_*^{-B+C_1+3C_0}|R|$. Choose the fixed integer $B>C_1+3C_0+2$; this is $O(e^{-k}|R|)$ and leaves the minimum by \eqref{eq:min-subadd}, before any power is taken.

The remaining amplitudes lie in $[D_*^{-B},P^{C_0}]$. Their quadruples number at most $C(1+\log D_*+\log P)^4\le P^C$. Thus a uniform per-bin estimate suffices.

Choose once and for all
\[
 \varepsilon=\frac{c}{4}=\frac1{1200},\qquad p=1-\varepsilon.
 \eqtag{eq:fixed-power}
\]
Apply \eqref{eq:distribution} to \eqref{eq:total-count} and \eqref{eq:applied-sparse}, for any pair $i\ne j$, to obtain directly
\[
 \sum_s(N_s^G)^p
 \le\frac43(L_0M_i^GM_j^G)^{3/4}(C\ell L_0^{12})^{1/4}(M')^{(1-c)/2}.
 \eqtag{eq:powered-count}
\]
Temporarily order the roles so $\sigma_1\ge\sigma_2\ge\sigma_3\ge\sigma_4$, and use pair $1,2$ in \eqref{eq:powered-count}. The path estimates give $\sum_GM_i^G\le CP^C\sigma_i^{-2}$ in every role. Expand $M'$ into its four possible maximizing roles. The total count exponent is $3/4+3/4+(1-c)/2=2p>1$. For nonnegative sequences and $b=\sum_i\beta_i\ge1$,
\[
 \sum_G\prod_i(m_i^G)^{\beta_i}
 \le\prod_i\left(\sum_Gm_i^G\right)^{\beta_i}.
\]
Indeed normalize each nonzero sum to one and apply H\"older with exponents $b/\beta_i$, omitting zero exponents; then $\sum_G(m_i^G)^b\le1$. Apply this to the preceding count factors, bounding the maximizing role by $CP^C\sigma_4^{-2}$. After multiplication by $(\prod_i\sigma_i)^p$, the amplitude factor is
\[
 \sigma_1^{-1/2-c/4}\sigma_2^{-1/2-c/4}
 \sigma_3^{1-c/4}\sigma_4^{3c/4}
 =\left(\frac{\sigma_2}{\sigma_1}\right)^{1/2+c/4}
  \left(\frac{\sigma_3}{\sigma_2}\right)^{1+c/2}
  \left(\frac{\sigma_4}{\sigma_3}\right)^{3c/4}\le1.
\]
The finite amplitude-bin count absorbs all their sums, including the $24$ orderings, into $P^C$. In particular neither the extra $D_*$ count cap nor a separate common-amplitude-scale argument is needed here.

The remaining frequency coefficient from \eqref{eq:powered-count} is at most $P^C(1+\log(\nu/\mu))^C$. The powered kernel factor is $\mu^p\nu^{-Jp}\Lambda^{-Jp}$. Extend the actual truncated bands to all $\mu=2^{-u}$, $\nu=2^v$, $\Lambda=2^w$ with $u,v,w\ge0$. Their sum is bounded by
\[
 C P^C\sum_{u,v,w\ge0}2^{-pu}2^{-Jpv}2^{-Jpw}(1+u+v)^C<\infty
\]
for a fixed sufficiently large derivative order $J$. For four cores omit the $w$ sum. Heights, types, role patterns, and search leaves have an allowed polynomial cost.

If $C_I$ is the sum of the absolute principal quantities \eqref{eq:principal-core}, subadditivity of $x^p$ and \eqref{eq:min-power} give
\[
 \min(\eta,C_I)\le\eta^\varepsilon\sum_{\text{configurations at }I}
 (\text{principal quantity})^p.
\]
The estimates just proved bound its path sum by $C_q\eta^\varepsilon P^{C_4}$. Integrating over $x_*$ uses \eqref{eq:path}, including when the point-dependent principal upper bounds vary within a cell. The original cell minimum was applied before this subadditivity, and $N_s^G$ remained inside its power. With the already separated absolute numerical errors, the core contribution is therefore
\[
 C_q\eta^{1/1200}P^{C_4}|R|+O(e^{-k}|R|).
 \eqtag{eq:core-final}
\]


\paragraph{Exported core estimate.}
For a retained attempt, the assertions \eqref{eq:total-count} and \eqref{eq:applied-sparse} are uniform in each path point, group, height, role type, amplitude tuple and frequency band. The distribution calculation is applied before the tuple count is split, while negligible amplitudes are removed before any power below one. The output \eqref{eq:core-final} has absolute exponent $1/1200$ and exports an absolute upper bound for $C_4$. Its degree is therefore available before the final choice of $q$.
\bp{core-estimate}{THT.Core.poweredCoreEstimate}{tested-counts, distribution, tops, side-caps, frequency-tests, path-packing, budget-rules}



\section{Normalized absorption and the operator estimate}\label{sec:completion}
\begin{proposition}[Normalized absorption]\label{prop:normalized}
There is $q\in(2,3)$ such that, for every sufficiently small fixed $c>0$, there is $A<\infty$, independent of $N$, with
\[
 \sum_{I\text{ used}}|I|\,|H_I(f)|\le(A+cC_N)|I_*|
 \eqtag{eq:normalized}
\]
on every finite rooted dyadic subtree with root $I_*$ and $\|f_j\|_{q,I}\le1$ on all its cells. The used cells may be any subset in at most $N$ consecutive depths.
\end{proposition}
\begin{proof}
All summation collections are finite and belong to the original $N$ depths, so every use of $C_N(q)$ in Lemma \ref{lem:shortblock} concerns the same finite constant.

\subsection*{One attempt and its restarts}
For a noninitial base tuple and active pattern, construct in order the main catalogs, colors and ball histories; the uniform first-match and mode parameter bounds; the offsets, by the measurable majorant on pre-existing main uses; the corresponding fixed modes and linear sides; and the helper catalogs on whole side states. All subsequent cuts act on these fixed objects. Marking failures have the length bound already proved, and every size, top-count, and hull cut is covered by Lemma \ref{lem:cuts}. The strong stack families cost $G_0$; polynomially many configurations in $P,k$ occur only in the permitted power-cost estimates. Choose their thresholds so the union of exceptional roots has length at most $|R|/4$.

At each such root restart on the original base summands, leaving the basic arguments and inherited lists unchanged. Their normalized sizes and rolling bounds hold there. Rebuild modes and sides using their root-uniform proofs; keep parent-attempt functions and extensions for the parent summands. A large-frozen-complement flag is still a powered error, not by itself a restart. Whole retained hulls and unchanged fixed stacks are exactly the conventions in Lemma \ref{lem:cuts}. The root packing proved in Section \ref{sec:reduction} gives total length at most $4|R|/3$. On a finite tree this recursion terminates, so a retained-attempt estimate sums with this fixed factor.

\subsection*{The quantitative inequality}
Use \eqref{eq:core-final} for at least three cores and \eqref{eq:three-sides} with $p=1-1/1200$ for at least three sides. With exactly two sides, restore the complements to whole current outputs and apply the transfer argument; all corrections have three sides. Its restored $Y$ pair gives the small main term, while \eqref{eq:near-band} and numerical errors are estimated absolutely. This is the decomposition established at the start of Section \ref{sec:core}.

After absorbing all power-cost exponents in one absolute $C_4$, the retained-attempt estimate is
\[
 C_q\left[\eta^\varepsilon P^{C_4}
 +(1+C_N)P^{-1/2}G_0(K+2)^{2(q-2)/(q-1)}\right]|R|
 +O(e^{-k}|R|),\qquad \varepsilon=1/1200.
 \eqtag{eq:before-absorption}
\]
Here the smaller main cost $\eta P^C$ is included in the first term. Fixed powers of $k$ can be included there as well. The three relevant types of bounds are
\[
 \log D_{\rm base},\ \log D_m\le C_qk^C,
\]
\[
 \log D_*,\ \log H,\ \log T\le C_qP^C,
 \qquad K\le C_qP^{C_5},
\]
\[
 G_0\le\exp\bigl(C_q(1+\log(2k))^C\bigr).
\]
All displayed degrees, including $C_4,C_5$, have absolute upper bounds before choosing $q$.

The absolute degrees follow from the local detection estimate and the multiplicative rank recurrence in Lemma \ref{lem:closure}; $q$ affects their constants, not their degrees. Main counts precede pair accuracy and edge height. The mode construction fixes full Fourier norms before cutoffs, then height and shift ranges, the horizontal gap, and finally $K$. Helpers use whole side states, not individual modes, so no mode count enters their recurrence. The two essential small-overhead facts are Lemma \ref{lem:lift-composition}, which compresses only main projections, and \eqref{eq:bad-density}, which removes the large height factor from the side-event chain length. These give precisely the $G_0$ bound above.

Choose $2<q<3$ sufficiently close to $2$ that
\[
 2C_5\frac{q-2}{q-1}<\frac18.
\]
After $C_4$ includes every powered side cost, set
\[
 \alpha=\frac{\varepsilon}{4(C_4+1)},\qquad P=\lceil\eta^{-\alpha}\rceil.
\]
For every such fixed $\alpha$ and sufficiently large $k$, $P$ dominates every fixed power of $k$ and $G_0\le P^{1/8}$. The block estimate and the strict choice of $q$ give
\[
 P^{-1/2}G_0(K+2)^{2(q-2)/(q-1)}\le C_qP^{-1/4}\le C_q\eta^{\alpha/4}.
\]
Also $P\le2\eta^{-\alpha}$, hence
\[
 \eta^\varepsilon P^{C_4}\le2^{C_4}\eta^{\varepsilon-\alpha C_4}
 \le2^{C_4}\eta^{\varepsilon/2}.
\]
Increase the starting logarithm to meet every earlier reserve and overhead requirement. Since $0<\theta<1$, $e^{-k}$ is smaller than the resulting quality gains. With $c_* =\min(\varepsilon/2,\alpha/4)>0$, \eqref{eq:before-absorption} and root packing prove
\[
 \sum_{I\text{ active}}|I|\min(\eta,|H_I(z)|)
 \le C_q\eta^{c_*}(1+C_N)|R|.
 \eqtag{eq:active-final}
\]
No groupwise pair-smallness assertion was used: every minimum was first on the original cell sum, with group and prescribed band sums in their stated order.

\subsection*{Summing the update levels}
For a later update level $t\ge12$, there are $O(t)$ choices of the earlier increment level and only a bounded number of slot patterns. Each tuple intersects only a bounded number of inheritance segment families, whose root lengths sum uniformly by the base reset packing. Further attempts preserve that bound. Thus the coefficient of $1+C_N$ from later levels is at most a fixed multiple of
\[
 \sum_{t\ge12}t\exp(-c_*k_t^\theta),\qquad k_{t+1}=k_t^A.
\]
This series converges and tends to zero as $k_1$ increases: $k_t=k_1^{A^{t-1}}$ grows faster than a geometric sequence, and for $k_1$ above any fixed threshold its summands are dominated by a convergent geometric sequence in $t$ whose first coefficient tends to zero. Absolute errors $O(e^{-k_t})$ have the same summable multiplicity. The finitely many leading levels have fixed budgets and bounded active multiplicity, so contribute a finite constant independent of $N$.

Given the small $c$ in Proposition \ref{prop:normalized}, take $k_1$ so the coefficient of $C_N$ is at most $c$, and include every remaining constant in $A$. First take the base expansion at finite accuracy cutoff; on any fixed finite tree its remainders tend to zero by \eqref{eq:stagger} and bounded $L^2$ sizes. The estimates just proved are uniform in that cutoff, so passage to the limit gives \eqref{eq:normalized}. This proves the asserted normalized estimate.


\end{proof}
\bp{normalized}{THT.LocalForm.normalizedAbsorption}{base, localize, linear-side, tops, side-estimate, core-estimate, cuts, budget-rules}

\begin{theorem}[Uniform local estimate]\label{thm:uniform-local}
For the $q$ supplied by Proposition \ref{prop:normalized}, there is $C_q<\infty$ such that $C_N(q)\le C_q$ for every $N$. Equivalently, all the finite local sums in \eqref{eq:CN} have this common bound, uniformly in the lattice, admissible kernels, and compactly supported $L^q$ inputs.
\end{theorem}
\begin{proof}
Take disjoint roots at the largest length in the given finite collection. At a root $J$ set $a_{j,J}=\|f_j\|_{q,J}$. If a root size vanishes then the corresponding input vanishes almost everywhere on $J$, so all its descendant forms vanish. For every other root, stop at the first strict descendants $J'$ at which some average of $|f_j|^q$ exceeds $16a_{j,J}^q$. For a fixed $j$ these first stops are disjoint and
\[
 16a_{j,J}^q\sum_{J'\text{ stopped for }j}|J'|
 <\sum_{J'\text{ stopped for }j}\int_{J'}|f_j|^q
 \le a_{j,J}^q|J|.
\]
Their union over four $j$ has length at most $|J|/4$. On the intervening rooted subtree, dividing each input by $16^{1/q}a_{j,J}$ makes every normalized size at most one. Proposition \ref{prop:normalized} bounds the contribution of this piece by
\[
 16^{4/q}(A+cC_N)|J|\prod_j a_{j,J}.
\]
Repeat on the stopped roots, using their own root sizes; the finite original depth range makes the recursion finite. The sets $E_J$ left after removing the next stopping roots are disjoint over the recursive family, have $|E_J|\ge3|J|/4$, and satisfy $M_qf_j\ge a_{j,J}$ there. Thus
\[
 \sum_J|J|\prod_j a_{j,J}
 \le\frac43\sum_J\int_{E_J}\prod_jM_qf_j
 \le\frac43\int_\R\prod_jM_qf_j.
\]
The coefficient $C=\frac43 16^{4/q}$ is bounded by an absolute constant for $2<q<3$. Taking the supremum of the finite ratios defining $C_N$, and remembering that its definition also imposes $C_N\ge1$, gives
\[
 C_N\le1+CA+CcC_N.
\]
Choose $c$ with $Cc<1/2$ and absorb; this is valid because $C_N$ is already finite by \eqref{eq:CN}. The result is $C_N\le2+2CA$, independent of $N$.
\end{proof}
\bp{uniform-local}{THT.LocalForm.uniformEstimate}{normalized, absolute, finite-CN, dyadic first stopping, permitted scalar maximal theory}

\begin{theorem}[The trilinear Hilbert transform estimate]\label{thm:main}
There is a finite $C$ such that, for Schwartz functions on $\R$,
\[
 \left\|\operatorname{p.v.}\int_\R f_1(x-t)f_2(x-2t)f_3(x-3t)\,\frac{dt}{t}\right\|_1
 \le C\prod_{j=1}^3\|f_j\|_3.
\]
The principal value exists pointwise for Schwartz inputs. It has a unique bounded trilinear extension from $L^3(\R)^3$ to $L^1(\R)$.
\end{theorem}
\begin{proof}
We first reconstruct the kernel with all four required marginal cancellations. Lemma \ref{lem:gevrey} gives an even compactly supported Gevrey cutoff $\chi$ equal to one near zero. Set $\varphi(S)=S^3\chi(S)$, so $\varphi$ is odd and $\varphi^{(3)}(0)=6$. Let $\omega$ be the nonnegative compactly supported Gevrey partition in that lemma, and put
\[
 W(X,S)=\prod_{j=0}^3(\partial_S-j\partial_X)[\omega(X)\varphi(S)].
\]
For each $j$, its corresponding derivative is differentiation along $S\mapsto(u-jS,S)$. The other constant-coefficient derivatives commute with it and preserve compact support. Integration in $S$ therefore proves
\[
 \int_\R W(u-jS,S)\,dS=0\qquad(u\in\R,
 \ j\in\{0,1,2,3\}).
\]
The partition sums and all its differentiated sums are locally finite. Summing in the translate index and commuting the derivatives thus gives
\[
 \sum_{m\in\Z}W(x/a-m,t/a)=\varphi^{(4)}(t/a).
\]
For $t>0$, substitution $s=t/a$ and the fundamental theorem of calculus give
\[
 \int_0^\infty\varphi^{(4)}(t/a)\,\frac{da}{a^2}
 =\frac1t\int_0^\infty\varphi^{(4)}(s)\,ds
 =-\frac{\varphi^{(3)}(0)}t.
 \eqtag{eq:kernel-reconstruction}
\]
Oddness of $\varphi^{(4)}$ gives the same equality for $t<0$. The integral of the absolute value is at most $C/|t|$.

Write $a=b2^n$, where $1\le b<2$ and $n\in\Z$. Then $da/a^2=(1/a)\,db/b$ on each piece. Fix $b$. Let $D$ be a sufficiently large fixed dyadic integer and let $r=Da$. A kernel $a^{-1}W(x/a-m,t/a)$ can be written $r^{-1}w_I(x,t)$, with $I$ a cell of length $r$ in one of three nested lattices. Here are explicit lattices and the support check. At depths $s\in\Z$ with lengths $r_s=r_*2^{-s}$ use boundaries
\[
 r_s\bigl(\Z+(-1)^s j/3\bigr),\qquad j\in\{-1,0,1\}.
\]
They are nested, because $2(-1)^sj/3-(-1)^{s+1}j/3=(-1)^sj$ is an integer. At any one depth the three boundary sets interlace with spacing $r_s/3$. The four progression coordinates on the support of the translated $W$ all lie in one interval of length $C_Wa$, with center at distance $O(a)$ from $ma$. Choose $D$ so that $C_Wa<r/12$. Enlarge this coordinate interval by $r/24$ at each endpoint. The enlarged interval has length less than $r/6$, so it meets at most one of the interlacing boundaries. In at least two lattices it has no boundary, and in either such lattice the original coordinate interval has margin at least $r/24$ inside one cell. Choose one cell by a fixed ordering of the lattices.

At a fixed scale a chosen cell can receive only boundedly many integer translates: its possible $ma$ belong to an interval of length at most $r+O(a)$, so their number is at most $D+O(1)$. Regard these as a fixed bounded number of separate kernel families. Define $w_I=(r/a)W(x/a-m,t/a)$ for each such family. Rescaling preserves the four zero marginals. Its derivatives of order $u$ have bounds
\[
 |\partial^u w_I|\le D^{u+1}(C(u+2))^{C(u+2)}r^{-u}.
\]
The fixed factor $D$ can be incorporated into the admissibility constants chosen at the start. The preceding local theorem was proved for arbitrary fixed admissibility parameters. Apply it to this resulting fixed class, and denote its exponent by $q$. Thus its bound is uniform in $b$.

Take a finite scale range and a bounded compactly supported $f_0$, with $f_1,f_2,f_3$ Schwartz. Only finitely many spatial translates are relevant to $f_0$ in that range, and all other progression coordinates then lie in a fixed compact interval. Truncate those three inputs there, so the inputs belong to the defining class of $C_N$. Apply the uniform local estimate to each lattice and translate family, then integrate in $b$ with measure $db/b$. By H\"older and the scalar maximal theorem,
\[
 \int_\R\prod_{j=0}^3M_qf_j
 \le\|f_0\|_\infty\prod_{j=1}^3\|M_qf_j\|_3
 \le C_q\|f_0\|_\infty\prod_{j=1}^3\|f_j\|_3,
\]
since $q<3$. Consequently the four-linear form with the integrated scale kernel has this bound, uniformly in the scale range.

After summing the spatial translates that kernel is odd. For $x$ in the compact support of $f_0$, differentiability and bounded derivatives of the Schwartz functions give
\[
 \left|\prod_{j=1}^3 f_j(x+jt)-\prod_{j=1}^3 f_j(x-jt)\right|\le C|t|
 \quad(|t|\le1).
\]
This follows by the three-term telescoping product identity and the mean value theorem for each factor; the other two factors are uniformly bounded. At large $|t|$, each $x\pm jt$ leaves every compact set uniformly in such $x$, and Schwartz decay makes the products divided by $|t|$ integrable. These facts also show directly the pointwise existence of the symmetric principal value for each real $x$.

For a finite scale range, the integrated kernel is bounded by $C/|t|$, independently of the range. Pair its positive and negative $t$ contributions. The preceding difference bound near zero and the integrable tail supply a common dominating function, also integrable in $x$ against $|f_0|$. Dominated convergence and \eqref{eq:kernel-reconstruction} recover $-\varphi^{(3)}(0)$ times the principal-value four-linear form. Reflection $t\mapsto-t$ changes the sign and gives the operator in the theorem. Division by the nonzero reconstruction constant yields
\[
 \left|\int f_0(x)T(f_1,f_2,f_3)(x)\,dx\right|
 \le C_q\|f_0\|_\infty\prod_{j=1}^3\|f_j\|_3.
\]
The operator is measurable as a pointwise limit of measurable truncated integrals and is locally bounded by the same local estimates used for principal-value existence. Take $f_0=\one_{[-R,R]}\overline{T}/|T|$, with value zero where $T=0$. This is a permitted bounded compactly supported test and makes the left side $\int_{[-R,R]}|T|$. Monotone convergence as $R\to\infty$ proves the desired $L^1$ inequality.

Finally Schwartz functions are dense in $L^3(\R)$; this is the specialization of \nolinkurl{SchwartzMap.denseRange_toLpCLM} in \cite{mathlibSchwartz}. Choose Schwartz sequences $f_{j,n}\to f_j$ in $L^3$. For $n,m$, trilinearity gives the exact difference
\[
 T(f_{1,n},f_{2,n},f_{3,n})-T(f_{1,m},f_{2,m},f_{3,m})
 =T(f_{1,n}-f_{1,m},f_{2,n},f_{3,n})
\]
\[
 \hspace{1em}+T(f_{1,m},f_{2,n}-f_{2,m},f_{3,n})
 +T(f_{1,m},f_{2,m},f_{3,n}-f_{3,m}).
\]
The proved inequality and boundedness of the approximating $L^3$ norms show that this is a Cauchy sequence in $L^1$. Completeness defines its limit. The same three-term estimate, with two different approximation sequences, proves independence of choice. Passage to the limit proves the inequality and trilinearity for general $L^3$ inputs. Any bounded extension agrees on dense Schwartz triples and hence, by the same estimate, everywhere. This proves uniqueness.
\end{proof}
\bp{main}{THT.trilinearHilbertTransform_L3_L3_L3_L1}{uniform-local, gevrey, permitted scalar maximal theory, Schwartz density, dominated convergence}

\appendix
\section{Declaration index and representation requirements}
The order in the following index is the order in which the declarations are established in the text. A dependency named by one of these identifiers occurs earlier in the index. Other entries denote the explicitly permitted foundations, elementary arithmetic or a locally proved finite counting step. The proposed Lean names are attached to the proofs themselves. The index is a dependency aid, not a report of a Lean build.

For a faithful translation, keep the following mathematical data in the relevant records. Finite trees carry their actual cells, ancestor relation, depth interval and disjoint partition identities. An inherited column system carries births, private-coordinate injections, finite masks and the path budget. A row family carries both its Bessel estimate and its exact input-locality equality. Group ranges carry their whole endpoint cells and downstream-use certificates; an indicator on a later use cell acts on the input of an adjoint endpoint projection, not on the entire allowed range space. The frozen-mode record keeps first-match choices immutable and keeps the bounded lift in its indicated operator order. A side stack is chosen before its hull cuts. These conditions are the hypotheses proved and used above, not extra regularity assumptions.

Every parameter estimate is an existential statement returning its exponent with the quantifier order in the first section. In particular a finite-$N$ theorem may use $C_N$ only after its finiteness proof, and an energy estimate used in the final absorption must expose whether its logarithmic budget is $D_m$ or $D_*$. All zero-size cases are removed by the stated vanishing arguments; sums with empty index sets use zero and empty products use one. Finite numerical choices and classical witness selections are allowed. No computable search for an unrestricted real chart parameter is asserted.

Measurability is preserved in the finite constructions by measurable scalar columns, fixed finite-dimensional Gram matrices, finite linear combinations and finite decisions on already specified cells. The branch-offset selection uses its explicit measurable majorant, so it does not require measurability of a chosen inverse witness. The first-success expansion is the finite-word/geometric-series formula in the core section. Limiting arguments occur only with their displayed dominated, monotone, $L^2$-continuity or completeness justifications. An eventual Lean project still has to encode these proofs; this document neither supplies a Lean executable nor uses an unproved analytic axiom to stand for one of its main constructions.

\begin{longtable}{@{}p{0.25\textwidth}p{0.68\textwidth}@{}}
\textbf{Declaration} & \textbf{Earlier dependencies}\\\hline
\endfirsthead
\textbf{Declaration} & \textbf{Earlier dependencies}\\\hline
\endhead
\hyperlink{bp:budget-rules}{budget-rules} & library arithmetic, real exponential limits\\[3pt]
\hyperlink{bp:gevrey}{gevrey} & library calculus\\[3pt]
\hyperlink{bp:fourier-foundation}{fourier-foundation} & \hyperlink{bp:gevrey}{gevrey}, library Fourier and integration by parts\\[3pt]
\hyperlink{bp:cube-definition}{cube-definition} & library Haar, finite products, nonnegative roots\\[3pt]
\hyperlink{bp:finite-fourier}{finite-fourier} & \hyperlink{bp:cube-definition}{cube-definition}, library finite sums, Holder, Minkowski\\[3pt]
\hyperlink{bp:cube-foundation}{cube-foundation} & \hyperlink{bp:cube-definition}{cube-definition}, library Haar, Holder\\[3pt]
\hyperlink{bp:vector-maximal}{vector-maximal} & permitted vector Hardy--Littlewood theorem\\[3pt]
\hyperlink{bp:integer-basis}{integer-basis} & integer Bezout identity, determinants, induction\\[3pt]
\hyperlink{bp:reduced-basis}{reduced-basis} & \hyperlink{bp:integer-basis}{integer-basis}, Gram--Schmidt, finite lattice intersections\\[3pt]
\hyperlink{bp:transference}{transference} & \hyperlink{bp:reduced-basis}{reduced-basis}, determinants, compact convex sets\\[3pt]
\hyperlink{bp:lattice-box}{lattice-box} & \hyperlink{bp:reduced-basis}{reduced-basis}\\[3pt]
\hyperlink{bp:torus-annihilator}{torus-annihilator} & \hyperlink{bp:integer-basis}{integer-basis}, Haar uniqueness, torus Fourier basis\\[3pt]
\hyperlink{bp:absolute}{absolute} & \hyperlink{bp:torus-annihilator}{torus-annihilator}, change of variables, Holder\\[3pt]
\hyperlink{bp:finite-CN}{finite-CN} & \hyperlink{bp:absolute}{absolute}, scalar maximal definition\\[3pt]
\hyperlink{bp:rational-linear}{rational-linear} & determinants, finite-dimensional projections\\[3pt]
\hyperlink{bp:path-packing}{path-packing} & finite Tonelli, geometric series\\[3pt]
\hyperlink{bp:cuts}{cuts} & \hyperlink{bp:path-packing}{path-packing}, permitted weak maximal inequality\\[3pt]
\hyperlink{bp:column-data}{column-data} & finite dyadic data, library Hilbert projections\\[3pt]
\hyperlink{bp:coeff}{coeff} & \hyperlink{bp:column-data}{column-data}, library Hilbert projections\\[3pt]
\hyperlink{bp:prefix}{prefix} & library orthogonal projections, dyadic index decomposition\\[3pt]
\hyperlink{bp:compression}{compression} & \hyperlink{bp:coeff}{coeff}, \hyperlink{bp:prefix}{prefix}, \hyperlink{bp:path-packing}{path-packing}, permitted scalar maximal theory\\[3pt]
\hyperlink{bp:height}{height} & \hyperlink{bp:rational-linear}{rational-linear}, finite independence, geometric series\\[3pt]
\hyperlink{bp:chart-definition}{chart-definition} & \hyperlink{bp:gevrey}{gevrey}, finite rational matrices\\[3pt]
\hyperlink{bp:chart-operations}{chart-operations} & \hyperlink{bp:chart-definition}{chart-definition}, \hyperlink{bp:rational-linear}{rational-linear}, \hyperlink{bp:gevrey}{gevrey}\\[3pt]
\hyperlink{bp:four-paths}{four-paths} & finite probability, Cauchy--Schwarz, deletion count\\[3pt]
\hyperlink{bp:BSG}{BSG} & \hyperlink{bp:four-paths}{four-paths}\\[3pt]
\hyperlink{bp:sumset-growth}{sumset-growth} & finite minimal choice, injection counting\\[3pt]
\hyperlink{bp:Chang}{Chang} & \hyperlink{bp:finite-fourier}{finite-fourier}, Jensen, finite maximal subsets\\[3pt]
\hyperlink{bp:almostperiods}{almostperiods} & \hyperlink{bp:finite-fourier}{finite-fourier}, finite cosets, Markov\\[3pt]
\hyperlink{bp:dense-Bohr}{dense-Bohr} & \hyperlink{bp:almostperiods}{almostperiods}, \hyperlink{bp:Chang}{Chang}, \hyperlink{bp:finite-fourier}{finite-fourier}\\[3pt]
\hyperlink{bp:derivative-graph}{derivative-graph} & \hyperlink{bp:finite-fourier}{finite-fourier}, Cauchy--Schwarz twice\\[3pt]
\hyperlink{bp:local-frequency}{local-frequency} & \hyperlink{bp:derivative-graph}{derivative-graph}, \hyperlink{bp:BSG}{BSG}, \hyperlink{bp:sumset-growth}{sumset-growth}, \hyperlink{bp:dense-Bohr}{dense-Bohr}\\[3pt]
\hyperlink{bp:Bohr-weight}{Bohr-weight} & \hyperlink{bp:gevrey}{gevrey}, \hyperlink{bp:Chang}{Chang}, character identities\\[3pt]
\hyperlink{bp:local-symmetry}{local-symmetry} & \hyperlink{bp:Bohr-weight}{Bohr-weight}, \hyperlink{bp:dense-Bohr}{dense-Bohr}, Cauchy--Schwarz\\[3pt]
\hyperlink{bp:local-chart}{local-chart} & \hyperlink{bp:lattice-box}{lattice-box}, \hyperlink{bp:integer-basis}{integer-basis}, \hyperlink{bp:gevrey}{gevrey}, \hyperlink{bp:chart-operations}{chart-operations}\\[3pt]
\hyperlink{bp:finite-inverse}{finite-inverse} & \hyperlink{bp:local-frequency}{local-frequency}, \hyperlink{bp:local-symmetry}{local-symmetry}, \hyperlink{bp:local-chart}{local-chart}, \hyperlink{bp:Bohr-weight}{Bohr-weight}, \hyperlink{bp:finite-fourier}{finite-fourier}, \hyperlink{bp:chart-operations}{chart-operations}\\[3pt]
\hyperlink{bp:inverse}{inverse} & \hyperlink{bp:finite-inverse}{finite-inverse}, \hyperlink{bp:cube-foundation}{cube-foundation}, \hyperlink{bp:gevrey}{gevrey}, \hyperlink{bp:chart-operations}{chart-operations}, library Hahn--Banach and Riesz\\[3pt]
\hyperlink{bp:return-splitting}{return-splitting} & \hyperlink{bp:transference}{transference}, \hyperlink{bp:rational-linear}{rational-linear}, \hyperlink{bp:chart-definition}{chart-definition}, finite torus pigeonhole, \hyperlink{bp:budget-rules}{budget-rules}\\[3pt]
\hyperlink{bp:weyl}{weyl} & \hyperlink{bp:chart-operations}{chart-operations}, \hyperlink{bp:fourier-foundation}{fourier-foundation}, \hyperlink{bp:return-splitting}{return-splitting}\\[3pt]
\hyperlink{bp:chart-fourier}{chart-fourier} & \hyperlink{bp:weyl}{weyl}, \hyperlink{bp:chart-operations}{chart-operations}, \hyperlink{bp:fourier-foundation}{fourier-foundation}, \hyperlink{bp:rational-linear}{rational-linear}\\[3pt]
\hyperlink{bp:separation}{separation} & \hyperlink{bp:chart-fourier}{chart-fourier}, \hyperlink{bp:chart-operations}{chart-operations}, \hyperlink{bp:gevrey}{gevrey}, \hyperlink{bp:absolute}{absolute}\\[3pt]
\hyperlink{bp:von-neumann}{von-neumann} & \hyperlink{bp:absolute}{absolute}, \hyperlink{bp:cube-foundation}{cube-foundation}, \hyperlink{bp:fourier-foundation}{fourier-foundation}, \hyperlink{bp:torus-annihilator}{torus-annihilator}\\[3pt]
\hyperlink{bp:detection}{detection} & \hyperlink{bp:inverse}{inverse}, \hyperlink{bp:separation}{separation}, \hyperlink{bp:von-neumann}{von-neumann}, \hyperlink{bp:absolute}{absolute}\\[3pt]
\hyperlink{bp:predictors}{predictors} & \hyperlink{bp:separation}{separation}, \hyperlink{bp:height}{height}, \hyperlink{bp:compression}{compression}, \hyperlink{bp:fourier-foundation}{fourier-foundation}, \hyperlink{bp:absolute}{absolute}\\[3pt]
\hyperlink{bp:delay}{delay} & \hyperlink{bp:predictors}{predictors}, \hyperlink{bp:coeff}{coeff}, \hyperlink{bp:prefix}{prefix}, conditional expectation\\[3pt]
\hyperlink{bp:adaptive-fit}{adaptive-fit} & \hyperlink{bp:detection}{detection}, \hyperlink{bp:coeff}{coeff}, \hyperlink{bp:prefix}{prefix}, \hyperlink{bp:path-packing}{path-packing}\\[3pt]
\hyperlink{bp:stagger}{stagger} & \hyperlink{bp:adaptive-fit}{adaptive-fit}, \hyperlink{bp:absolute}{absolute}, \hyperlink{bp:budget-rules}{budget-rules}\\[3pt]
\hyperlink{bp:base}{base} & \hyperlink{bp:adaptive-fit}{adaptive-fit}, \hyperlink{bp:stagger}{stagger}, \hyperlink{bp:predictors}{predictors}, \hyperlink{bp:delay}{delay}, \hyperlink{bp:detection}{detection}, \hyperlink{bp:height}{height}, \hyperlink{bp:cuts}{cuts}, \hyperlink{bp:path-packing}{path-packing}\\[3pt]
\hyperlink{bp:graph}{graph} & \hyperlink{bp:weyl}{weyl}, \hyperlink{bp:return-splitting}{return-splitting}, \hyperlink{bp:chart-operations}{chart-operations}, \hyperlink{bp:rational-linear}{rational-linear}, \hyperlink{bp:fourier-foundation}{fourier-foundation}\\[3pt]
\hyperlink{bp:coloring}{coloring} & \hyperlink{bp:graph}{graph}, finite geometric probabilities, finite dyadic cohorts, \hyperlink{bp:path-packing}{path-packing}\\[3pt]
\hyperlink{bp:masks}{masks} & \hyperlink{bp:graph}{graph}, \hyperlink{bp:coloring}{coloring}\\[3pt]
\hyperlink{bp:whole-ranges}{whole-ranges} & \hyperlink{bp:masks}{masks}, \hyperlink{bp:graph}{graph}, \hyperlink{bp:coeff}{coeff}\\[3pt]
\hyperlink{bp:same-tag}{same-tag} & \hyperlink{bp:chart-operations}{chart-operations}, \hyperlink{bp:fourier-foundation}{fourier-foundation}, \hyperlink{bp:absolute}{absolute}, \hyperlink{bp:gevrey}{gevrey}\\[3pt]
\hyperlink{bp:closure}{closure} & \hyperlink{bp:same-tag}{same-tag}, \hyperlink{bp:prefix}{prefix}, \hyperlink{bp:base}{base}, \hyperlink{bp:coloring}{coloring}, \hyperlink{bp:masks}{masks}, \hyperlink{bp:coeff}{coeff}, \hyperlink{bp:budget-rules}{budget-rules}\\[3pt]
\hyperlink{bp:localize}{localize} & \hyperlink{bp:closure}{closure}, \hyperlink{bp:masks}{masks}, \hyperlink{bp:graph}{graph}, \hyperlink{bp:base}{base}, \hyperlink{bp:coeff}{coeff}\\[3pt]
\hyperlink{bp:variation}{variation} & \hyperlink{bp:localize}{localize}, \hyperlink{bp:whole-ranges}{whole-ranges}, \hyperlink{bp:compression}{compression}, \hyperlink{bp:prefix}{prefix}\\[3pt]
\hyperlink{bp:first-match}{first-match} & \hyperlink{bp:graph}{graph}, \hyperlink{bp:chart-fourier}{chart-fourier}, \hyperlink{bp:coeff}{coeff}, \hyperlink{bp:localize}{localize}, \hyperlink{bp:gevrey}{gevrey}\\[3pt]
\hyperlink{bp:mode-schedule}{mode-schedule} & \hyperlink{bp:first-match}{first-match}, \hyperlink{bp:height}{height}, \hyperlink{bp:rational-linear}{rational-linear}, \hyperlink{bp:budget-rules}{budget-rules}\\[3pt]
\hyperlink{bp:tiny-space}{tiny-space} & \hyperlink{bp:mode-schedule}{mode-schedule}, \hyperlink{bp:height}{height}, \hyperlink{bp:rational-linear}{rational-linear}\\[3pt]
\hyperlink{bp:stabilize}{stabilize} & \hyperlink{bp:mode-schedule}{mode-schedule}, \hyperlink{bp:tiny-space}{tiny-space}, \hyperlink{bp:rational-linear}{rational-linear}, \hyperlink{bp:integer-basis}{integer-basis}\\[3pt]
\hyperlink{bp:coherence}{coherence} & \hyperlink{bp:stabilize}{stabilize}, \hyperlink{bp:mode-schedule}{mode-schedule}, finite intersections\\[3pt]
\hyperlink{bp:mode-rows}{mode-rows} & \hyperlink{bp:mode-schedule}{mode-schedule}, \hyperlink{bp:coeff}{coeff}, \hyperlink{bp:fourier-foundation}{fourier-foundation}\\[3pt]
\hyperlink{bp:lift-composition}{lift-composition} & \hyperlink{bp:first-match}{first-match}, \hyperlink{bp:compression}{compression}, \hyperlink{bp:whole-ranges}{whole-ranges}\\[3pt]
\hyperlink{bp:linear-side}{linear-side} & \hyperlink{bp:stabilize}{stabilize}, \hyperlink{bp:coherence}{coherence}, \hyperlink{bp:tiny-space}{tiny-space}, \hyperlink{bp:mode-rows}{mode-rows}, \hyperlink{bp:lift-composition}{lift-composition}\\[3pt]
\hyperlink{bp:offset-selection}{offset-selection} & \hyperlink{bp:first-match}{first-match}, \hyperlink{bp:mode-schedule}{mode-schedule}, \hyperlink{bp:linear-side}{linear-side}, finite probability, Tonelli\\[3pt]
\hyperlink{bp:side-caps}{side-caps} & \hyperlink{bp:offset-selection}{offset-selection}, \hyperlink{bp:linear-side}{linear-side}, \hyperlink{bp:variation}{variation}, \hyperlink{bp:base}{base}, \hyperlink{bp:cuts}{cuts}\\[3pt]
\hyperlink{bp:helpers}{helpers} & \hyperlink{bp:closure}{closure}, \hyperlink{bp:side-caps}{side-caps}, \hyperlink{bp:localize}{localize}\\[3pt]
\hyperlink{bp:representatives}{representatives} & \hyperlink{bp:coherence}{coherence}, \hyperlink{bp:stabilize}{stabilize}, \hyperlink{bp:mode-schedule}{mode-schedule}, finite dyadic size classes\\[3pt]
\hyperlink{bp:component-sizes}{component-sizes} & \hyperlink{bp:mode-rows}{mode-rows}, \hyperlink{bp:first-match}{first-match}, \hyperlink{bp:linear-side}{linear-side}\\[3pt]
\hyperlink{bp:tops}{tops} & \hyperlink{bp:representatives}{representatives}, \hyperlink{bp:component-sizes}{component-sizes}, \hyperlink{bp:mode-rows}{mode-rows}, \hyperlink{bp:lift-composition}{lift-composition}, \hyperlink{bp:linear-side}{linear-side}, \hyperlink{bp:side-caps}{side-caps}, \hyperlink{bp:cuts}{cuts}\\[3pt]
\hyperlink{bp:transfer}{transfer} & \hyperlink{bp:helpers}{helpers}, \hyperlink{bp:localize}{localize}, \hyperlink{bp:masks}{masks}, \hyperlink{bp:graph}{graph}, \hyperlink{bp:coeff}{coeff}\\[3pt]
\hyperlink{bp:shortblock}{shortblock} & \hyperlink{bp:finite-CN}{finite-CN}, \hyperlink{bp:finite-fourier}{finite-fourier}, \hyperlink{bp:vector-maximal}{vector-maximal}, \hyperlink{bp:absolute}{absolute}, \hyperlink{bp:fourier-foundation}{fourier-foundation}\\[3pt]
\hyperlink{bp:power-costs}{power-costs} & \hyperlink{bp:side-caps}{side-caps}, \hyperlink{bp:variation}{variation}, \hyperlink{bp:transfer}{transfer}, \hyperlink{bp:base}{base}, \hyperlink{bp:cuts}{cuts}, \hyperlink{bp:absolute}{absolute}\\[3pt]
\hyperlink{bp:fixed-stacks}{fixed-stacks} & \hyperlink{bp:prefix}{prefix}, \hyperlink{bp:variation}{variation}, \hyperlink{bp:linear-side}{linear-side}, \hyperlink{bp:base}{base}, \hyperlink{bp:helpers}{helpers}, \hyperlink{bp:finite-fourier}{finite-fourier}\\[3pt]
\hyperlink{bp:side-estimate}{side-estimate} & \hyperlink{bp:transfer}{transfer}, \hyperlink{bp:shortblock}{shortblock}, \hyperlink{bp:power-costs}{power-costs}, \hyperlink{bp:fixed-stacks}{fixed-stacks}, \hyperlink{bp:cuts}{cuts}, \hyperlink{bp:side-caps}{side-caps}\\[3pt]
\hyperlink{bp:common-height}{common-height} & \hyperlink{bp:coherence}{coherence}, \hyperlink{bp:representatives}{representatives}, \hyperlink{bp:mode-schedule}{mode-schedule}, finite probability, geometric series\\[3pt]
\hyperlink{bp:frequency-tests}{frequency-tests} & \hyperlink{bp:common-height}{common-height}, \hyperlink{bp:torus-annihilator}{torus-annihilator}, \hyperlink{bp:component-sizes}{component-sizes}, \hyperlink{bp:absolute}{absolute}, \hyperlink{bp:fourier-foundation}{fourier-foundation}, \hyperlink{bp:mode-schedule}{mode-schedule}\\[3pt]
\hyperlink{bp:sparse-count}{sparse-count} & \hyperlink{bp:four-paths}{four-paths}, \hyperlink{bp:sumset-growth}{sumset-growth}, finite packing, integer arithmetic\\[3pt]
\hyperlink{bp:tested-counts}{tested-counts} & \hyperlink{bp:frequency-tests}{frequency-tests}, \hyperlink{bp:representatives}{representatives}, \hyperlink{bp:tops}{tops}, \hyperlink{bp:sparse-count}{sparse-count}, \hyperlink{bp:height}{height}, \hyperlink{bp:side-caps}{side-caps}\\[3pt]
\hyperlink{bp:distribution}{distribution} & Tonelli, one-variable power integration\\[3pt]
\hyperlink{bp:core-estimate}{core-estimate} & \hyperlink{bp:tested-counts}{tested-counts}, \hyperlink{bp:distribution}{distribution}, \hyperlink{bp:tops}{tops}, \hyperlink{bp:side-caps}{side-caps}, \hyperlink{bp:frequency-tests}{frequency-tests}, \hyperlink{bp:path-packing}{path-packing}, \hyperlink{bp:budget-rules}{budget-rules}\\[3pt]
\hyperlink{bp:normalized}{normalized} & \hyperlink{bp:base}{base}, \hyperlink{bp:localize}{localize}, \hyperlink{bp:linear-side}{linear-side}, \hyperlink{bp:tops}{tops}, \hyperlink{bp:side-estimate}{side-estimate}, \hyperlink{bp:core-estimate}{core-estimate}, \hyperlink{bp:cuts}{cuts}, \hyperlink{bp:budget-rules}{budget-rules}\\[3pt]
\hyperlink{bp:uniform-local}{uniform-local} & \hyperlink{bp:normalized}{normalized}, \hyperlink{bp:absolute}{absolute}, \hyperlink{bp:finite-CN}{finite-CN}, dyadic first stopping, permitted scalar maximal theory\\[3pt]
\hyperlink{bp:main}{main} & \hyperlink{bp:uniform-local}{uniform-local}, \hyperlink{bp:gevrey}{gevrey}, permitted scalar maximal theory, Schwartz density, dominated convergence\\[3pt]
\end{longtable}

\begin{thebibliography}{9}
\bibitem{M} OpenAI, \emph{An $L^3$ bound for the trilinear Hilbert transform}, supplied manuscript dated 5 October 2026, 93 pages. Source of the singular-integral construction.
\bibitem{R} \emph{An $L^3$ estimate for the trilinear Hilbert transform: A revised, shorter argument}, supplied \nolinkurl{shorter_argument.tex} and corresponding 57-page PDF, 8 October 2026. Source version reorganized here; its external structural imports are replaced by the proofs in this blueprint.
\bibitem{mathlibLp} The mathlib community, \emph{Mathlib.Analysis.Fourier.LpSpace}, documentation and source at the reference revision specified in the first section. \url{https://leanprover-community.github.io/mathlib4_docs/Mathlib/Analysis/Fourier/LpSpace.html}.
\bibitem{mathlibTorus} The mathlib community, \emph{Mathlib.Analysis.Fourier.AddCircleMulti}, documentation and source at the same reference revision. \url{https://leanprover-community.github.io/mathlib4_docs/Mathlib/Analysis/Fourier/AddCircleMulti.html}.
\bibitem{mathlibSchwartz} The mathlib community, \emph{Mathlib.Analysis.Distribution.SchwartzSpace.Basic}, especially \nolinkurl{SchwartzMap.denseRange_toLpCLM}. \url{https://leanprover-community.github.io/mathlib4_docs/Mathlib/Analysis/Distribution/SchwartzSpace/Basic.html}.
\bibitem{mathlibGeometry} The mathlib community, \emph{Mathlib.MeasureTheory.Group.GeometryOfNumbers}, source revision \texttt{6be3d352cec3ceabd28bf85266d28646e02f1975}. \url{https://leanprover-community.github.io/mathlib4_docs/Mathlib/MeasureTheory/Group/GeometryOfNumbers.html}. The nonsharp successive-minima estimate in this blueprint is proved independently.
\bibitem{CS} E. Croot and O. Sisask, \emph{A probabilistic technique for finding almost-periods of convolutions}, arXiv:1003.2978. Background attribution for the sampling method; no theorem from the paper is an import.
\bibitem{GT} B. Green and T. Tao, \emph{An inverse theorem for the Gowers $U^3$ norm}, arXiv:math/0503014, version 4. Background attribution for the inverse-theorem strategy; no theorem from the paper is an import. The lattice coordinate argument used here is proved explicitly.
\end{thebibliography}

\end{document}
